\part*{Part I — Motivation}
\label{navigation:I}
\addcontentsline{toc}{part}{Part I — Motivation}
Stakeholders are cross-cutting Motivation perspectives. Existing associations are retained; no new stakeholder-to-concern associations are inferred.\par
Read the axioms as constraints on the rest of the architecture. Harmonia owns its semantics while using underlying engines for machinery; native engine behaviour does not establish information authority. Semantic distinctions and explicit uncertainty must survive downstream representation. These points explain why the following chapters distinguish currentness, acceptance, responsibility, persistence and execution rather than treating them as interchangeable.

{\small Reading sources: \hyperref[element:54508cda80484a6caf95b901780eaf7d]{AX-04 Harmonia Owns the Semantics; Engines Provide the Machinery}; \hyperref[element:d3fc6025dc7646d1b2020b51f7a44bfa]{AX-14 Semantic Distinctions Are Preserved}; \hyperref[element:8882573df7ce43c6bc0f66da0a676139]{AX-17 Architectural Authority and Explicit Uncertainty}.}\par
\begin{center}\begin{tikzpicture}[>=Stealth,every node/.style={draw=Ink!40,rounded corners,align=left,font=\small,text width=5.2cm,inner sep=6pt}]
\node (source) at (0,0) {\hyperref[element:f4a3060eba464f4ab89b50aec4628e5b]{[RequirementGroup] Non-Functional-Requirements}};
\node (target0) at (8,1.7) {\hyperref[element:abf16578801a466b8f08d80ae292e644]{[RequirementGroup] Availability}};
\draw[->] (source.east) -- (3,0) |- (target0.west);
\node[draw=none,text width=1.8cm,font=\tiny,align=center,inner sep=0pt] at (4.15,2.08) {Composition\newline{}source-qualified / unclassified proposal};
\node (target1) at (8,0.0) {\hyperref[element:cb090e968a25409fb8ceb0fcc9a11ee7]{[RequirementGroup] Disaster Recovery Requirements}};
\draw[->] (source.east) -- (3,0) |- (target1.west);
\node[draw=none,text width=1.8cm,font=\tiny,align=center,inner sep=0pt] at (4.15,0.38) {Composition\newline{}source-qualified / unclassified proposal};
\node (target2) at (8,-1.7) {\hyperref[element:dc3f2a927e284da88734df4f7e0f9311]{[RequirementGroup] Performance Requirements}};
\draw[->] (source.east) -- (3,0) |- (target2.west);
\node[draw=none,text width=1.8cm,font=\tiny,align=center,inner sep=0pt] at (4.15,-1.32) {Composition\newline{}source-qualified / unclassified proposal};
\end{tikzpicture}\end{center}
{\small Selected outgoing Composition assertions; arrows retain source-to-target direction. Labels retain assertion standing and record state; element standing accompanies each description. \href{structure.json}{Complete relationship qualifications and provenance} remain source-defined. No omitted edge or semantic effect is inferred.}\par
\chapter*{I.1 Drivers}
\phantomsection
\label{navigation:I.1}
\addcontentsline{toc}{chapter}{I.1 Drivers}

\section*{I.1.1 24/7 Availability}
\phantomsection
\label{navigation:I.1.1}
\addcontentsline{toc}{section}{I.1.1 24/7 Availability}

\label{element:782f53e2190f4a8a8d0dc0a33cd6fbb3}
\index{24/7 Availability (I.1.1)}
{\small Driver · Historical — current applicability unresolved · historical\_reference}\par
\href{../../../../../../../continuum/motivation/driver/782f53e2190f4a8a8d0dc0a33cd6fbb3.md}{Historical canonical definition and provenance}.\par
\subsection*{I.1.1.1 Stakeholder references — Reference navigation}
\phantomsection
\label{navigation:I.1.1.1}
\addcontentsline{toc}{subsection}{I.1.1.1 Stakeholder references — Reference navigation}

Context from existing canonical Association relationships only; these references do not allocate Stakeholders.\par
\subsubsection*{I.1.1.1.1 Healthcare Provider — Stakeholder reference [historical\_reference]}
\phantomsection
\label{navigation:I.1.1.1.1}
\addcontentsline{toc}{subsubsection}{I.1.1.1.1 Healthcare Provider — Stakeholder reference [historical\_reference]}

{\small Canonical Stakeholder: R.1.6 \hyperref[element:2018200a493d458096b72031cc2ae779]{Healthcare Provider}; Canonical ID: \texttt{\seqsplit{2018200a493d458096b72031cc2ae779}}. Historical — current applicability unresolved · historical\_reference.}\par
Contextual reference only; no additional architectural instance or semantic relationship.\par
Existing canonical associations: id-78a34270ced656e7901d367bfd405999\par
\section*{I.1.2 Audit \& PHI Privacy}
\phantomsection
\label{navigation:I.1.2}
\addcontentsline{toc}{section}{I.1.2 Audit \& PHI Privacy}

\label{element:499313a0a5d84520a3061120ae9b6dac}
\index{Audit \& PHI Privacy (I.1.2)}
{\small Driver · Historical — current applicability unresolved · historical\_reference}\par
\href{../../../../../../../continuum/motivation/driver/499313a0a5d84520a3061120ae9b6dac.md}{Historical canonical definition and provenance}.\par
\subsection*{I.1.2.1 Stakeholder references — Reference navigation}
\phantomsection
\label{navigation:I.1.2.1}
\addcontentsline{toc}{subsection}{I.1.2.1 Stakeholder references — Reference navigation}

Context from existing canonical Association relationships only; these references do not allocate Stakeholders.\par
\subsubsection*{I.1.2.1.1 Patient — Stakeholder reference [historical\_reference]}
\phantomsection
\label{navigation:I.1.2.1.1}
\addcontentsline{toc}{subsubsection}{I.1.2.1.1 Patient — Stakeholder reference [historical\_reference]}

{\small Canonical Stakeholder: R.1.7 \hyperref[element:5f6d333e94634420b8a8bbf8b9cb6e2a]{Patient}; Canonical ID: \texttt{\seqsplit{5f6d333e94634420b8a8bbf8b9cb6e2a}}. Historical — current applicability unresolved · historical\_reference.}\par
Contextual reference only; no additional architectural instance or semantic relationship.\par
Existing canonical associations: id-daee593684c55e55bf1ce04c223cf585\par
\section*{I.1.3 Clinical Interoperability}
\phantomsection
\label{navigation:I.1.3}
\addcontentsline{toc}{section}{I.1.3 Clinical Interoperability}

\label{element:babb9f5d63d8435890c1db3579ec78d9}
\index{Clinical Interoperability (I.1.3)}
{\small Driver · Historical — current applicability unresolved · historical\_reference}\par
\href{../../../../../../../continuum/motivation/driver/babb9f5d63d8435890c1db3579ec78d9.md}{Historical canonical definition and provenance}.\par
\subsection*{I.1.3.1 Stakeholder references — Reference navigation}
\phantomsection
\label{navigation:I.1.3.1}
\addcontentsline{toc}{subsection}{I.1.3.1 Stakeholder references — Reference navigation}

Context from existing canonical Association relationships only; these references do not allocate Stakeholders.\par
\subsubsection*{I.1.3.1.1 Health Network — Stakeholder reference [historical\_reference]}
\phantomsection
\label{navigation:I.1.3.1.1}
\addcontentsline{toc}{subsubsection}{I.1.3.1.1 Health Network — Stakeholder reference [historical\_reference]}

{\small Canonical Stakeholder: R.1.3 \hyperref[element:55e57c01dfe146a2a225228b31196862]{Health Network}; Canonical ID: \texttt{\seqsplit{55e57c01dfe146a2a225228b31196862}}. Historical — current applicability unresolved · historical\_reference.}\par
Contextual reference only; no additional architectural instance or semantic relationship.\par
Existing canonical associations: id-d5fa8c64975852f88132df579c806836\par
\subsubsection*{I.1.3.1.2 Care Team — Stakeholder reference [historical\_reference]}
\phantomsection
\label{navigation:I.1.3.1.2}
\addcontentsline{toc}{subsubsection}{I.1.3.1.2 Care Team — Stakeholder reference [historical\_reference]}

{\small Canonical Stakeholder: R.1.1 \hyperref[element:6f5a203d81e24a81b384a80781987daa]{Care Team}; Canonical ID: \texttt{\seqsplit{6f5a203d81e24a81b384a80781987daa}}. Historical — current applicability unresolved · historical\_reference.}\par
Contextual reference only; no additional architectural instance or semantic relationship.\par
Existing canonical associations: id-457a3dd70bb358ff8989bc3c9c725c70\par
\section*{I.1.4 Clinical Interoperability Across Heterogeneous Standards}
\phantomsection
\label{navigation:I.1.4}
\addcontentsline{toc}{section}{I.1.4 Clinical Interoperability Across Heterogeneous Standards}

\label{element:921d7cf84d2b4abca0deb0d300fa75ea}
\index{Clinical Interoperability Across Heterogeneous Standards (I.1.4)}
{\small Driver · Agreed · current}\par

\par\medskip\textbf{Context \& Pressure}\par


Regional healthcare environments are characterized by decades of technical layering. Acute care hospitals, outpatient clinics, pathology laboratories, and diagnostic imaging centers operate software solutions spanning multiple technological eras:
- Legacy hospital information systems communicating via HL7 v2.x pipe-and-hat (\texttt{\seqsplit{ER7}}) messaging over Minimal Lower Layer Protocol (MLLP).
- Modern clinical repositories exposing HL7 FHIR RESTful APIs (R4/R5) using JSON and XML encodings.
- Specialized imaging platforms exchanging DICOM objects.
- Proprietary file drops, CSV exports, and ad-hoc database integrations.



\par\medskip\textbf{Architectural Imperative}\par


Harmonia must bridge these heterogeneous standards without forcing disruptive, multi-million-dollar system replacements across delivery organizations. The platform must accept standards-based messages natively at its integration membrane while maintaining canonical semantic coherence internally, preventing message corruption or loss of clinical intent during translation.



\begin{itemize}


\item \textbf{Primary Stakeholders}: \href{\detokenize{../../../../../../../../docs/markdown/01-motivation/stakeholders/enterprise-stakeholders.md}}{Regional Health Network Operator}, \href{\detokenize{../../../../../../../../docs/markdown/01-motivation/stakeholders/enterprise-stakeholders.md}}{Healthcare Delivery Organizations}

\item \textbf{External Constraint}: \href{\detokenize{../../../../../../../../docs/markdown/01-motivation/requirements-constraints/external-constraints.md}}{Mandated External Interoperability Contracts}

\item \textbf{Associated Goal}: \href{\detokenize{../../../../../../../../docs/markdown/01-motivation/goals-outcomes/strategic-goals.md}}{Vendor-Independent Longitudinal Clinical Information Coherence}

\item \textbf{Governing Axioms}: \hyperref[element:eee8c3783d964d18a7f9dcea85fcd4cf]{AX-02 (Standards at Boundary)}, \hyperref[element:e99a4e2bf99f476d876394300d6b9c3d]{AX-03 (Native Standards Representations)}

\end{itemize}

{\footnotesize \href{../../../../../../../continuum/motivation/driver/921d7cf84d2b4abca0deb0d300fa75ea.md}{Canonical record: 921d7cf84d2b4abca0deb0d300fa75ea}; full qualifications and provenance.}\par
\section*{I.1.5 Continuous Clinical Service Availability}
\phantomsection
\label{navigation:I.1.5}
\addcontentsline{toc}{section}{I.1.5 Continuous Clinical Service Availability}

\label{element:80e2d5f142da4dae8b7bb70e9ddcab06}
\index{Continuous Clinical Service Availability (I.1.5)}
{\small Driver · Agreed · current}\par

\par\medskip\textbf{Context \& Pressure}\par


Healthcare delivery is continuous, mission-critical, and time-sensitive (24/7/365). In emergency departments, intensive care units, and operating theaters, clinical decisions cannot be paused for integration system downtime, database maintenance, or middleware restarts.



A failure in the integration environment directly impacts patient care: an emergency physician cannot view critical drug allergy records; a surgeon cannot inspect pre-operative pathology results; an inpatient unit cannot receive bed transfer notifications.



\par\medskip\textbf{Architectural Imperative}\par


Harmonia must provide continuous clinical service availability. The platform must withstand hardware failures, network partitions, and downstream destination outages without dropping incoming clinical events or blocking point-of-care clinical queries.



\begin{itemize}


\item \textbf{Primary Stakeholders}: \href{\detokenize{../../../../../../../../docs/markdown/01-motivation/stakeholders/enterprise-stakeholders.md}}{Clinicians \& Care Teams}, \href{\detokenize{../../../../../../../../docs/markdown/01-motivation/stakeholders/enterprise-stakeholders.md}}{Platform Operations \& Integration Engineers}, \href{\detokenize{../../../../../../../../docs/markdown/01-motivation/stakeholders/enterprise-stakeholders.md}}{Healthcare Delivery Organizations}

\item \textbf{Key Assessments}: \href{\detokenize{../../../../../../../../docs/markdown/01-motivation/drivers-assessments/assessments.md}}{Silent Data Loss via False Acceptance}, \href{\detokenize{../../../../../../../../docs/markdown/01-motivation/drivers-assessments/assessments.md}}{Centralised Synchronous Persistence Can Constrain Concurrency}

\item \textbf{Associated Goals}: \href{\detokenize{../../../../../../../../docs/markdown/01-motivation/goals-outcomes/strategic-goals.md}}{Durable Acceptance \& Preservation of Clinical Events}, \href{\detokenize{../../../../../../../../docs/markdown/01-motivation/goals-outcomes/strategic-goals.md}}{Responsive Access to Managed Information}

\item \textbf{Governing Axioms}: \hyperref[element:0d9673833eb840eab03804aa076c1b0c]{AX-10 (Failure as Normal Condition)}, \hyperref[element:8a1e995f16954484a78ae1c5c114d546]{AX-11 (Responsive \& Resilient Access)}, \hyperref[element:0412a2084f0145babc4f022c71fea65f]{AX-05 (State Separation)}

\end{itemize}

{\footnotesize \href{../../../../../../../continuum/motivation/driver/80e2d5f142da4dae8b7bb70e9ddcab06.md}{Canonical record: 80e2d5f142da4dae8b7bb70e9ddcab06}; full qualifications and provenance.}\par
\section*{I.1.6 Coordinated Operational Activity and Entity State}
\phantomsection
\label{navigation:I.1.6}
\addcontentsline{toc}{section}{I.1.6 Coordinated Operational Activity and Entity State}

\label{element:b3fb294d95024e4382b2c251c389fa66}
\index{Coordinated Operational Activity and Entity State (I.1.6)}
{\small Driver · Agreed · current}\par

\par\medskip\textbf{Context \& Pressure}\par


Healthcare integration is rarely a simple point-to-point message transmission. A single real-world clinical event—such as a patient admission or an emergency department transfer—triggers a cascade of multi-stage operational activities across multiple downstream systems (e.g., notifying pharmacy, updating bed management, triggering pre-admission lab orders, broadcasting HL7 ADT messages to billing and clinical portals).



When integration middleware treats these complex interactions simply as a stream of disconnected, fire-and-forget messages, failures in downstream destinations cause silent state divergence. The hospital's billing system may record the patient as discharged while the pharmacy continues dispensing medication.



\par\medskip\textbf{Architectural Imperative}\par


Harmonia must coordinate multi-stage operational activity in lockstep with the governed state of the real-world clinical entity. Activities must progress through explicit, observable state transitions, ensuring that fan-out status, acknowledgements, and retry progressions are fully trackable and recoverable.



\begin{itemize}


\item \textbf{Primary Stakeholders}: \href{\detokenize{../../../../../../../../docs/markdown/01-motivation/stakeholders/enterprise-stakeholders.md}}{Healthcare Delivery Organizations}, \href{\detokenize{../../../../../../../../docs/markdown/01-motivation/stakeholders/enterprise-stakeholders.md}}{Platform Operations \& Integration Engineers}

\item \textbf{Key Assessment}: \href{\detokenize{../../../../../../../../docs/markdown/01-motivation/drivers-assessments/assessments.md}}{Unmonitored Destination Failure / Fan-Out Divergence}

\item \textbf{Associated Goal}: \href{\detokenize{../../../../../../../../docs/markdown/01-motivation/goals-outcomes/strategic-goals.md}}{Coordinated Progression of Operational Activities \& Associated Entity State}

\item \textbf{Governing Axiom}: \hyperref[element:de7e90f46127455a9ec9ad3b9b48ba54]{AX-16 (Operational Activity \& Entity State Progress Together)}

\end{itemize}

{\footnotesize \href{../../../../../../../continuum/motivation/driver/b3fb294d95024e4382b2c251c389fa66.md}{Canonical record: b3fb294d95024e4382b2c251c389fa66}; full qualifications and provenance.}\par
\section*{I.1.7 Durable Clinical Information Independence}
\phantomsection
\label{navigation:I.1.7}
\addcontentsline{toc}{section}{I.1.7 Durable Clinical Information Independence}

\label{element:3d81912dd8c94603ba0e10c98b7a62d8}
\index{Durable Clinical Information Independence (I.1.7)}
{\small Driver · Agreed · current}\par

\par\medskip\textbf{Context \& Pressure}\par


Clinical information must remain intelligible, authoritative, and actionable across a patient's entire lifetime (often 80+ years). In contrast, commercial Electronic Medical Record (EMR) software products, proprietary database schemas, and integration broker technologies have commercial lifecycles of only 5 to 10 years.



When healthcare networks store longitudinal patient data exclusively inside proprietary EMR schemas, they suffer severe vendor lock-in, exorbitant data extraction fees, and the perpetual risk of data obsolescence during vendor transitions.



\par\medskip\textbf{Architectural Imperative}\par


Harmonia must establish durable clinical information independence. Patient health data managed by Harmonia must exist as a vendor-neutral, sovereign regional asset. The clinical representation must outlive the commercial life of any individual software vendor, database engine, or middleware framework.



\begin{itemize}


\item \textbf{Primary Stakeholders}: \href{\detokenize{../../../../../../../../docs/markdown/01-motivation/stakeholders/enterprise-stakeholders.md}}{Clinicians \& Care Teams}, \href{\detokenize{../../../../../../../../docs/markdown/01-motivation/stakeholders/enterprise-stakeholders.md}}{Healthcare Delivery Organizations}

\item \textbf{Associated Goal}: \href{\detokenize{../../../../../../../../docs/markdown/01-motivation/goals-outcomes/strategic-goals.md}}{Vendor-Independent Longitudinal Clinical Information Coherence}

\item \textbf{Governing Axioms}: \hyperref[element:e9bd61acaaea4577bc4086a5639e52c3]{AX-01 (Health-Information Centricity)}, \hyperref[element:952c3b6148ce4fce94b69f99a0c2fec7]{AX-06 (Explicit Information Authority)}

\end{itemize}

{\footnotesize \href{../../../../../../../continuum/motivation/driver/3d81912dd8c94603ba0e10c98b7a62d8.md}{Canonical record: 3d81912dd8c94603ba0e10c98b7a62d8}; full qualifications and provenance.}\par
\section*{I.1.8 National Healthcare Directory / Endpoint Governance}
\phantomsection
\label{navigation:I.1.8}
\addcontentsline{toc}{section}{I.1.8 National Healthcare Directory / Endpoint Governance}

\label{element:5ea7e57dcde04ef88c19a88ade9af6ca}
\index{National Healthcare Directory / Endpoint Governance (I.1.8)}
{\small Driver · Agreed · current}\par

\par\medskip\textbf{Context \& Pressure}\par


Healthcare delivery requires seamless coordination between hundreds of clinics, hospitals, diagnostic practices, and thousands of accredited practitioners. Clinical correspondence, diagnostic reports, and e-referrals frequently go astray due to outdated contact details, retired practitioners, or misconfigured electronic messaging endpoints.



National regulatory frameworks (such as the Australian HI Service and AHPRA) publish authoritative registers of practitioners and healthcare organizations. Electronic messaging networks require verified service endpoints to ensure confidential messages are delivered only to legitimate recipients.



\par\medskip\textbf{Architectural Imperative}\par


Harmonia must govern practitioner credentials and electronic service endpoints through authoritative directory mechanisms. Integration workflows must validate that practitioner identities, roles, and destination addresses are authenticated and currently authorized.



\begin{itemize}


\item \textbf{Primary Stakeholders}: \href{\detokenize{../../../../../../../../docs/markdown/01-motivation/stakeholders/enterprise-stakeholders.md}}{Healthcare Directory Stewards \& Registrars}, \href{\detokenize{../../../../../../../../docs/markdown/01-motivation/stakeholders/enterprise-stakeholders.md}}{Clinicians \& Care Teams}

\item \textbf{External Authorities}: \href{\detokenize{../../../../../../../../docs/markdown/01-motivation/stakeholders/external-authorities.md}}{HI Service}, \href{\detokenize{../../../../../../../../docs/markdown/01-motivation/stakeholders/external-authorities.md}}{AHPRA}

\item \textbf{Associated Goal}: \href{\detokenize{../../../../../../../../docs/markdown/01-motivation/goals-outcomes/strategic-goals.md}}{Reliable Subject Identity \& Referential Integrity}

\item \textbf{Governing Axiom}: \hyperref[element:d3fc6025dc7646d1b2020b51f7a44bfa]{AX-14 (Preservation of Semantic Distinctions)}

\end{itemize}

{\footnotesize \href{../../../../../../../continuum/motivation/driver/5ea7e57dcde04ef88c19a88ade9af6ca.md}{Canonical record: 5ea7e57dcde04ef88c19a88ade9af6ca}; full qualifications and provenance.}\par
\section*{I.1.9 Patient Safety \& Identity Integrity}
\phantomsection
\label{navigation:I.1.9}
\addcontentsline{toc}{section}{I.1.9 Patient Safety \& Identity Integrity}

\label{element:b835c7b1055e4dbc9e1cd3a0c545c1c8}
\index{Patient Safety \& Identity Integrity (I.1.9)}
{\small Driver · Agreed · current}\par

\par\medskip\textbf{Context \& Pressure}\par


Patient safety is a foundational operational priority in healthcare. Accurately associating a clinical observation, pathology result, or medication order with the correct individual human subject is essential.



Mismatched patient records lead directly to severe clinical harm: administering contraindicated blood transfusions, giving incorrect medication doses to patients with unrecorded allergies, or performing wrong-patient interventions. In distributed regional networks, patients move between multiple facilities where different local medical record numbers (MRNs) are assigned.



\par\medskip\textbf{Architectural Imperative}\par


Harmonia must rigorously preserve and validate subject referential integrity across all managed clinical events. It must prevent ambiguous, invalid, or contradictory subject associations from being silently treated as valid.



\emph{Note on Architectural Boundary}: Harmonia does not act as an enterprise Master Patient Index (MPI) or deterministic golden-record matcher. Rather, it enforces strict referential validation and transparently captures identifier provenance without lossy unification.



\begin{itemize}


\item \textbf{Primary Stakeholders}: \href{\detokenize{../../../../../../../../docs/markdown/01-motivation/stakeholders/enterprise-stakeholders.md}}{Patients \& Care Recipients}, \href{\detokenize{../../../../../../../../docs/markdown/01-motivation/stakeholders/enterprise-stakeholders.md}}{Clinicians \& Care Teams}

\item \textbf{External Constraint}: \href{\detokenize{../../../../../../../../docs/markdown/01-motivation/requirements-constraints/external-constraints.md}}{Applicable National / Jurisdictional Healthcare Identifier Obligations}

\item \textbf{Associated Goal}: \href{\detokenize{../../../../../../../../docs/markdown/01-motivation/goals-outcomes/strategic-goals.md}}{Reliable Subject Identity \& Referential Integrity}

\item \textbf{Governing Axioms}: \hyperref[element:952c3b6148ce4fce94b69f99a0c2fec7]{AX-06 (Explicit Information Authority)}, \hyperref[element:d3fc6025dc7646d1b2020b51f7a44bfa]{AX-14 (Preservation of Semantic Distinctions)}

\end{itemize}

{\footnotesize \href{../../../../../../../continuum/motivation/driver/b835c7b1055e4dbc9e1cd3a0c545c1c8.md}{Canonical record: b835c7b1055e4dbc9e1cd3a0c545c1c8}; full qualifications and provenance.}\par
\section*{I.1.10 Statutory Health Information Privacy \& Protection}
\phantomsection
\label{navigation:I.1.10}
\addcontentsline{toc}{section}{I.1.10 Statutory Health Information Privacy \& Protection}

\label{element:1dbf34b452b6497f914c704e037aec24}
\index{Statutory Health Information Privacy \& Protection (I.1.10)}
{\small Driver · Agreed · current}\par

\par\medskip\textbf{Context \& Pressure}\par


Health records contain the most sensitive category of personal data recognized under law (medical histories, psychiatric evaluations, sexual health diagnoses, genetic profiles). Healthcare networks are subject to rigorous statutory privacy regimes (e.g., the Australian \emph{Privacy Act 1988}, My Health Records Act, HIPAA, GDPR).



Statutory authorities enforce severe civil penalties, mandatory breach reporting, and reputational sanctions for unauthorized disclosure, eavesdropping, or accidental data leakage.



\par\medskip\textbf{Architectural Imperative}\par


Harmonia must implement intrinsic, default-deny security and privacy controls across all operations. Sensitive Protected Health Information (PHI) must never be emitted into non-clinical logging streams, telemetry systems, or debug channels. Access must require explicit cryptographic authentication and fine-grained authorization.



\begin{itemize}


\item \textbf{Primary Stakeholders}: \href{\detokenize{../../../../../../../../docs/markdown/01-motivation/stakeholders/enterprise-stakeholders.md}}{Patients \& Care Recipients}, \href{\detokenize{../../../../../../../../docs/markdown/01-motivation/stakeholders/enterprise-stakeholders.md}}{Platform Operations \& Integration Engineers}

\item \textbf{External Constraint}: \href{\detokenize{../../../../../../../../docs/markdown/01-motivation/requirements-constraints/external-constraints.md}}{Applicable Health-Information Privacy and Data-Protection Obligations}

\item \textbf{Key Assessments}: \href{\detokenize{../../../../../../../../docs/markdown/01-motivation/drivers-assessments/assessments.md}}{Implicit Perimeter Trust}, \href{\detokenize{../../../../../../../../docs/markdown/01-motivation/drivers-assessments/assessments.md}}{PHI Leakage through Operational Logging}

\item \textbf{Governing Axioms}: \hyperref[element:5de2ae2c3b794c2d9f3a14efd51d564f]{AX-07 (Security Is Intrinsic to Managed Operations)}, \hyperref[element:21d6864657964ba5b03e3f6b5eb29cd5]{AX-08 (Evidence Records Meaning)}

\end{itemize}

{\footnotesize \href{../../../../../../../continuum/motivation/driver/1dbf34b452b6497f914c704e037aec24.md}{Canonical record: 1dbf34b452b6497f914c704e037aec24}; full qualifications and provenance.}\par
\chapter*{I.2 Strategic Goals}
\phantomsection
\label{navigation:I.2}
\addcontentsline{toc}{chapter}{I.2 Strategic Goals}

\section*{I.2.1 Coordinated Progression of Operational Activities and Associated Entity State}
\phantomsection
\label{navigation:I.2.1}
\addcontentsline{toc}{section}{I.2.1 Coordinated Progression of Operational Activities and Associated Entity State}

\label{element:221eaf405d454c749112ece28ac595f1}
\index{Coordinated Progression of Operational Activities and Associated Entity State (I.2.1)}
{\small Goal · Agreed · current}\par

\par\medskip\textbf{Goal Statement}\par


Harmonia shall coordinate multi-stage operational activities through explicit, observable state transitions that remain strictly synchronized with the governed state of the associated real-world clinical entity.



\par\medskip\textbf{Rationale \& Context}\par


Healthcare integration activities (such as ADT distribution, lab order dispatch, and multi-destination fan-out) involve multi-stage distributed operations. Treating integration simply as a stream of fire-and-forget message transfers causes destination divergence and unmonitored failures (as diagnosed in \href{\detokenize{../../../../../../../../docs/markdown/01-motivation/drivers-assessments/assessments.md}}{Unmonitored Destination Failure}).



\par\medskip\textbf{Architectural Boundary \& Intent}\par


\begin{itemize}


\item \textbf{Beyond Conventional Workflow Engines}: Avoid language implying that Harmonia is merely a generic workflow engine. Harmonia does not simply execute arbitrary BPMN scripts; it coordinates distributed integration activities in lockstep with the real-world lifecycle of clinical entities (e.g., patient admission state, lab specimen lifecycle, encounter status).

\item Requires explicit, observable progression state across multi-stage activities (\href{\detokenize{../../../../../../../../docs/markdown/01-motivation/requirements-constraints/foundational-requirements.md}}{\texttt{\seqsplit{REQ-FND-002}}}).

\item Governed by \hyperref[element:de7e90f46127455a9ec9ad3b9b48ba54]{AX-16 (Operational Activity \& Entity State Progress Together)}.

\end{itemize}

{\footnotesize \href{../../../../../../../continuum/motivation/goal/221eaf405d454c749112ece28ac595f1.md}{Canonical record: 221eaf405d454c749112ece28ac595f1}; full qualifications and provenance.}\par
\section*{I.2.2 Durable Acceptance and Preservation of Clinical Events}
\phantomsection
\label{navigation:I.2.2}
\addcontentsline{toc}{section}{I.2.2 Durable Acceptance and Preservation of Clinical Events}

\label{element:44daa77317aa44b7967b62b295d7288d}
\index{Durable Acceptance and Preservation of Clinical Events (I.2.2)}
{\small Goal · Agreed · current}\par

\par\medskip\textbf{Goal Statement}\par


Harmonia shall ensure that once an incoming clinical event is accepted by the platform, that event is guaranteed to be durably preserved and cannot be silently lost due to subsequent crashes, power interruptions, hardware faults, or transient network partitions.



\par\medskip\textbf{Rationale \& Context}\par


Clinical events—such as patient admissions, pathology test orders, critical laboratory alerts, and clinical discharge summaries—carry direct patient safety implications. The platform must eliminate the critical vulnerability identified in \href{\detokenize{../../../../../../../../docs/markdown/01-motivation/drivers-assessments/assessments.md}}{Silent Data Loss via False Acceptance}, ensuring that acceptance at the ingress membrane is backed by authoritative durability before responsibility is released.



\par\medskip\textbf{Architectural Boundary}\par


\begin{itemize}


\item This goal defines the requirement for durable boundary crossing.

\item It does \textbf{not} mandate a specific storage technology (such as a relational database, append-only log, or distributed file system).

\item Downstream realization is governed by \href{\detokenize{../../../../../../../../docs/markdown/01-motivation/requirements-constraints/foundational-requirements.md}}{\texttt{\seqsplit{REQ-FND-001}} Durable Ingress Acceptance Boundary}.

\end{itemize}

{\footnotesize \href{../../../../../../../continuum/motivation/goal/44daa77317aa44b7967b62b295d7288d.md}{Canonical record: 44daa77317aa44b7967b62b295d7288d}; full qualifications and provenance.}\par
\section*{I.2.3 Real-Time Ingestion}
\phantomsection
\label{navigation:I.2.3}
\addcontentsline{toc}{section}{I.2.3 Real-Time Ingestion}

\label{element:d43c345839e548589e2b6aa677022f8b}
\index{Real-Time Ingestion (I.2.3)}
{\small Goal · Historical — current applicability unresolved · historical\_reference}\par
\href{../../../../../../../continuum/motivation/goal/d43c345839e548589e2b6aa677022f8b.md}{Historical canonical definition and provenance}.\par
\section*{I.2.4 Reliable Subject Identity and Referential Integrity}
\phantomsection
\label{navigation:I.2.4}
\addcontentsline{toc}{section}{I.2.4 Reliable Subject Identity and Referential Integrity}

\label{element:e2f8b62e2cbd486282d68b6009c06de6}
\index{Reliable Subject Identity and Referential Integrity (I.2.4)}
{\small Goal · Agreed · current}\par

\par\medskip\textbf{Goal Statement}\par


Harmonia shall ensure that all managed clinical events, observations, and diagnostic records maintain verified, internally consistent, and referentially sound associations with the correct human patient subject.



\par\medskip\textbf{Rationale \& Context}\par


Accurate patient identification is fundamental to patient safety. Cross-associating clinical records between different individuals leads directly to contraindicated treatments and severe diagnostic errors. In regional health networks where patients possess multiple local Medical Record Numbers (MRNs) across different hospitals, the platform must preserve and validate identity references across every managed event.



\par\medskip\textbf{Architectural Boundaries \& Explicit Anti-Goals}\par


Harmonia explicitly defines what this goal \textbf{does NOT entail}:
- \textbf{Harmonia is NOT a Master Patient Index (MPI)}: Harmonia does not maintain an enterprise patient demographic master record.
- \textbf{Harmonia is NOT a Golden-Record Authority}: The platform does not synthesize or assert a single "golden" demographic record.
- \textbf{Harmonia is NOT a Deterministic Identity Matcher}: Harmonia does not execute probabilistic or deterministic patient matching algorithms to merge patient identities.
- \emph{Downstream Handoff}: Questions of enterprise identity matching, demographic curation, and golden-record master indices are intentionally deferred to \textbf{Domain 02 Strategy} and \textbf{Domain 04 Information Architecture}. In Domain 01, this goal enforces strict referential integrity, identifier validation, and non-destructive provenance capture (\href{\detokenize{../../../../../../../../docs/markdown/01-motivation/requirements-constraints/foundational-requirements.md}}{\texttt{\seqsplit{REQ-FND-003}}}).


{\footnotesize \href{../../../../../../../continuum/motivation/goal/e2f8b62e2cbd486282d68b6009c06de6.md}{Canonical record: e2f8b62e2cbd486282d68b6009c06de6}; full qualifications and provenance.}\par
\section*{I.2.5 Responsive Access to Managed Information}
\phantomsection
\label{navigation:I.2.5}
\addcontentsline{toc}{section}{I.2.5 Responsive Access to Managed Information}

\label{element:c7b8639a922f4d1cb2fcc68933cf7a90}
\index{Responsive Access to Managed Information (I.2.5)}
{\small Goal · Agreed · current}\par

\par\medskip\textbf{Goal Statement}\par


Harmonia shall provide consumers and internal services with timely, highly available, and responsive access to managed information and active coordination state, decoupled from the latency and lock contention of central synchronous persistence.



\par\medskip\textbf{Rationale \& Context}\par


In a regional network processing thousands of concurrent clinical interactions, forcing every active query and coordination check through central database locks introduces severe contention and availability bottlenecks (as diagnosed in \href{\detokenize{../../../../../../../../docs/markdown/01-motivation/drivers-assessments/assessments.md}}{Centralised Synchronous Persistence Can Constrain Concurrent Processing}). Consumers require responsive, concurrent read paths that remain available even during background batch processing or heavy ingestion spikes.



\par\medskip\textbf{Architectural Boundary \& Technology Independence}\par


\begin{itemize}


\item \textbf{Strictly Technology-Independent}: This goal does \textbf{not} mandate caching, in-memory grids, Infinispan, Redis, or non-blocking reactive frameworks.

\item It defines the architectural intent: separating active, responsive query and coordination paths from heavy transactional persistence locks.

\item Realized downstream under \hyperref[element:8a1e995f16954484a78ae1c5c114d546]{AX-11 (Responsive \& Resilient Access)} and \hyperref[element:0412a2084f0145babc4f022c71fea65f]{AX-05 (State Separation)}.

\end{itemize}

{\footnotesize \href{../../../../../../../continuum/motivation/goal/c7b8639a922f4d1cb2fcc68933cf7a90.md}{Canonical record: c7b8639a922f4d1cb2fcc68933cf7a90}; full qualifications and provenance.}\par
\section*{I.2.6 Sub-ms Cache Access}
\phantomsection
\label{navigation:I.2.6}
\addcontentsline{toc}{section}{I.2.6 Sub-ms Cache Access}

\label{element:f699096b8d0f4f5690eb57b25001b5e5}
\index{Sub-ms Cache Access (I.2.6)}
{\small Goal · Historical — current applicability unresolved · historical\_reference}\par
\href{../../../../../../../continuum/motivation/goal/f699096b8d0f4f5690eb57b25001b5e5.md}{Historical canonical definition and provenance}.\par
\section*{I.2.7 Timely Availability of Clinical Information}
\phantomsection
\label{navigation:I.2.7}
\addcontentsline{toc}{section}{I.2.7 Timely Availability of Clinical Information}

\label{element:26e84fb0acfc49f1ac501b64a6e4bbf7}
\index{Timely Availability of Clinical Information (I.2.7)}
{\small Goal · Agreed · current}\par

\par\medskip\textbf{Goal Statement}\par


Harmonia shall make clinical information, event notifications, and diagnostic updates available to care teams and consuming hospital systems within operational latency thresholds appropriate to patient care.



\par\medskip\textbf{Rationale \& Context}\par


Healthcare delivery operates under tight clinical timelines. Point-of-care decisions in emergency departments and intensive care units require rapid turnaround of diagnostic findings. While high throughput is vital, low latency for critical alerts (e.g., life-threatening pathology results) is paramount to prevent diagnostic delays.



\par\medskip\textbf{Architectural Boundary}\par


\begin{itemize}


\item Governs the latency and throughput profile of integration pipelines.

\item Differentiates between urgent real-time clinical feeds and bulk historical synchronizations.

\item Hands over directly to \href{\detokenize{../../../../../../../../docs/markdown/01-motivation/goals-outcomes/business-outcomes.md}}{Continuous and Resilient Regional Health Information Exchange (O2)}.

\end{itemize}

{\footnotesize \href{../../../../../../../continuum/motivation/goal/26e84fb0acfc49f1ac501b64a6e4bbf7.md}{Canonical record: 26e84fb0acfc49f1ac501b64a6e4bbf7}; full qualifications and provenance.}\par
\section*{I.2.8 Vendor-Independent Longitudinal Clinical Information Coherence}
\phantomsection
\label{navigation:I.2.8}
\addcontentsline{toc}{section}{I.2.8 Vendor-Independent Longitudinal Clinical Information Coherence}

\label{element:905aca03c85748d79c3c065f6f16fc0a}
\index{Vendor-Independent Longitudinal Clinical Information Coherence (I.2.8)}
{\small Goal · Agreed · current}\par

\par\medskip\textbf{Goal Statement}\par


Harmonia shall establish and maintain a vendor-neutral, semantically coherent representation of a patient's longitudinal health record that outlives individual EMR software lifecycles and remains independent of proprietary vendor schemas.



\par\medskip\textbf{Rationale \& Context}\par


Patient care spans decades across multiple disparate healthcare institutions. Commercial EMR platforms have short market lifecycles and store data in proprietary, closed schemas that create vendor lock-in. To ensure continuity of care, the regional health network requires an enduring clinical asset that preserves the semantic meaning, units of measure, clinical nuance, and temporal sequence of health events across disparate provider systems.



\par\medskip\textbf{Architectural Boundary}\par


\begin{itemize}


\item Governed by \hyperref[element:e9bd61acaaea4577bc4086a5639e52c3]{AX-01 (Health-Information Centricity)}, \hyperref[element:eee8c3783d964d18a7f9dcea85fcd4cf]{AX-02 (Standards at Boundary)}, and \hyperref[element:952c3b6148ce4fce94b69f99a0c2fec7]{AX-06 (Explicit Information Authority)}.

\item Hands over directly to \href{\detokenize{../../../../../../../../docs/markdown/01-motivation/goals-outcomes/business-outcomes.md}}{Reduced Clinical Risk (O1)} and \href{\detokenize{../../../../../../../../docs/markdown/01-motivation/goals-outcomes/business-outcomes.md}}{Continuous Regional Exchange (O2)}.

\end{itemize}

{\footnotesize \href{../../../../../../../continuum/motivation/goal/905aca03c85748d79c3c065f6f16fc0a.md}{Canonical record: 905aca03c85748d79c3c065f6f16fc0a}; full qualifications and provenance.}\par
\section*{I.2.9 Zero Message Loss}
\phantomsection
\label{navigation:I.2.9}
\addcontentsline{toc}{section}{I.2.9 Zero Message Loss}

\label{element:8f70dbb9c28f4dd9afe91c21243e5914}
\index{Zero Message Loss (I.2.9)}
{\small Goal · Historical — current applicability unresolved · historical\_reference}\par
\href{../../../../../../../continuum/motivation/goal/8f70dbb9c28f4dd9afe91c21243e5914.md}{Historical canonical definition and provenance}.\par
\section*{I.2.10 Architecture-Implementation Alignment}
\phantomsection
\label{navigation:I.2.10}
\addcontentsline{toc}{section}{I.2.10 Architecture-Implementation Alignment}

\label{element:adf5bd2a409b44099253fbed486d736f}
\index{Architecture-Implementation Alignment (I.2.10)}
{\small Goal · Solution Architect-reviewed Iteration 001; authorised ADD; resulting increment review pending · current}\par

Maintain demonstrable alignment between Harmonia's governed architecture and its implemented solutions throughout the solution lifecycle.



Source entry: \textbf{NFR-152 — Architecture-Implementation Alignment}. Primary subject: \texttt{\seqsplit{G01.04}}.



The source classification is Goal. The source identifier is a scoped register reference and does not determine architectural type. Applicability, release scope and unresolved thresholds remain as stated above.


{\footnotesize \href{../../../../../../../continuum/motivation/goal/adf5bd2a409b44099253fbed486d736f.md}{Canonical record: adf5bd2a409b44099253fbed486d736f}; full qualifications and provenance.}\par
\section*{I.2.11 External Dependency Resilience}
\phantomsection
\label{navigation:I.2.11}
\addcontentsline{toc}{section}{I.2.11 External Dependency Resilience}

\label{element:150b3b1eea9b484b89ddc609bf550485}
\index{External Dependency Resilience (I.2.11)}
{\small Goal · Solution Architect-reviewed Iteration 001; authorised ADD; resulting increment review pending · current}\par

Harmonia aims to explicitly identify and document its external service dependencies for each solution release and, where practical, support continued operation during the unavailability of external services for a defined period, with resilience objectives established according to the requirements and characteristics of each deployment.



Source entry: \textbf{NFR-156 — External Dependency Resilience}. Primary subject: \texttt{\seqsplit{G07.05}}.



The source classification is Goal. The source identifier is a scoped register reference and does not determine architectural type. Applicability, release scope and unresolved thresholds remain as stated above.


{\footnotesize \href{../../../../../../../continuum/motivation/goal/150b3b1eea9b484b89ddc609bf550485.md}{Canonical record: 150b3b1eea9b484b89ddc609bf550485}; full qualifications and provenance.}\par
\chapter*{I.3 Business Outcomes}
\phantomsection
\label{navigation:I.3}
\addcontentsline{toc}{chapter}{I.3 Business Outcomes}

\section*{I.3.1 Continuous and Resilient Regional Health Information Exchange (O2)}
\phantomsection
\label{navigation:I.3.1}
\addcontentsline{toc}{section}{I.3.1 Continuous and Resilient Regional Health Information Exchange (O2)}

\label{element:aed90b4a88f24b78adfeadbf0bca8980}
\index{Continuous and Resilient Regional Health Information Exchange (O2) (I.3.1)}
{\small Outcome · Agreed · current}\par

\par\medskip\textbf{Outcome Definition}\par


Regional healthcare facilities operate within an uninterrupted, highly available, and fault-tolerant digital integration fabric that reliably exchanges clinical events across heterogeneous systems without data loss or performance degradation.



\par\medskip\textbf{Clinical \& Operational Value}\par


\begin{itemize}


\item \textbf{Zero Silent Data Loss}: Admission notifications, emergency referrals, and diagnostic alerts are guaranteed to cross durable boundaries before positive transport acknowledgement, ensuring that integration failures never result in dropped clinical events.

\item \textbf{High Concurrency \& Responsiveness}: Point-of-care clinical queries and vital notification streams remain highly responsive, unblocked by background batch synchronization or database transaction locks.

\item \textbf{Coordinated Inter-Facility Workflows}: Multi-stage clinical activities (e.g., patient transfers, regional lab dispatch, pharmacy updates) progress deterministically across multiple destination systems, preventing divergence of clinical state between facilities.

\item \textbf{Seamless Standards Interoperability}: Seamless bridging of legacy HL7 v2 systems and modern FHIR RESTful repositories without expensive, disruptive rip-and-replace software upgrades.

\end{itemize}


\par\medskip\textbf{Direct Contributing Goals}\par


\begin{itemize}


\item \href{\detokenize{../../../../../../../../docs/markdown/01-motivation/goals-outcomes/strategic-goals.md}}{Vendor-Independent Longitudinal Clinical Information Coherence} (Thread 1)

\item \href{\detokenize{../../../../../../../../docs/markdown/01-motivation/goals-outcomes/strategic-goals.md}}{Durable Acceptance and Preservation of Clinical Events} (Thread 3)

\item \href{\detokenize{../../../../../../../../docs/markdown/01-motivation/goals-outcomes/strategic-goals.md}}{Responsive Access to Managed Information} (Thread 4)

\item \href{\detokenize{../../../../../../../../docs/markdown/01-motivation/goals-outcomes/strategic-goals.md}}{Coordinated Progression of Operational Activities and Associated Entity State} (Thread 7)

\end{itemize}

{\footnotesize \href{../../../../../../../continuum/motivation/outcome/aed90b4a88f24b78adfeadbf0bca8980.md}{Canonical record: aed90b4a88f24b78adfeadbf0bca8980}; full qualifications and provenance.}\par
\section*{I.3.2 Demonstrable Protection and Accountable Handling of Health Information (O3)}
\phantomsection
\label{navigation:I.3.2}
\addcontentsline{toc}{section}{I.3.2 Demonstrable Protection and Accountable Handling of Health Information (O3)}

\label{element:75519ed8f7de4a619cc9065f525d703a}
\index{Demonstrable Protection and Accountable Handling of Health Information (O3) (I.3.2)}
{\small Outcome · Agreed · current}\par

\par\medskip\textbf{Outcome Definition}\par


Sensitive health information is demonstrably safeguarded across all integration operations, backed by default-deny authorization, zero-PHI diagnostic logging, and tamper-evident audit provenance that satisfies rigorous statutory data protection regimes.



\par\medskip\textbf{Clinical \& Operational Value}\par


\begin{itemize}


\item \textbf{Patient Trust \& Confidentiality}: Patients are assured that their intimate health data is protected from unauthorized access, eavesdropping, and accidental leakage into operational logs.

\item \textbf{Statutory Compliance \& Legal Assurance}: The regional health network demonstrably complies with national privacy legislation (e.g., the Australian \emph{Privacy Act 1988}, My Health Records Act, and Notifiable Data Breaches scheme), eliminating legal liability and multi-million-dollar breach penalties.

\item \textbf{Forensic Auditability}: Every state-changing clinical transaction maintains unambiguous, tamper-evident cryptographic provenance (who authorized it, what information was conveyed, when it occurred), enabling rigorous clinical incident retrospectives.

\end{itemize}


\par\medskip\textbf{Relationship to Motivational Principles}\par


Unlike Outcomes O1 and O2 (which receive direct handovers from specific strategic goals), Outcome O3 is an \textbf{intrinsic platform property}:
- It is realized through the foundational architectural principles \hyperref[element:5de2ae2c3b794c2d9f3a14efd51d564f]{AX-07 (Security Is Intrinsic to Managed Operations)} and \hyperref[element:21d6864657964ba5b03e3f6b5eb29cd5]{AX-08 (Evidence Records Meaning)}.
- Security and privacy are enforced intrinsically across all governed ingress, processing, and egress paths rather than being isolated to a single integration workflow.
- As established in the \href{\detokenize{../../../../../../../../docs/markdown/01-motivation/orientation-view.md}}{Orientation View}, AX-07 and AX-08 do not maintain causal arrows to outcomes; O3 acts as a common consequence sink for the entire platform.


{\footnotesize \href{../../../../../../../continuum/motivation/outcome/75519ed8f7de4a619cc9065f525d703a.md}{Canonical record: 75519ed8f7de4a619cc9065f525d703a}; full qualifications and provenance.}\par
\section*{I.3.3 Reduced Clinical Risk from Unavailable, Fragmented or Incorrectly Associated Information (O1)}
\phantomsection
\label{navigation:I.3.3}
\addcontentsline{toc}{section}{I.3.3 Reduced Clinical Risk from Unavailable, Fragmented or Incorrectly Associated Information (O1)}

\label{element:0582153fb1f1402b83e5358b61fd60af}
\index{Reduced Clinical Risk from Unavailable, Fragmented or Incorrectly Associated Information (O1) (I.3.3)}
{\small Outcome · Agreed · current}\par

\par\medskip\textbf{Outcome Definition}\par


Patients experience significantly safer care and reduced preventable harm across the regional healthcare network, achieved through the elimination of clinical data blind spots, the prevention of erroneous subject associations, and the provision of coherent longitudinal patient records.



\par\medskip\textbf{Clinical \& Operational Value}\par


\begin{itemize}


\item \textbf{Prevention of Adverse Medical Events}: Clinicians have point-of-care access to complete allergy lists, current medications, and past medical histories, preventing severe adverse drug events and contraindicated treatments.

\item \textbf{Elimination of Diagnostic Blind Spots}: Hospital emergency departments can immediately view recent laboratory results and diagnostic imaging reports performed at external facilities, avoiding redundant radiation exposure and delays in critical therapy.

\item \textbf{Accurate Subject Attribution}: Clinical records are referentially validated, ensuring that test results and treatment orders are never erroneously attributed to the wrong patient.

\end{itemize}


\par\medskip\textbf{Direct Contributing Goals}\par


\begin{itemize}


\item \href{\detokenize{../../../../../../../../docs/markdown/01-motivation/goals-outcomes/strategic-goals.md}}{Vendor-Independent Longitudinal Clinical Information Coherence} (Thread 2)

\item \href{\detokenize{../../../../../../../../docs/markdown/01-motivation/goals-outcomes/strategic-goals.md}}{Reliable Subject Identity and Referential Integrity} (Thread 5)

\end{itemize}

{\footnotesize \href{../../../../../../../continuum/motivation/outcome/0582153fb1f1402b83e5358b61fd60af.md}{Canonical record: 0582153fb1f1402b83e5358b61fd60af}; full qualifications and provenance.}\par
\chapter*{I.4 Assessments}
\phantomsection
\label{navigation:I.4}
\addcontentsline{toc}{chapter}{I.4 Assessments}

\section*{I.4.1 Centralised Synchronous Persistence Can Constrain Concurrent Processing}
\phantomsection
\label{navigation:I.4.1}
\addcontentsline{toc}{section}{I.4.1 Centralised Synchronous Persistence Can Constrain Concurrent Processing}

\label{element:4f7b18cecb08415bad2b4247068a511e}
\index{Centralised Synchronous Persistence Can Constrain Concurrent Processing (I.4.1)}
{\small Assessment · Agreed · current}\par

\par\medskip\textbf{Assessment Summary}\par


Requiring central synchronous persistence interaction for every unit of active processing can introduce contention, latency, unnecessary I/O and coordination bottlenecks that limit concurrent throughput and scalability.



\par\medskip\textbf{Operational Hazard}\par


In a high-volume regional integration hub, thousands of clinical messages arrive concurrently from emergency rooms, pathology analyzers, and patient portals. If the architecture forces every active query, routing step, and transient coordination check to execute synchronous disk I/O and transactional row locking against a single central database:
- Transaction lock contention and wait states escalate as concurrency increases.
- Increased I/O wait times can cause ingress queues to fill and upstream connections to time out.
- High-priority emergency point-of-care clinical queries become stalled behind long-running background batch imports.
- Scalability is constrained to the physical I/O limits of the central persistence appliance.



\emph{Note on Technology Neutrality}: This assessment diagnoses the fundamental hazards of synchronous persistence contention, latency, and coordination bottlenecks. It is an architectural evaluation of concurrency dynamics; it does not mandate any specific product, cache, database, or queueing technology.



\par\medskip\textbf{Architectural Implication \& Response}\par


\begin{itemize}


\item The architecture must decouple active operational coordination from authoritative durable persistence.

\item High-frequency operational checks and client-facing queries must be served through responsive, concurrent access paths without imposing transactional lock contention on authoritative storage.

\item \textbf{Related Driver}: \href{\detokenize{../../../../../../../../docs/markdown/01-motivation/drivers-assessments/drivers.md}}{Continuous Clinical Service Availability}

\item \textbf{Target Goal}: \href{\detokenize{../../../../../../../../docs/markdown/01-motivation/goals-outcomes/strategic-goals.md}}{Responsive Access to Managed Information}

\item \textbf{Governing Axioms}: \hyperref[element:8a1e995f16954484a78ae1c5c114d546]{AX-11 (Responsive \& Resilient Information Access)}, \hyperref[element:0412a2084f0145babc4f022c71fea65f]{AX-05 (Active Coordination ≠ Authoritative State)}, \hyperref[element:9a5cb9e8876844e58e660f909f7df9fc]{AX-09 (Default Ephemerality of Operational State)}

\end{itemize}

{\footnotesize \href{../../../../../../../continuum/motivation/assessment/4f7b18cecb08415bad2b4247068a511e.md}{Canonical record: 4f7b18cecb08415bad2b4247068a511e}; full qualifications and provenance.}\par
\section*{I.4.2 Implicit Perimeter Trust}
\phantomsection
\label{navigation:I.4.2}
\addcontentsline{toc}{section}{I.4.2 Implicit Perimeter Trust}

\label{element:f1b3b73329ae46b9b5105fb8ff236be5}
\index{Implicit Perimeter Trust (I.4.2)}
{\small Assessment · Agreed · current}\par

\par\medskip\textbf{Assessment Summary}\par


Assuming that internal network segments, container fabrics, or service-to-service calls are trustworthy by default leaves the platform vulnerable to lateral movement, credential abuse, and unauthorized clinical data manipulation.



\par\medskip\textbf{Operational Hazard}\par


Traditional enterprise architectures frequently rely on perimeter firewalls: once a request crosses the external network boundary, all internal service-to-service communication is unauthenticated, unencrypted, and fully trusted.



In modern distributed and containerized integration environments, perimeter-only trust is hazardous:
- Compromise of a single peripheral container or gateway allows an attacker to make unhindered calls to core clinical repositories.
- Internal integration routes can bypass clinical consent rules and authorization checks because the caller is presumed to be "internal."
- Audit logs cannot prove who authorized a specific data retrieval or state mutation because individual operations lack verifiable caller security context.



\par\medskip\textbf{Architectural Implication \& Response}\par


\begin{itemize}


\item Security must be intrinsic to every governed operation rather than dependent upon voluntary caller behaviour or perimeter assumptions.

\item Every internal request, service invocation, and state mutation must execute within an established, verifiable security context subject to platform-enforced security policy.

\item Downstream security architecture (Domain 08) derives specific enforcement patterns, such as default-deny authorization and authenticated service communication, from governing axiom AX-07.

\item \textbf{Related Driver}: \href{\detokenize{../../../../../../../../docs/markdown/01-motivation/drivers-assessments/drivers.md}}{Statutory Health Information Privacy \& Protection}

\item \textbf{Governing Axiom}: \hyperref[element:5de2ae2c3b794c2d9f3a14efd51d564f]{AX-07 (Security Is Intrinsic to Managed Operations)}

\item \textbf{Downstream Requirements}: \texttt{\seqsplit{SEC-001 / SEC-002}} (Trusted Authentication \& Default-Deny Authorization).

\end{itemize}

{\footnotesize \href{../../../../../../../continuum/motivation/assessment/f1b3b73329ae46b9b5105fb8ff236be5.md}{Canonical record: f1b3b73329ae46b9b5105fb8ff236be5}; full qualifications and provenance.}\par
\section*{I.4.3 PHI Leakage through Operational Logging}
\phantomsection
\label{navigation:I.4.3}
\addcontentsline{toc}{section}{I.4.3 PHI Leakage through Operational Logging}

\label{element:6d8848b700df4f6eacc0347993ac541a}
\index{PHI Leakage through Operational Logging (I.4.3)}
{\small Assessment · Agreed · current}\par

\par\medskip\textbf{Assessment Summary}\par


Emitting unmasked Protected Health Information (PHI) into operational diagnostic logs, monitoring systems, and tracing tools creates severe data breach exposures and violates statutory privacy obligations.



\par\medskip\textbf{Operational Hazard}\par


Software developers and integration engineers require diagnostic logging to troubleshoot routing errors, transform bugs, and network timeouts. A common developer practice is logging raw message payloads (e.g., logging an entire raw HL7 message string or complete FHIR resource JSON) at \texttt{\seqsplit{DEBUG}} or \texttt{\seqsplit{INFO}} levels.



In a healthcare environment, these logs are ingested into centralized search clusters, log forwarders, monitoring dashboards, and cloud observability tools:
- Patient names, home addresses, phone numbers, Medicare numbers, and intimate clinical diagnoses become stored in unencrypted, broadly accessible log repositories.
- Operational support personnel and external vendors gain unauthorized access to sensitive clinical data without clinical need-to-know authorization.
- Routine log retention and aggregation triggers substantial statutory data breach notifications under national privacy legislation (e.g., the Australian \emph{Privacy Act 1988} and Notifiable Data Breaches scheme).



\par\medskip\textbf{Architectural Implication \& Response}\par


\begin{itemize}


\item Operational log streams must be structurally segregated from clinical data payloads.

\item Diagnostic logs must emit only non-sensitive operational metadata (transaction UUIDs, message control IDs, timestamps, masked identifiers) and strictly prohibit unmasked clinical observation values, demographic text, or free-text clinical notes.

\item \textbf{Related Driver}: \href{\detokenize{../../../../../../../../docs/markdown/01-motivation/drivers-assessments/drivers.md}}{Statutory Health Information Privacy \& Protection}

\item \textbf{Governing Axiom}: \hyperref[element:5de2ae2c3b794c2d9f3a14efd51d564f]{AX-07 (Security Is Intrinsic to Managed Operations)}

\item \textbf{Derived Guardrail \& Requirement}: Classifies downstream as a mandatory Security Architecture requirement (\texttt{\seqsplit{SEC-010}} PHI-Safe Operational Logging), with formal governance guardrails owned under Domain 13.

\end{itemize}

{\footnotesize \href{../../../../../../../continuum/motivation/assessment/6d8848b700df4f6eacc0347993ac541a.md}{Canonical record: 6d8848b700df4f6eacc0347993ac541a}; full qualifications and provenance.}\par
\section*{I.4.4 Silent Data Loss via False Acceptance}
\phantomsection
\label{navigation:I.4.4}
\addcontentsline{toc}{section}{I.4.4 Silent Data Loss via False Acceptance}

\label{element:9f185cd26a0f4118a0b28bb08134f472}
\index{Silent Data Loss via False Acceptance (I.4.4)}
{\small Assessment · Agreed · current}\par

\par\medskip\textbf{Assessment Summary}\par


Prematurely acknowledging receipt of an incoming clinical event before responsibility for that event has crossed an authoritative durable boundary creates a critical risk of unrecoverable data loss.



\par\medskip\textbf{Operational Hazard}\par


When an ingress component positively acknowledges an incoming operational or clinical event before responsibility for that event has crossed an authoritative durable boundary, upstream senders legitimately consider their delivery obligation satisfied and release local responsibility.



If an unrecoverable failure occurs (such as process termination, node failure, network partition, or power loss) after positive acknowledgement but before durable persistence or guaranteed downstream delivery:
- The upstream sender marks the event as delivered and will not retransmit.
- The receiving platform loses the event representation without having persisted or processed it.
- The clinical event is silently dropped without any error indication or notification.



In clinical practice, this leads directly to missing critical diagnostic results, unrecorded medication orders, or undetected admission events, creating severe risks for patient safety.



\par\medskip\textbf{Downstream Integration Architecture Example}\par


In traditional HL7 v2 and MLLP transport middleware, inbound gateways often return an immediate positive acknowledgement (\texttt{\seqsplit{MSA-1 = AA}}, Application Accept) upon successful memory parsing or volatile buffering. If the gateway process crashes or hardware fails before the message is committed to durable storage, silent data loss occurs because the sender received an \texttt{\seqsplit{AA}} and will never retry. Downstream integration adapters (Domain 06) must enforce durable acceptance boundaries and emit negative acknowledgements (\texttt{\seqsplit{AE}}) when durability is not established.



\par\medskip\textbf{Architectural Implication \& Response}\par


\begin{itemize}


\item Positive acceptance must be strictly coupled to crossing an authoritative durable boundary.

\item Ingress boundaries must guarantee that incoming event responsibility is durably accepted before emitting a positive acceptance response to the sender. If durable acceptance cannot be established, an explicit negative or indeterminate response must be emitted so the sender retains retry responsibility.

\item \textbf{Related Driver}: \href{\detokenize{../../../../../../../../docs/markdown/01-motivation/drivers-assessments/drivers.md}}{Continuous Clinical Service Availability}

\item \textbf{Target Goal}: \href{\detokenize{../../../../../../../../docs/markdown/01-motivation/goals-outcomes/strategic-goals.md}}{Durable Acceptance \& Preservation of Clinical Events}

\item \textbf{Governing Axioms}: \hyperref[element:0d9673833eb840eab03804aa076c1b0c]{AX-10 (Failure as Normal Condition)}, \hyperref[element:0412a2084f0145babc4f022c71fea65f]{AX-05 (State Separation)}, \hyperref[element:ef8abbaec37e4890ad78b9f263e67906]{AX-15 (Explicit Representation of Uncertainty)}

\item \textbf{Foundational Requirement}: \href{\detokenize{../../../../../../../../docs/markdown/01-motivation/requirements-constraints/foundational-requirements.md}}{\texttt{\seqsplit{REQ-FND-001}} Durable Ingress Acceptance Boundary}

\end{itemize}

{\footnotesize \href{../../../../../../../continuum/motivation/assessment/9f185cd26a0f4118a0b28bb08134f472.md}{Canonical record: 9f185cd26a0f4118a0b28bb08134f472}; full qualifications and provenance.}\par
\section*{I.4.5 Unmonitored Destination Failure / Fan-Out Divergence}
\phantomsection
\label{navigation:I.4.5}
\addcontentsline{toc}{section}{I.4.5 Unmonitored Destination Failure / Fan-Out Divergence}

\label{element:3357fcd2f7ce457386e33a2f06ae2182}
\index{Unmonitored Destination Failure / Fan-Out Divergence (I.4.5)}
{\small Assessment · Agreed · current}\par

\par\medskip\textbf{Assessment Summary}\par


Broadcasting clinical events across multiple downstream destination systems without tracking granular, per-destination delivery status leads to silent divergence between destination systems and fractured clinical state.



\par\medskip\textbf{Operational Hazard}\par


A single clinical event often requires distribution to multiple disparate destinations (fan-out). For example, an HL7 \texttt{\seqsplit{ADT\textasciicircum{}A08}} (Update Patient Information) or \texttt{\seqsplit{ADT\textasciicircum{}A02}} (Transfer Patient) may need to update:
1. The regional clinical data repository.
2. The local pharmacy dispensing system.
3. The radiology information system (RIS).
4. The hospital billing and claims engine.
5. An external primary health network notifications feed.



If the integration platform considers the parent activity "completed" merely because the message was posted to an internal dispatch mechanism, failures at individual destination endpoints go unmonitored:
- The clinical repository is updated, but the pharmacy feed times out.
- The parent workflow reports success.
- The pharmacy continues operating on obsolete patient medication allergy data.
- Systemic clinical state diverges silently across the healthcare network.



\par\medskip\textbf{Architectural Implication \& Response}\par


\begin{itemize}


\item Operational workflows must track explicit progression state per destination.

\item Parent activities must maintain fine-grained sub-status across the entire fan-out lifecycle until every individual destination delivery is deterministically confirmed or explicitly flagged as failed/indeterminate.

\item \textbf{Related Driver}: \href{\detokenize{../../../../../../../../docs/markdown/01-motivation/drivers-assessments/drivers.md}}{Coordinated Operational Activity and Entity State}

\item \textbf{Target Goal}: \href{\detokenize{../../../../../../../../docs/markdown/01-motivation/goals-outcomes/strategic-goals.md}}{Coordinated Progression of Operational Activities \& Associated Entity State}

\item \textbf{Governing Axiom}: \hyperref[element:de7e90f46127455a9ec9ad3b9b48ba54]{AX-16 (Operational Activity \& Entity State Progress Together)}

\item \textbf{Foundational Requirement}: \href{\detokenize{../../../../../../../../docs/markdown/01-motivation/requirements-constraints/foundational-requirements.md}}{\texttt{\seqsplit{REQ-FND-002}} Operational Activity Progression State}

\end{itemize}

{\footnotesize \href{../../../../../../../continuum/motivation/assessment/3357fcd2f7ce457386e33a2f06ae2182.md}{Canonical record: 3357fcd2f7ce457386e33a2f06ae2182}; full qualifications and provenance.}\par
\chapter*{I.5 Requirements}
\phantomsection
\label{navigation:I.5}
\addcontentsline{toc}{chapter}{I.5 Requirements}

\section*{I.5.1 AI Assurance Architecture Discovery Before Implementation}
\phantomsection
\label{navigation:I.5.1}
\addcontentsline{toc}{section}{I.5.1 AI Assurance Architecture Discovery Before Implementation}

\label{element:24a609117a4f40839f8b669ebf720cde}
\index{AI Assurance Architecture Discovery Before Implementation (I.5.1)}
{\small Requirement · Adjudicated Domain05 direction — U02/U06/U19; holistic result acceptance pending · current}\par

Domain05 must undertake explicit Information Systems Layer Architecture discovery before implementation models are finalised, including: AI model/service selection and abstraction; context/input construction; information disclosure/security; model/service provenance and versioning; confidence/uncertainty/result representation; deterministic versus non-deterministic behaviour; reproducibility; result validation; human review/escalation; assurance evidence representation; workflow integration; failure/timeout/unavailability behaviour; and reporting/observability consistent with Dokimasia's terminal-boundary rules.



This assertion is normalised from the governed DesignNote under U02/U06/U19 and the ARB release Praxis. Original candidate/release qualifications and separately deferred implementation details remain in force.



Architectural subject: \hyperref[element:fb3f37a6ac614415ac486a6de9d35994]{Energeia}.


{\footnotesize \href{../../../../../../../continuum/motivation/requirement/24a609117a4f40839f8b669ebf720cde.md}{Canonical record: 24a609117a4f40839f8b669ebf720cde}; full qualifications and provenance.}\par
\section*{I.5.2 Artemis HA Failover}
\phantomsection
\label{navigation:I.5.2}
\addcontentsline{toc}{section}{I.5.2 Artemis HA Failover}

\label{element:8c0f08d0d12f48b9ae6a001c19ede377}
\index{Artemis HA Failover (I.5.2)}
{\small Requirement · Historical — current applicability unresolved · historical\_reference}\par
\href{../../../../../../../continuum/motivation/requirement/8c0f08d0d12f48b9ae6a001c19ede377.md}{Historical canonical definition and provenance}.\par
\section*{I.5.3 Audit Corrections Create New Linked Events}
\phantomsection
\label{navigation:I.5.3}
\addcontentsline{toc}{section}{I.5.3 Audit Corrections Create New Linked Events}

\label{element:ddaa9bff2b684af7a6ad18fe4f2da8ad}
\index{Audit Corrections Create New Linked Events (I.5.3)}
{\small Requirement · Adjudicated Domain05 direction — U07/U15; holistic result acceptance pending · current}\par

Where an earlier AuditEvent is subsequently determined to be incorrect, incomplete, misleading, or otherwise requires qualification, Kleio SHALL create a \textbf{new AuditEvent} that refers back to the earlier AuditEvent and records the subsequent assertion.



This assertion is normalised from the governed DesignNote under U07/U15 and the ARB release Praxis. Original candidate/release qualifications and separately deferred implementation details remain in force.



Architectural subject: \hyperref[element:fefda110e2274884ad58906733606efd]{Kleio}.


{\footnotesize \href{../../../../../../../continuum/motivation/requirement/ddaa9bff2b684af7a6ad18fe4f2da8ad.md}{Canonical record: ddaa9bff2b684af7a6ad18fe4f2da8ad}; full qualifications and provenance.}\par
\section*{I.5.4 CLIN-001 — Vendor-Neutral Clinical Information Layer}
\phantomsection
\label{navigation:I.5.4}
\addcontentsline{toc}{section}{I.5.4 CLIN-001 — Vendor-Neutral Clinical Information Layer}

\label{element:44eec8cf17f44593824a1f723d949352}
\index{CLIN-001 — Vendor-Neutral Clinical Information Layer (I.5.4)}
{\small Requirement · Historical — current owning-domain standing not inferred · historical\_reference}\par
\href{../../../../../../../continuum/motivation/requirement/44eec8cf17f44593824a1f723d949352.md}{Historical canonical definition and provenance}.\par
\section*{I.5.5 CLIN-002 — FHIR-Based Longitudinal Clinical Representation}
\phantomsection
\label{navigation:I.5.5}
\addcontentsline{toc}{section}{I.5.5 CLIN-002 — FHIR-Based Longitudinal Clinical Representation}

\label{element:9ad88c6b7f934923800dd46c8522b9a1}
\index{CLIN-002 — FHIR-Based Longitudinal Clinical Representation (I.5.5)}
{\small Requirement · Historical — current owning-domain standing not inferred · historical\_reference}\par
\href{../../../../../../../continuum/motivation/requirement/9ad88c6b7f934923800dd46c8522b9a1.md}{Historical canonical definition and provenance}.\par
\section*{I.5.6 CLIN-003 — Medication as a First-Class Clinical Domain}
\phantomsection
\label{navigation:I.5.6}
\addcontentsline{toc}{section}{I.5.6 CLIN-003 — Medication as a First-Class Clinical Domain}

\label{element:f5a6dd8f23f047bb9f4c327bbd3e85f2}
\index{CLIN-003 — Medication as a First-Class Clinical Domain (I.5.6)}
{\small Requirement · Historical — current owning-domain standing not inferred · historical\_reference}\par
\href{../../../../../../../continuum/motivation/requirement/f5a6dd8f23f047bb9f4c327bbd3e85f2.md}{Historical canonical definition and provenance}.\par
\section*{I.5.7 CLIN-004 — Source-System Independence}
\phantomsection
\label{navigation:I.5.7}
\addcontentsline{toc}{section}{I.5.7 CLIN-004 — Source-System Independence}

\label{element:a5a3d2ddb02d4ac4b3eabdfb2ecd22bb}
\index{CLIN-004 — Source-System Independence (I.5.7)}
{\small Requirement · Historical — current owning-domain standing not inferred · historical\_reference}\par
\href{../../../../../../../continuum/motivation/requirement/a5a3d2ddb02d4ac4b3eabdfb2ecd22bb.md}{Historical canonical definition and provenance}.\par
\section*{I.5.8 CLIN-005 — Source Authority Preservation}
\phantomsection
\label{navigation:I.5.8}
\addcontentsline{toc}{section}{I.5.8 CLIN-005 — Source Authority Preservation}

\label{element:5f399a338688438c922df969e724a24a}
\index{CLIN-005 — Source Authority Preservation (I.5.8)}
{\small Requirement · Historical — current owning-domain standing not inferred · historical\_reference}\par
\href{../../../../../../../continuum/motivation/requirement/5f399a338688438c922df969e724a24a.md}{Historical canonical definition and provenance}.\par
\section*{I.5.9 CLIN-006 — Authority-Aware Clinical Reconciliation}
\phantomsection
\label{navigation:I.5.9}
\addcontentsline{toc}{section}{I.5.9 CLIN-006 — Authority-Aware Clinical Reconciliation}

\label{element:3d7c709a435145e3821184b2c9b980fa}
\index{CLIN-006 — Authority-Aware Clinical Reconciliation (I.5.9)}
{\small Requirement · Historical — current owning-domain standing not inferred · historical\_reference}\par
\href{../../../../../../../continuum/motivation/requirement/3d7c709a435145e3821184b2c9b980fa.md}{Historical canonical definition and provenance}.\par
\section*{I.5.10 CLIN-007 — Preservation of Conflicting Assertions}
\phantomsection
\label{navigation:I.5.10}
\addcontentsline{toc}{section}{I.5.10 CLIN-007 — Preservation of Conflicting Assertions}

\label{element:cd669acb9a76499db70b5d092917a56b}
\index{CLIN-007 — Preservation of Conflicting Assertions (I.5.10)}
{\small Requirement · Historical — current owning-domain standing not inferred · historical\_reference}\par
\href{../../../../../../../continuum/motivation/requirement/cd669acb9a76499db70b5d092917a56b.md}{Historical canonical definition and provenance}.\par
\section*{I.5.11 CLIN-008 — Clinical Source Context}
\phantomsection
\label{navigation:I.5.11}
\addcontentsline{toc}{section}{I.5.11 CLIN-008 — Clinical Source Context}

\label{element:591d9b8007504d2382ae7977a6589742}
\index{CLIN-008 — Clinical Source Context (I.5.11)}
{\small Requirement · Historical — current owning-domain standing not inferred · historical\_reference}\par
\href{../../../../../../../continuum/motivation/requirement/591d9b8007504d2382ae7977a6589742.md}{Historical canonical definition and provenance}.\par
\section*{I.5.12 CLIN-009 — Governed Patient Identity}
\phantomsection
\label{navigation:I.5.12}
\addcontentsline{toc}{section}{I.5.12 CLIN-009 — Governed Patient Identity}

\label{element:4dded2374b9346e4b8d0e326e7d9c97c}
\index{CLIN-009 — Governed Patient Identity (I.5.12)}
{\small Requirement · Historical — current owning-domain standing not inferred · historical\_reference}\par
\href{../../../../../../../continuum/motivation/requirement/4dded2374b9346e4b8d0e326e7d9c97c.md}{Historical canonical definition and provenance}.\par
\section*{I.5.13 CLIN-010 — Read-Oriented Initial LHR}
\phantomsection
\label{navigation:I.5.13}
\addcontentsline{toc}{section}{I.5.13 CLIN-010 — Read-Oriented Initial LHR}

\label{element:96dad3bf0c764ca0849227eae64848aa}
\index{CLIN-010 — Read-Oriented Initial LHR (I.5.13)}
{\small Requirement · Historical — current owning-domain standing not inferred · historical\_reference}\par
\href{../../../../../../../continuum/motivation/requirement/96dad3bf0c764ca0849227eae64848aa.md}{Historical canonical definition and provenance}.\par
\section*{I.5.14 CLIN-011 — Selective Clinical Authoring}
\phantomsection
\label{navigation:I.5.14}
\addcontentsline{toc}{section}{I.5.14 CLIN-011 — Selective Clinical Authoring}

\label{element:bd593ed8f3b3465ca499151f18767585}
\index{CLIN-011 — Selective Clinical Authoring (I.5.14)}
{\small Requirement · Historical — current owning-domain standing not inferred · historical\_reference}\par
\href{../../../../../../../continuum/motivation/requirement/bd593ed8f3b3465ca499151f18767585.md}{Historical canonical definition and provenance}.\par
\section*{I.5.15 CORE-001 — Information Authority Classification}
\phantomsection
\label{navigation:I.5.15}
\addcontentsline{toc}{section}{I.5.15 CORE-001 — Information Authority Classification}

\label{element:6e76bb1db8f2464496355d47d2f981ef}
\index{CORE-001 — Information Authority Classification (I.5.15)}
{\small Requirement · Historical — current owning-domain standing not inferred · historical\_reference}\par
\href{../../../../../../../continuum/motivation/requirement/6e76bb1db8f2464496355d47d2f981ef.md}{Historical canonical definition and provenance}.\par
\section*{I.5.16 CORE-002 — Separation of Architectural Concerns}
\phantomsection
\label{navigation:I.5.16}
\addcontentsline{toc}{section}{I.5.16 CORE-002 — Separation of Architectural Concerns}

\label{element:a5c09af151984e989590b9198e2def70}
\index{CORE-002 — Separation of Architectural Concerns (I.5.16)}
{\small Requirement · Historical — current owning-domain standing not inferred · historical\_reference}\par
\href{../../../../../../../continuum/motivation/requirement/a5c09af151984e989590b9198e2def70.md}{Historical canonical definition and provenance}.\par
\section*{I.5.17 CORE-003 — Governed Canonical Semantics}
\phantomsection
\label{navigation:I.5.17}
\addcontentsline{toc}{section}{I.5.17 CORE-003 — Governed Canonical Semantics}

\label{element:006383a5cce64780960fda32b6682d8d}
\index{CORE-003 — Governed Canonical Semantics (I.5.17)}
{\small Requirement · Historical — current owning-domain standing not inferred · historical\_reference}\par
\href{../../../../../../../continuum/motivation/requirement/006383a5cce64780960fda32b6682d8d.md}{Historical canonical definition and provenance}.\par
\section*{I.5.18 [AssuranceRequirement] CORE-004 — Correlation and Causation}
\phantomsection
\label{navigation:I.5.18}
\addcontentsline{toc}{section}{I.5.18 [AssuranceRequirement] CORE-004 — Correlation and Causation}

\label{element:f2e504e7d6374e41a81303265adda6de}
\index{[AssuranceRequirement] CORE-004 — Correlation and Causation (I.5.18)}
{\small Requirement / AssuranceRequirement · Solution Architect-reviewed Iteration 001; authorised ADD; resulting increment review pending · current}\par

\par\medskip\textbf{Accepted Iteration 001 contribution}\par


\textbf{NFR-058 — Cross-Service Activity Correlation}



Correlate processing activity across services



Primary subject: \texttt{\seqsplit{G09.04}}. Native Requirement; Harmonia specialisation AssuranceRequirement.



\par\medskip\textbf{Preserved historical contribution}\par


Original semantic standing: Historical — current owning-domain standing not inferred. Original currentness: \texttt{\seqsplit{historical\_reference}}. The following original content remains historical; the new contribution does not independently approve its additional details.



Harmonia shall retain sufficient correlation and causation information to relate an originating interaction to subsequent messages, tasks, workflow activities, persistence operations, and egress interactions.


{\footnotesize \href{../../../../../../../continuum/motivation/requirement/assurance-requirement/f2e504e7d6374e41a81303265adda6de.md}{Canonical record: f2e504e7d6374e41a81303265adda6de}; full qualifications and provenance.}\par
\section*{I.5.19 DATA-001 — Durable Governed State}
\phantomsection
\label{navigation:I.5.19}
\addcontentsline{toc}{section}{I.5.19 DATA-001 — Durable Governed State}

\label{element:15c71ccd5d714503891aece2a23445e1}
\index{DATA-001 — Durable Governed State (I.5.19)}
{\small Requirement · Historical — current owning-domain standing not inferred · historical\_reference}\par
\href{../../../../../../../continuum/motivation/requirement/15c71ccd5d714503891aece2a23445e1.md}{Historical canonical definition and provenance}.\par
\section*{I.5.20 DATA-002 — Authoritative Read Semantics}
\phantomsection
\label{navigation:I.5.20}
\addcontentsline{toc}{section}{I.5.20 DATA-002 — Authoritative Read Semantics}

\label{element:bc55838493c24617823410798c8d7598}
\index{DATA-002 — Authoritative Read Semantics (I.5.20)}
{\small Requirement · Historical — current owning-domain standing not inferred · historical\_reference}\par
\href{../../../../../../../continuum/motivation/requirement/bc55838493c24617823410798c8d7598.md}{Historical canonical definition and provenance}.\par
\section*{I.5.21 DATA-003 — Cache Transparency}
\phantomsection
\label{navigation:I.5.21}
\addcontentsline{toc}{section}{I.5.21 DATA-003 — Cache Transparency}

\label{element:ff04c3dab8a04359b9c4d62939839be7}
\index{DATA-003 — Cache Transparency (I.5.21)}
{\small Requirement · Historical — current owning-domain standing not inferred · historical\_reference}\par
\href{../../../../../../../continuum/motivation/requirement/ff04c3dab8a04359b9c4d62939839be7.md}{Historical canonical definition and provenance}.\par
\section*{I.5.22 DATA-004 — Search Scalability}
\phantomsection
\label{navigation:I.5.22}
\addcontentsline{toc}{section}{I.5.22 DATA-004 — Search Scalability}

\label{element:f27e74b6cb13454aa4339932436971e3}
\index{DATA-004 — Search Scalability (I.5.22)}
{\small Requirement · Historical — current owning-domain standing not inferred · historical\_reference}\par
\href{../../../../../../../continuum/motivation/requirement/f27e74b6cb13454aa4339932436971e3.md}{Historical canonical definition and provenance}.\par
\section*{I.5.23 DATA-005 — Concurrency and Version Integrity}
\phantomsection
\label{navigation:I.5.23}
\addcontentsline{toc}{section}{I.5.23 DATA-005 — Concurrency and Version Integrity}

\label{element:7f40b2f4947845d287564b43f812b6c8}
\index{DATA-005 — Concurrency and Version Integrity (I.5.23)}
{\small Requirement · Historical — current owning-domain standing not inferred · historical\_reference}\par
\href{../../../../../../../continuum/motivation/requirement/7f40b2f4947845d287564b43f812b6c8.md}{Historical canonical definition and provenance}.\par
\section*{I.5.24 DATA-006 — Referential Integrity}
\phantomsection
\label{navigation:I.5.24}
\addcontentsline{toc}{section}{I.5.24 DATA-006 — Referential Integrity}

\label{element:4a8a79f520394cfe863f473ec2a31b34}
\index{DATA-006 — Referential Integrity (I.5.24)}
{\small Requirement · Historical — current owning-domain standing not inferred · historical\_reference}\par
\href{../../../../../../../continuum/motivation/requirement/4a8a79f520394cfe863f473ec2a31b34.md}{Historical canonical definition and provenance}.\par
\section*{I.5.25 DATA-007 — Persistence Failure Safety}
\phantomsection
\label{navigation:I.5.25}
\addcontentsline{toc}{section}{I.5.25 DATA-007 — Persistence Failure Safety}

\label{element:40ed2565d0d1467cbbfcf5706924de94}
\index{DATA-007 — Persistence Failure Safety (I.5.25)}
{\small Requirement · Historical — current owning-domain standing not inferred · historical\_reference}\par
\href{../../../../../../../continuum/motivation/requirement/40ed2565d0d1467cbbfcf5706924de94.md}{Historical canonical definition and provenance}.\par
\section*{I.5.26 DATA-008 — Recovery Fidelity}
\phantomsection
\label{navigation:I.5.26}
\addcontentsline{toc}{section}{I.5.26 DATA-008 — Recovery Fidelity}

\label{element:7c7a33be31d042459adb83c71ae044fa}
\index{DATA-008 — Recovery Fidelity (I.5.26)}
{\small Requirement · Historical — current owning-domain standing not inferred · historical\_reference}\par
\href{../../../../../../../continuum/motivation/requirement/7c7a33be31d042459adb83c71ae044fa.md}{Historical canonical definition and provenance}.\par
\section*{I.5.27 Development Representations of Runtime Governed Datasets}
\phantomsection
\label{navigation:I.5.27}
\addcontentsline{toc}{section}{I.5.27 Development Representations of Runtime Governed Datasets}

\label{element:0a4a5c73048c4e49a82d40e349e70a95}
\index{Development Representations of Runtime Governed Datasets (I.5.27)}
{\small Requirement · Adjudicated Domain05 direction — U13/U15; holistic result acceptance pending · current}\par

\begin{quote}


\textbf{Calliope SHALL provide development-time representations of known
runtime-managed DataSets, CodeSets and ValueSets sufficient to support
efficient human, IDE and AI-assisted software development, while
preserving runtime resolution against authoritative Calliope-DataSets
state.}



\end{quote}


This assertion is normalised from the governed DesignNote under U13/U15 and the ARB release Praxis. Original candidate/release qualifications and separately deferred implementation details remain in force.



Architectural subject: \hyperref[element:466ad8f75a954c0581b5bba090443849]{Calliope}.


{\footnotesize \href{../../../../../../../continuum/motivation/requirement/0a4a5c73048c4e49a82d40e349e70a95.md}{Canonical record: 0a4a5c73048c4e49a82d40e349e70a95}; full qualifications and provenance.}\par
\section*{I.5.28 DIR-001 — Standards-Based Provider Directory}
\phantomsection
\label{navigation:I.5.28}
\addcontentsline{toc}{section}{I.5.28 DIR-001 — Standards-Based Provider Directory}

\label{element:c1b1d82f8a9b4c649a3216c5e09a22c6}
\index{DIR-001 — Standards-Based Provider Directory (I.5.28)}
{\small Requirement · Historical — current owning-domain standing not inferred · historical\_reference}\par
\href{../../../../../../../continuum/motivation/requirement/c1b1d82f8a9b4c649a3216c5e09a22c6.md}{Historical canonical definition and provenance}.\par
\section*{I.5.29 DIR-002 — Provider Directory Discovery}
\phantomsection
\label{navigation:I.5.29}
\addcontentsline{toc}{section}{I.5.29 DIR-002 — Provider Directory Discovery}

\label{element:27998c2a8f2e497a814b32e7a0bc5e03}
\index{DIR-002 — Provider Directory Discovery (I.5.29)}
{\small Requirement · Historical — current owning-domain standing not inferred · historical\_reference}\par
\href{../../../../../../../continuum/motivation/requirement/27998c2a8f2e497a814b32e7a0bc5e03.md}{Historical canonical definition and provenance}.\par
\section*{I.5.30 DIR-003 — Governed Directory Change}
\phantomsection
\label{navigation:I.5.30}
\addcontentsline{toc}{section}{I.5.30 DIR-003 — Governed Directory Change}

\label{element:934462c11133409cb3aa59bf2502fa81}
\index{DIR-003 — Governed Directory Change (I.5.30)}
{\small Requirement · Historical — current owning-domain standing not inferred · historical\_reference}\par
\href{../../../../../../../continuum/motivation/requirement/934462c11133409cb3aa59bf2502fa81.md}{Historical canonical definition and provenance}.\par
\section*{I.5.31 DIR-004 — Directory Referential Integrity}
\phantomsection
\label{navigation:I.5.31}
\addcontentsline{toc}{section}{I.5.31 DIR-004 — Directory Referential Integrity}

\label{element:f866ae5adb08499794331bc600db827c}
\index{DIR-004 — Directory Referential Integrity (I.5.31)}
{\small Requirement · Historical — current owning-domain standing not inferred · historical\_reference}\par
\href{../../../../../../../continuum/motivation/requirement/f866ae5adb08499794331bc600db827c.md}{Historical canonical definition and provenance}.\par
\section*{I.5.32 DIR-005 — Healthcare Identifier Governance}
\phantomsection
\label{navigation:I.5.32}
\addcontentsline{toc}{section}{I.5.32 DIR-005 — Healthcare Identifier Governance}

\label{element:097368b4b0b4468dbf1c57c11d9ed86c}
\index{DIR-005 — Healthcare Identifier Governance (I.5.32)}
{\small Requirement · Historical — current owning-domain standing not inferred · historical\_reference}\par
\href{../../../../../../../continuum/motivation/requirement/097368b4b0b4468dbf1c57c11d9ed86c.md}{Historical canonical definition and provenance}.\par
\section*{I.5.33 DIR-006 — Provider Self-Service Governance}
\phantomsection
\label{navigation:I.5.33}
\addcontentsline{toc}{section}{I.5.33 DIR-006 — Provider Self-Service Governance}

\label{element:4230f93461a4470383fb99fddc35f8bf}
\index{DIR-006 — Provider Self-Service Governance (I.5.33)}
{\small Requirement · Historical — current owning-domain standing not inferred · historical\_reference}\par
\href{../../../../../../../continuum/motivation/requirement/4230f93461a4470383fb99fddc35f8bf.md}{Historical canonical definition and provenance}.\par
\section*{I.5.34 Efficient Materialisation of Anticipated Context Resources}
\phantomsection
\label{navigation:I.5.34}
\addcontentsline{toc}{section}{I.5.34 Efficient Materialisation of Anticipated Context Resources}

\label{element:1f3431ca25284783ac69f8f6affcd103}
\index{Efficient Materialisation of Anticipated Context Resources (I.5.34)}
{\small Requirement · Adjudicated Domain05 direction — U10/U18; holistic result acceptance pending · current}\par

Mneme SHALL support efficient materialisation of the anticipated working resource set declared by the Ergo definition.



This assertion is normalised from the governed DesignNote under U10/U18 and the ARB release Praxis. Original candidate/release qualifications and separately deferred implementation details remain in force.



Architectural subject: \hyperref[element:1af8a56358f7429fabfb57be66542a99]{Mneme}.


{\footnotesize \href{../../../../../../../continuum/motivation/requirement/1f3431ca25284783ac69f8f6affcd103.md}{Canonical record: 1f3431ca25284783ac69f8f6affcd103}; full qualifications and provenance.}\par
\section*{I.5.35 Establish Effective Security Attributes through Governed Rules}
\phantomsection
\label{navigation:I.5.35}
\addcontentsline{toc}{section}{I.5.35 Establish Effective Security Attributes through Governed Rules}

\label{element:d14061c03b9042948ba33cae91e02efb}
\index{Establish Effective Security Attributes through Governed Rules (I.5.35)}
{\small Requirement · Adjudicated Domain05 direction — U12/U15/U16; holistic result acceptance pending · current}\par

Where required Information.Compartment and/or Information.Classification attributes are absent, the applicable enforcement mechanism SHALL establish effective values through governed, deterministic rules.



This assertion is normalised from the governed DesignNote under U12/U15/U16 and the ARB release Praxis. Original candidate/release qualifications and separately deferred implementation details remain in force.



Architectural subject: \hyperref[element:899d34ec491948b68aad09e507dea6d1]{Themis}.


{\footnotesize \href{../../../../../../../continuum/motivation/requirement/d14061c03b9042948ba33cae91e02efb.md}{Canonical record: d14061c03b9042948ba33cae91e02efb}; full qualifications and provenance.}\par
\section*{I.5.36 Explain Effective PEP Criteria Derivation}
\phantomsection
\label{navigation:I.5.36}
\addcontentsline{toc}{section}{I.5.36 Explain Effective PEP Criteria Derivation}

\label{element:f027de050273400eaac0de083b820238}
\index{Explain Effective PEP Criteria Derivation (I.5.36)}
{\small Requirement · Adjudicated Domain05 direction — U12/U15/U16; holistic result acceptance pending · current}\par

Themis SHALL retain sufficient provenance to explain the derivation of effective PEP-Criteria, including the authoritative Criteria, Criteria collections, AuthorisationRole assignments, direct assignments and applicable policy state from which authorisation was derived.



This assertion is normalised from the governed DesignNote under U12/U15/U16 and the ARB release Praxis. Original candidate/release qualifications and separately deferred implementation details remain in force.



Architectural subject: \hyperref[element:899d34ec491948b68aad09e507dea6d1]{Themis}.


{\footnotesize \href{../../../../../../../continuum/motivation/requirement/f027de050273400eaac0de083b820238.md}{Canonical record: f027de050273400eaac0de083b820238}; full qualifications and provenance.}\par
\section*{I.5.37 FHIR R5 Schemas}
\phantomsection
\label{navigation:I.5.37}
\addcontentsline{toc}{section}{I.5.37 FHIR R5 Schemas}

\label{element:8d8c62b8cb7b40b0be2ce23f7f41f154}
\index{FHIR R5 Schemas (I.5.37)}
{\small Requirement · Historical / Reclassified — Domain04/06 · historical\_reference}\par
\href{../../../../../../../continuum/motivation/requirement/8d8c62b8cb7b40b0be2ce23f7f41f154.md}{Historical canonical definition and provenance}.\par
\section*{I.5.38 INT-001 — Standards-Based Healthcare Interoperability}
\phantomsection
\label{navigation:I.5.38}
\addcontentsline{toc}{section}{I.5.38 INT-001 — Standards-Based Healthcare Interoperability}

\label{element:44dba056ada04f6583b77859aac9ac68}
\index{INT-001 — Standards-Based Healthcare Interoperability (I.5.38)}
{\small Requirement · Historical — current owning-domain standing not inferred · historical\_reference}\par
\href{../../../../../../../continuum/motivation/requirement/44dba056ada04f6583b77859aac9ac68.md}{Historical canonical definition and provenance}.\par
\section*{I.5.39 INT-002 — Reliable Messaging}
\phantomsection
\label{navigation:I.5.39}
\addcontentsline{toc}{section}{I.5.39 INT-002 — Reliable Messaging}

\label{element:aac9b66911e849cb9792c3a2b7eacb88}
\index{INT-002 — Reliable Messaging (I.5.39)}
{\small Requirement · Historical — current owning-domain standing not inferred · historical\_reference}\par
\href{../../../../../../../continuum/motivation/requirement/aac9b66911e849cb9792c3a2b7eacb88.md}{Historical canonical definition and provenance}.\par
\section*{I.5.40 INT-003 — Messaging Recovery}
\phantomsection
\label{navigation:I.5.40}
\addcontentsline{toc}{section}{I.5.40 INT-003 — Messaging Recovery}

\label{element:6e72cf09e7cc4acf9466dd39e1a15ae6}
\index{INT-003 — Messaging Recovery (I.5.40)}
{\small Requirement · Historical — current owning-domain standing not inferred · historical\_reference}\par
\href{../../../../../../../continuum/motivation/requirement/6e72cf09e7cc4acf9466dd39e1a15ae6.md}{Historical canonical definition and provenance}.\par
\section*{I.5.41 INT-004 — FHIR R5 Boundary}
\phantomsection
\label{navigation:I.5.41}
\addcontentsline{toc}{section}{I.5.41 INT-004 — FHIR R5 Boundary}

\label{element:927b0cf4710e424391120b124067051e}
\index{INT-004 — FHIR R5 Boundary (I.5.41)}
{\small Requirement · Historical — current owning-domain standing not inferred · historical\_reference}\par
\href{../../../../../../../continuum/motivation/requirement/927b0cf4710e424391120b124067051e.md}{Historical canonical definition and provenance}.\par
\section*{I.5.42 INT-005 — Opaque Resource Identity}
\phantomsection
\label{navigation:I.5.42}
\addcontentsline{toc}{section}{I.5.42 INT-005 — Opaque Resource Identity}

\label{element:91bd7a8732b3468b84be2493103433e7}
\index{INT-005 — Opaque Resource Identity (I.5.42)}
{\small Requirement · Historical — current owning-domain standing not inferred · historical\_reference}\par
\href{../../../../../../../continuum/motivation/requirement/91bd7a8732b3468b84be2493103433e7.md}{Historical canonical definition and provenance}.\par
\section*{I.5.43 INT-006 — Governed Validation}
\phantomsection
\label{navigation:I.5.43}
\addcontentsline{toc}{section}{I.5.43 INT-006 — Governed Validation}

\label{element:6c14251b5d4e4eff96d2fd812834d124}
\index{INT-006 — Governed Validation (I.5.43)}
{\small Requirement · Historical — current owning-domain standing not inferred · historical\_reference}\par
\href{../../../../../../../continuum/motivation/requirement/6c14251b5d4e4eff96d2fd812834d124.md}{Historical canonical definition and provenance}.\par
\section*{I.5.44 INT-007 — Terminology Abstraction}
\phantomsection
\label{navigation:I.5.44}
\addcontentsline{toc}{section}{I.5.44 INT-007 — Terminology Abstraction}

\label{element:6fe86be142044e5aa7720e2e6493a72a}
\index{INT-007 — Terminology Abstraction (I.5.44)}
{\small Requirement · Historical — current owning-domain standing not inferred · historical\_reference}\par
\href{../../../../../../../continuum/motivation/requirement/6fe86be142044e5aa7720e2e6493a72a.md}{Historical canonical definition and provenance}.\par
\section*{I.5.45 Normal Reporting Does Not Require Raw PubSub Topology Interpretation}
\phantomsection
\label{navigation:I.5.45}
\addcontentsline{toc}{section}{I.5.45 Normal Reporting Does Not Require Raw PubSub Topology Interpretation}

\label{element:aa569bb1448648f685bb73cf9c6563e3}
\index{Normal Reporting Does Not Require Raw PubSub Topology Interpretation (I.5.45)}
{\small Requirement · Adjudicated Domain05 direction — U04/U06/U14/U18; holistic result acceptance pending · current}\par

Harmonia SHALL NOT require humans to reason directly from raw Pub/Sub topology and telemetry for normal reporting and assurance purposes.



This assertion is normalised from the governed DesignNote under U04/U06/U14/U18 and the ARB release Praxis. Original candidate/release qualifications and separately deferred implementation details remain in force.



Architectural subject: \hyperref[element:fb3f37a6ac614415ac486a6de9d35994]{Energeia}.


{\footnotesize \href{../../../../../../../continuum/motivation/requirement/aa569bb1448648f685bb73cf9c6563e3.md}{Canonical record: aa569bb1448648f685bb73cf9c6563e3}; full qualifications and provenance.}\par
\section*{I.5.46 Operations and Assurance Supports Semantic Preserving Partitioning}
\phantomsection
\label{navigation:I.5.46}
\addcontentsline{toc}{section}{I.5.46 Operations and Assurance Supports Semantic Preserving Partitioning}

\label{element:3238d37e27a7421e905b705c5f30fc85}
\index{Operations and Assurance Supports Semantic Preserving Partitioning (I.5.46)}
{\small Requirement · Adjudicated Domain05 direction — U11; holistic result acceptance pending · current}\par

\textbf{Harmonia Operations \& Assurance (O\&A)} SHALL permit horizontal partitioning without changing semantic identity, provenance, relationships or governed meaning.



This assertion is normalised from the governed DesignNote under U11 and the ARB release Praxis. Original candidate/release qualifications and separately deferred implementation details remain in force.



Architectural subject: \hyperref[element:7985e647bf9e4ca3ba70faf9e3e4ae02]{Mnemosyne}.


{\footnotesize \href{../../../../../../../continuum/motivation/requirement/3238d37e27a7421e905b705c5f30fc85.md}{Canonical record: 3238d37e27a7421e905b705c5f30fc85}; full qualifications and provenance.}\par
\section*{I.5.47 OPS-001 — Reproducible Deployment}
\phantomsection
\label{navigation:I.5.47}
\addcontentsline{toc}{section}{I.5.47 OPS-001 — Reproducible Deployment}

\label{element:c4373b1fd8444f10b502a5c3216152d7}
\index{OPS-001 — Reproducible Deployment (I.5.47)}
{\small Requirement · Historical — current owning-domain standing not inferred · historical\_reference}\par
\href{../../../../../../../continuum/motivation/requirement/c4373b1fd8444f10b502a5c3216152d7.md}{Historical canonical definition and provenance}.\par
\section*{I.5.48 OPS-002 — Durable Infrastructure State}
\phantomsection
\label{navigation:I.5.48}
\addcontentsline{toc}{section}{I.5.48 OPS-002 — Durable Infrastructure State}

\label{element:d5d6a9ea450047b1bd62b13177e8fca6}
\index{OPS-002 — Durable Infrastructure State (I.5.48)}
{\small Requirement · Historical — current owning-domain standing not inferred · historical\_reference}\par
\href{../../../../../../../continuum/motivation/requirement/d5d6a9ea450047b1bd62b13177e8fca6.md}{Historical canonical definition and provenance}.\par
\section*{I.5.49 [AssuranceRequirement] OPS-003 — Meaningful Health and Readiness}
\phantomsection
\label{navigation:I.5.49}
\addcontentsline{toc}{section}{I.5.49 [AssuranceRequirement] OPS-003 — Meaningful Health and Readiness}

\label{element:4c02fec3af264b878ee3e0cfaecbf55f}
\index{[AssuranceRequirement] OPS-003 — Meaningful Health and Readiness (I.5.49)}
{\small Requirement / AssuranceRequirement · Solution Architect-reviewed Iteration 001; authorised ADD; resulting increment review pending · current}\par

\par\medskip\textbf{Accepted Iteration 001 contribution}\par


\textbf{NFR-059 — Component and Dependency Health Distinction}



Distinguish component health from dependency health



Primary subject: \texttt{\seqsplit{G09.06}}. Native Requirement; Harmonia specialisation AssuranceRequirement.



\par\medskip\textbf{Preserved historical contribution}\par


Original semantic standing: Historical — current owning-domain standing not inferred. Original currentness: \texttt{\seqsplit{historical\_reference}}. The following original content remains historical; the new contribution does not independently approve its additional details.



Services shall expose health and readiness information that reflects material dependencies rather than merely indicating that a process or network port exists.


{\footnotesize \href{../../../../../../../continuum/motivation/requirement/assurance-requirement/4c02fec3af264b878ee3e0cfaecbf55f.md}{Canonical record: 4c02fec3af264b878ee3e0cfaecbf55f}; full qualifications and provenance.}\par
\section*{I.5.50 [AssuranceRequirement] OPS-004 — Operational Observability}
\phantomsection
\label{navigation:I.5.50}
\addcontentsline{toc}{section}{I.5.50 [AssuranceRequirement] OPS-004 — Operational Observability}

\label{element:b5d3447a05b249a8af154136fc36c4d9}
\index{[AssuranceRequirement] OPS-004 — Operational Observability (I.5.50)}
{\small Requirement / AssuranceRequirement · Solution Architect-reviewed Iteration 001; authorised ADD; resulting increment review pending · current}\par

\par\medskip\textbf{Accepted Iteration 001 contribution}\par


\textbf{NFR-057 — Operational Failure Diagnostics}



Diagnose operational failures without direct access to implementation internals



Primary subject: \texttt{\seqsplit{G09.02}}. Native Requirement; Harmonia specialisation AssuranceRequirement.



\par\medskip\textbf{Preserved historical contribution}\par


Original semantic standing: Historical — current owning-domain standing not inferred. Original currentness: \texttt{\seqsplit{historical\_reference}}. The following original content remains historical; the new contribution does not independently approve its additional details.



Harmonia shall expose sufficient operational telemetry to diagnose subsystem, interface, queue, workflow, persistence, event, and task health without requiring direct inspection of implementation internals.


{\footnotesize \href{../../../../../../../continuum/motivation/requirement/assurance-requirement/b5d3447a05b249a8af154136fc36c4d9.md}{Canonical record: b5d3447a05b249a8af154136fc36c4d9}; full qualifications and provenance.}\par
\section*{I.5.51 OPS-005 — Production Resilience}
\phantomsection
\label{navigation:I.5.51}
\addcontentsline{toc}{section}{I.5.51 OPS-005 — Production Resilience}

\label{element:01e7d6168d554048a424ef75a6556040}
\index{OPS-005 — Production Resilience (I.5.51)}
{\small Requirement · Historical — current owning-domain standing not inferred · historical\_reference}\par
\href{../../../../../../../continuum/motivation/requirement/01e7d6168d554048a424ef75a6556040.md}{Historical canonical definition and provenance}.\par
\section*{I.5.52 OPS-006 — Architecture Conformance}
\phantomsection
\label{navigation:I.5.52}
\addcontentsline{toc}{section}{I.5.52 OPS-006 — Architecture Conformance}

\label{element:dee1792d83bf413f974fc24802084872}
\index{OPS-006 — Architecture Conformance (I.5.52)}
{\small Requirement · Historical — current owning-domain standing not inferred · historical\_reference}\par
\href{../../../../../../../continuum/motivation/requirement/dee1792d83bf413f974fc24802084872.md}{Historical canonical definition and provenance}.\par
\section*{I.5.53 Preserve Necessary Evidence When Historical Resolution Is Unavailable}
\phantomsection
\label{navigation:I.5.53}
\addcontentsline{toc}{section}{I.5.53 Preserve Necessary Evidence When Historical Resolution Is Unavailable}

\label{element:81ce6a25725e41bb8e92c945f53c8b4c}
\index{Preserve Necessary Evidence When Historical Resolution Is Unavailable (I.5.53)}
{\small Requirement · Adjudicated Domain05 direction — U07/U15; holistic result acceptance pending · current}\par

Where the authoritative source cannot guarantee longitudinal retrieval of the information as it existed at the relevant audited point in time, Kleio SHALL preserve the information necessary to maintain the integrity of the audit evidence.



This assertion is normalised from the governed DesignNote under U07/U15 and the ARB release Praxis. Original candidate/release qualifications and separately deferred implementation details remain in force.



Architectural subject: \hyperref[element:fefda110e2274884ad58906733606efd]{Kleio}.


{\footnotesize \href{../../../../../../../continuum/motivation/requirement/81ce6a25725e41bb8e92c945f53c8b4c.md}{Canonical record: 81ce6a25725e41bb8e92c945f53c8b4c}; full qualifications and provenance.}\par
\section*{I.5.54 [AssuranceRequirement] PROV-001 — End-to-End Provenance}
\phantomsection
\label{navigation:I.5.54}
\addcontentsline{toc}{section}{I.5.54 [AssuranceRequirement] PROV-001 — End-to-End Provenance}

\label{element:3e45e60affb44bf384d54de26d3add6a}
\index{[AssuranceRequirement] PROV-001 — End-to-End Provenance (I.5.54)}
{\small Requirement / AssuranceRequirement · Solution Architect-reviewed Iteration 001; authorised ADD; resulting increment review pending · current}\par

\par\medskip\textbf{Accepted Iteration 001 contribution}\par


\textbf{NFR-019 — Information Provenance}



Preserve information provenance where required



Primary subject: \texttt{\seqsplit{G02.02}}. Native Requirement; Harmonia specialisation AssuranceRequirement.



\par\medskip\textbf{Preserved historical contribution}\par


Original semantic standing: Historical — current owning-domain standing not inferred. Original currentness: \texttt{\seqsplit{historical\_reference}}. The following original content remains historical; the new contribution does not independently approve its additional details.



Harmonia shall retain sufficient provenance to determine the origin, actors, transformations, processing history, and resulting artefacts associated with clinically or operationally material information.


{\footnotesize \href{../../../../../../../continuum/motivation/requirement/assurance-requirement/3e45e60affb44bf384d54de26d3add6a.md}{Canonical record: 3e45e60affb44bf384d54de26d3add6a}; full qualifications and provenance.}\par
\section*{I.5.55 PROV-002 — Source Context}
\phantomsection
\label{navigation:I.5.55}
\addcontentsline{toc}{section}{I.5.55 PROV-002 — Source Context}

\label{element:929ed42779a843dd84b683e3dcc35eb3}
\index{PROV-002 — Source Context (I.5.55)}
{\small Requirement · Historical — current owning-domain standing not inferred · historical\_reference}\par
\href{../../../../../../../continuum/motivation/requirement/929ed42779a843dd84b683e3dcc35eb3.md}{Historical canonical definition and provenance}.\par
\section*{I.5.56 PROV-003 — Dual-Actor Provenance}
\phantomsection
\label{navigation:I.5.56}
\addcontentsline{toc}{section}{I.5.56 PROV-003 — Dual-Actor Provenance}

\label{element:c016286faa3d4dfbbad133f73e855f1c}
\index{PROV-003 — Dual-Actor Provenance (I.5.56)}
{\small Requirement · Historical — current owning-domain standing not inferred · historical\_reference}\par
\href{../../../../../../../continuum/motivation/requirement/c016286faa3d4dfbbad133f73e855f1c.md}{Historical canonical definition and provenance}.\par
\section*{I.5.57 PROV-004 — Provenance and Authority Independence}
\phantomsection
\label{navigation:I.5.57}
\addcontentsline{toc}{section}{I.5.57 PROV-004 — Provenance and Authority Independence}

\label{element:3fb9a89d345a4cb6a56d62823316664e}
\index{PROV-004 — Provenance and Authority Independence (I.5.57)}
{\small Requirement · Historical — current owning-domain standing not inferred · historical\_reference}\par
\href{../../../../../../../continuum/motivation/requirement/3fb9a89d345a4cb6a56d62823316664e.md}{Historical canonical definition and provenance}.\par
\section*{I.5.58 PROV-005 — Immutable Audit}
\phantomsection
\label{navigation:I.5.58}
\addcontentsline{toc}{section}{I.5.58 PROV-005 — Immutable Audit}

\label{element:0f7b23cfd94e450c8a2772879bc9bb1e}
\index{PROV-005 — Immutable Audit (I.5.58)}
{\small Requirement · Historical — current owning-domain standing not inferred · historical\_reference}\par
\href{../../../../../../../continuum/motivation/requirement/0f7b23cfd94e450c8a2772879bc9bb1e.md}{Historical canonical definition and provenance}.\par
\section*{I.5.59 PROV-006 — Audit Discoverability}
\phantomsection
\label{navigation:I.5.59}
\addcontentsline{toc}{section}{I.5.59 PROV-006 — Audit Discoverability}

\label{element:927c65d2d665451396782b4c1d7c8c86}
\index{PROV-006 — Audit Discoverability (I.5.59)}
{\small Requirement · Historical — current owning-domain standing not inferred · historical\_reference}\par
\href{../../../../../../../continuum/motivation/requirement/927c65d2d665451396782b4c1d7c8c86.md}{Historical canonical definition and provenance}.\par
\section*{I.5.60 REQ-FND-001: Durable Ingress Acceptance Boundary}
\phantomsection
\label{navigation:I.5.60}
\addcontentsline{toc}{section}{I.5.60 REQ-FND-001: Durable Ingress Acceptance Boundary}

\label{element:01c834537a1e41dfb523d48f287d93e9}
\index{REQ-FND-001: Durable Ingress Acceptance Boundary (I.5.60)}
{\small Requirement · Accepted · current}\par

\par\medskip\textbf{Normative Requirement Statement}\par


\begin{quote}


\textbf{Harmonia shall not positively acknowledge acceptance of an incoming event until responsibility for that event has crossed the required durable acceptance boundary.}



\end{quote}


\par\medskip\textbf{Architectural Intent \& Rationale}\par


This requirement eliminates the critical vulnerability identified in \href{\detokenize{../../../../../../../../docs/markdown/01-motivation/drivers-assessments/assessments.md}}{Silent Data Loss via False Acceptance}. When an upstream external system transmits a clinical event (such as an emergency admission, critical lab result, or medication order), it remains the responsibility of that upstream sender until Harmonia has durably secured the event.



Releasing an upstream sender from responsibility (by returning a positive transport acknowledgement) before the event is committed to durable storage creates a vulnerable window where gateway crashes, memory exhaustions, or hardware reboots cause silent, unrecoverable data loss.



\par\medskip\textbf{Protocol-Neutral Delineation \& Downstream Realisations}\par


\begin{itemize}


\item \textbf{Normative Rule}: The requirement mandates that the durable acceptance boundary must be crossed before positive acknowledgement. It is completely protocol-neutral.

\item \textbf{Downstream Realisations}:

\item In HL7 v2 over MLLP integration (owned by Domain 06 Integration Architecture), this rule is realised by forbidding an \texttt{\seqsplit{MSA-1 = AA}} (Application Accept) acknowledgement until downstream persistence or queue publication succeeds. If persistence fails, the gateway MUST return \texttt{\seqsplit{MSA-1 = AE}} (Application Error) to trigger upstream sender retry.

\item In FHIR RESTful APIs, this rule is realised by returning \texttt{\seqsplit{201 Created}} or \texttt{\seqsplit{200 OK}} only after transactional database commitment, returning \texttt{\seqsplit{202 Accepted}} only when durable asynchronous queueing is verified, and returning \texttt{\seqsplit{500 Internal Server Error}} or \texttt{\seqsplit{503 Service Unavailable}} on durable failure.

\item 

In Ingress Dual-Write safety (\texttt{\seqsplit{REC-001}}), this rule enforces that in-memory active caching alone is insufficient to satisfy positive acceptance.




\item 

\textbf{Primary Motivation}: \href{\detokenize{../../../../../../../../docs/markdown/01-motivation/drivers-assessments/assessments.md}}{Silent Data Loss via False Acceptance}, \href{\detokenize{../../../../../../../../docs/markdown/01-motivation/drivers-assessments/drivers.md}}{Continuous Clinical Service Availability}




\item \textbf{Governing Principles}: \hyperref[element:0412a2084f0145babc4f022c71fea65f]{AX-05 (State Separation)}, \hyperref[element:0d9673833eb840eab03804aa076c1b0c]{AX-10 (Failure as Normal Condition)}, \hyperref[element:ef8abbaec37e4890ad78b9f263e67906]{AX-15 (Explicit Representation of Uncertainty)}

\item \textbf{Historical Identifier Reference}: \texttt{\seqsplit{REQ-INGRESS-001}}

\end{itemize}

{\footnotesize \href{../../../../../../../continuum/motivation/requirement/01c834537a1e41dfb523d48f287d93e9.md}{Canonical record: 01c834537a1e41dfb523d48f287d93e9}; full qualifications and provenance.}\par
\section*{I.5.61 REQ-FND-002: Operational Activity Progression State}
\phantomsection
\label{navigation:I.5.61}
\addcontentsline{toc}{section}{I.5.61 REQ-FND-002: Operational Activity Progression State}

\label{element:eaabcdb96aad4f32a050c94f9ecec049}
\index{REQ-FND-002: Operational Activity Progression State (I.5.61)}
{\small Requirement · Accepted · current}\par

\par\medskip\textbf{Normative Requirement Statement}\par


\begin{quote}


\textbf{Harmonia shall maintain sufficient explicit and observable state to determine the progression and outcome of managed multi-stage operational activities.}



\end{quote}


\par\medskip\textbf{Architectural Intent \& Rationale}\par


Complex health integration involves distributed, multi-stage activities—such as routing ADT events, dispatching laboratory orders, and executing multi-destination fan-out deliveries. If integration middleware treats these operations as disconnected fire-and-forget messages, failures at individual destination endpoints cause unmonitored divergence between hospital systems (as diagnosed in \href{\detokenize{../../../../../../../../docs/markdown/01-motivation/drivers-assessments/assessments.md}}{Unmonitored Destination Failure}).



Harmonia mandates that every multi-stage operational activity must maintain explicit, queryable, and observable progression state across its entire lifecycle. Operational state must capture checkpoints, destination delivery statuses, error codes, and completion outcomes, enabling operational engineers to monitor, audit, and reconcile distributed activities.



\par\medskip\textbf{Downstream Scope}\par


\begin{itemize}


\item Realised in Domain 05 Application Architecture and Domain 06 Integration Architecture through Pragma envelopes, FHIR \texttt{\seqsplit{Task}} resources, and task progression tracking.

\item 

Does not mandate a specific workflow engine; it mandates observable progression state synchronized with real-world entity state.




\item 

\textbf{Primary Motivation}: \href{\detokenize{../../../../../../../../docs/markdown/01-motivation/drivers-assessments/assessments.md}}{Unmonitored Destination Failure / Fan-Out Divergence}, \href{\detokenize{../../../../../../../../docs/markdown/01-motivation/drivers-assessments/drivers.md}}{Coordinated Operational Activity and Entity State}




\item \textbf{Governing Principle}: \hyperref[element:de7e90f46127455a9ec9ad3b9b48ba54]{AX-16 (Operational Activity \& Entity State Progress Together)}

\item \textbf{Historical Identifier Reference}: \texttt{\seqsplit{REQ-COORD-001}}

\end{itemize}

{\footnotesize \href{../../../../../../../continuum/motivation/requirement/eaabcdb96aad4f32a050c94f9ecec049.md}{Canonical record: eaabcdb96aad4f32a050c94f9ecec049}; full qualifications and provenance.}\par
\section*{I.5.62 REQ-FND-003: Subject Referential Integrity}
\phantomsection
\label{navigation:I.5.62}
\addcontentsline{toc}{section}{I.5.62 REQ-FND-003: Subject Referential Integrity}

\label{element:43ecbfe8f9b14fcbbc59ee0c03c8774b}
\index{REQ-FND-003: Subject Referential Integrity (I.5.62)}
{\small Requirement · Accepted · current}\par

\par\medskip\textbf{Normative Requirement Statement}\par


\begin{quote}


\textbf{Harmonia shall preserve and validate the integrity of subject references and associations used within managed information, preventing known-invalid or internally inconsistent associations from being silently treated as valid.}



\end{quote}


\par\medskip\textbf{Architectural Intent \& Rationale}\par


Clinical records derive their entire medical validity from their association with the correct human patient subject. In regional health networks, patient identifiers arrive from disparate sources (national IHIs, regional MRNs, local facility patient numbers). Mismatched or corrupted subject references lead directly to patient harm.



Harmonia mandates that any subject reference attached to a managed clinical event must be structurally validated, referentially verified, and checked for internal consistency. Known-invalid identifiers (e.g., failed check-digits, malformed namespaces, or conflicting patient identifiers within a single clinical bundle) must be rejected or explicitly flagged rather than silently propagated as valid.



\par\medskip\textbf{Explicit Boundary: Non-MPI Demarcation}\par


\begin{itemize}


\item \textbf{Not a Master Patient Index (MPI)}: This requirement does \textbf{not} establish Harmonia as an MPI or enterprise demographic matching engine.

\item \textbf{Not a Golden-Record Determiner}: Harmonia does not execute probabilistic matching or merge patient demographics into a synthetic "golden record."

\item 

\textbf{Intent}: Preserving the referential integrity, authenticity, and provenance of subject identifiers without performing lossy or speculative demographic unification.




\item 

\textbf{Primary Motivation}: \href{\detokenize{../../../../../../../../docs/markdown/01-motivation/drivers-assessments/drivers.md}}{Patient Safety \& Identity Integrity}




\item \textbf{External Constraint}: \href{\detokenize{../../../../../../../../docs/markdown/01-motivation/requirements-constraints/external-constraints.md}}{Applicable National / Jurisdictional Healthcare Identifier Obligations}

\item \textbf{Governing Principles}: \hyperref[element:952c3b6148ce4fce94b69f99a0c2fec7]{AX-06 (Explicit Information Authority)}, \hyperref[element:d3fc6025dc7646d1b2020b51f7a44bfa]{AX-14 (Preservation of Semantic Distinctions)}

\item \textbf{Historical Identifier Reference}: \texttt{\seqsplit{REQ-IDENT-001}}

\end{itemize}

{\footnotesize \href{../../../../../../../continuum/motivation/requirement/43ecbfe8f9b14fcbbc59ee0c03c8774b.md}{Canonical record: 43ecbfe8f9b14fcbbc59ee0c03c8774b}; full qualifications and provenance.}\par
\section*{I.5.63 REQ-FND-004: Explicit Indeterminate Outcome}
\phantomsection
\label{navigation:I.5.63}
\addcontentsline{toc}{section}{I.5.63 REQ-FND-004: Explicit Indeterminate Outcome}

\label{element:fc732b64fc564127b73174e28d9ccce8}
\index{REQ-FND-004: Explicit Indeterminate Outcome (I.5.63)}
{\small Requirement · Accepted · current}\par

\par\medskip\textbf{Normative Requirement Statement}\par


\begin{quote}


\textbf{Where Harmonia cannot establish the outcome of a managed operation with sufficient certainty, the outcome shall be represented explicitly as indeterminate and shall not be interpreted as either success or failure without subsequent reconciliation or another authoritative determination.}



\end{quote}


\par\medskip\textbf{Architectural Intent \& Rationale}\par


In distributed networks, communication timeouts, socket disconnects, and unacknowledged dispatches leave the status of remote operations uncertain. In the presence of a network partition or destination crash, a client cannot know whether an update was processed and the ACK was lost, or whether the update was never received.



A frequent architectural failure mode is "binary collapsing": assuming an unconfirmed operation failed (triggering duplicate side effects on blind retry) or assuming it succeeded (ignoring unperformed clinical orders). Harmonia mandates that uncertainty must be preserved and explicitly represented as \texttt{\seqsplit{INDETERMINATE}}.



\par\medskip\textbf{Reconciliation \& Retry Semantics}\par


\begin{itemize}


\item This requirement does \textbf{not} prohibit automated retries or resilient recovery.

\item 

It mandates that retry mechanisms must be safe, idempotent, and reconciliation-aware. An indeterminate outcome must be resolved through explicit reconciliation (such as querying destination state or verifying idempotency tokens) rather than blind, speculative assumption.




\item 

\textbf{Primary Motivation}: \href{\detokenize{../../../../../../../../docs/markdown/01-motivation/drivers-assessments/assessments.md}}{Silent Data Loss via False Acceptance}, \href{\detokenize{../../../../../../../../docs/markdown/01-motivation/drivers-assessments/drivers.md}}{Continuous Clinical Service Availability}




\item \textbf{Governing Principle}: \hyperref[element:ef8abbaec37e4890ad78b9f263e67906]{AX-15 (Uncertainty Is Preserved Until Resolved)}

\item \textbf{Historical Identifier Reference}: \texttt{\seqsplit{REQ-UNCERT-001}}

\end{itemize}

{\footnotesize \href{../../../../../../../continuum/motivation/requirement/fc732b64fc564127b73174e28d9ccce8.md}{Canonical record: fc732b64fc564127b73174e28d9ccce8}; full qualifications and provenance.}\par
\section*{I.5.64 REQ-FND-005: Independent Assurance of Governed Activity}
\phantomsection
\label{navigation:I.5.64}
\addcontentsline{toc}{section}{I.5.64 REQ-FND-005: Independent Assurance of Governed Activity}

\label{element:bbb7438e78604308abba93afc14f4e2e}
\index{REQ-FND-005: Independent Assurance of Governed Activity (I.5.64)}
{\small Requirement · APPROVED · current}\par

\textbf{Classification:} Foundational Platform Requirement. \textbf{Status:} APPROVED — human architectural approval recorded on 2026-10-08; normative wording and relationships unchanged. The authorised Independent Assurance reconciliation is APPROVED / CLOSED; Domain01 Motivation Architecture is CLOSED / FROZEN.



\par\medskip\textbf{Normative Requirement Statement}\par


\begin{quote}


\textbf{Harmonia shall support independent governed assurance of activities, information and outcomes for which assurance is required, such that assurance conclusions are established from sufficient and trustworthy evidence against applicable obligations, expectations or controls.}



\textbf{The activity or mechanism responsible for performing or managing the subject being assured shall not be solely responsible for determining, suppressing, manufacturing or retrospectively altering its assurance outcome. Harmonia shall establish sufficient separation of assurance responsibility and authority from responsibility for performance or control of the subject so that assurance is not merely self-attestation by the subject being assured.}



\textbf{Assurance may identify findings, exceptions and the need for operational response, but shall not thereby assume responsibility for management, assignment, delegation, escalation or remediation of the activity being assured.}



\textbf{Assurance may manage its own assurance activity without thereby assuming management of the subject activity.}



\textbf{Where available evidence is insufficient to establish an assurance conclusion with the required degree of confidence, that insufficiency shall remain explicit and shall not be interpreted as either assurance or non-assurance.}



\end{quote}


\par\medskip\textbf{Architectural Intent \& Rationale}\par


Those accountable for health-information handling and operational activity need credible grounds for demonstrating that applicable obligations, expectations or controls have been met. The \hyperref[element:e76216c494fa4500bb80eda6f83513c1]{Regional Health Network Operator's compliance and forensic-review concerns}, \hyperref[element:c20aa17ce4aa4a58be5eb3cff2ecdd1d]{engineers' incident review and deterministic verification concerns}, and \hyperref[element:75519ed8f7de4a619cc9065f525d703a]{O3's demonstrable protection and accountable handling} already establish the value of substantive review and trustworthy evidence.



Those existing concerns do not themselves establish performer independence. The approved reconciliation adds that obligation explicitly: a conclusion controlled solely by the activity or mechanism being assured can conceal findings or manufacture confidence and cannot provide the required independent assurance. Evidence availability and operational completion alone cannot establish whether the applicable obligations, expectations or controls have been satisfied.



\par\medskip\textbf{Applicability \& Responsibility Boundary}\par


The requirement applies where assurance is required. It does not require assurance of every activity, information item or outcome. The criteria establishing applicability, applicable obligations and the required degree of confidence remain for subsequent governed derivation; this requirement does not define an applicability model.



Independence concerns responsibility and authority. It does not prescribe physical, technical or organisational separation, nor require external assurance. It does not allocate an assessor or prohibit the subject's performer from contributing evidence. Assurance manages its own work and may inform an operational response; responsibility for managing that response remains distinct. The explanatory shorthand is: \textbf{the assurance capability is the QA service, not the Job Foreman.}



\par\medskip\textbf{Semantic \& Uncertainty Boundaries}\par


The following distinctions constrain this requirement; they do not define new information types or an assurance outcome taxonomy:



\begin{quote}\small\ttfamily
Provenance ≠ Assurance\\
Audit information ≠ Assurance\\
Evidence ≠ Assurance conclusion\\
Operational progression ≠ Assurance\\
Workflow completion ≠ Assurance\\
Business Outcome ≠ Assurance Outcome\\
Operational uncertainty ≠ Insufficient assurance evidence\\
Independent assurance ≠ External assurance\\
Assurance finding ≠ Operational response\\
Assurance responsibility ≠ Operational management responsibility\\

\end{quote}


\hyperref[element:fc732b64fc564127b73174e28d9ccce8]{\texttt{\seqsplit{REQ-FND-004}}} and \hyperref[element:ef8abbaec37e4890ad78b9f263e67906]{AX-15} continue to govern uncertainty about an operation's state, effect or outcome, including an assurance activity's own execution where applicable. Separately, an assurance activity may execute correctly yet lack sufficient evidence to establish a conclusion about its subject. This requirement preserves that insufficiency without treating it as an uncertain execution outcome, execution failure, assurance or non-assurance. Neither existing operational rule is rewritten or extended into an assurance-evidence sufficiency rule.



\par\medskip\textbf{Motivation Relationships}\par


These are branching relationships, not a sequential motivation pipeline. The independence, management boundary and evidence-insufficiency obligations are newly approved intent for reconciliation, not meanings retrospectively attributed to the reused elements.



\begin{longtable}{>{\raggedright\arraybackslash}p{\dimexpr\linewidth/3-2\tabcolsep\relax}>{\raggedright\arraybackslash}p{\dimexpr\linewidth/3-2\tabcolsep\relax}>{\raggedright\arraybackslash}p{\dimexpr\linewidth/3-2\tabcolsep\relax}}
Existing element & Relationship to \texttt{\seqsplit{REQ-FND-005}} & Semantic limit \\
\hyperref[element:e76216c494fa4500bb80eda6f83513c1]{Regional Health Network Operator} & Compliance demonstration and forensic-review concerns motivate the requirement directly. & No new Operator-to-Privacy-Driver relationship or statutory independence mandate is inferred. \\
\hyperref[element:c20aa17ce4aa4a58be5eb3cff2ecdd1d]{Platform Operations \& Integration Engineers} & Incident review and deterministic verification concerns motivate the requirement directly. & Operational diagnosis and monitoring remain distinct from independent assurance. \\
\hyperref[element:1dbf34b452b6497f914c704e037aec24]{Statutory Health Information Privacy \& Protection Driver} & Supports the requirement's need to evaluate handling against applicable obligations and controls. & Privacy supplies one valid motivation, not the entire assurance scope or independence obligation. \\
\hyperref[element:75519ed8f7de4a619cc9065f525d703a]{O3} & The requirement supports O3's existing demonstrability and accountability purpose. & Textual contribution only; no new causal handover, intermediate Goal or assurance-outcome equivalence. \\
\hyperref[element:0ed6fb9fcc9547659ce76b917881274e]{CST-EXT-001} & Bounds assurance concerning health-information privacy and protection by applicable obligations. & No universal external requirement for independent assurance is claimed. \\
\hyperref[element:952c3b6148ce4fce94b69f99a0c2fec7]{AX-06} & Co-governs interpretation of evidence authority and credibility independently of technical state. & Information authority does not establish assessor independence. \\
\hyperref[element:5de2ae2c3b794c2d9f3a14efd51d564f]{AX-07} & Co-governs security context and policy for assurance activity. & Security enforcement does not establish independent conclusion authority. \\
\hyperref[element:21d6864657964ba5b03e3f6b5eb29cd5]{AX-08} & Co-governs meaningful provenance and audit evidence used in assurance. & Preserving evidence does not perform assurance or establish a conclusion. \\
\hyperref[element:d3fc6025dc7646d1b2020b51f7a44bfa]{AX-14} & Co-governs preservation of the requirement's semantic distinctions. & Distinct meanings do not prescribe separate technologies or information types. \\
\end{longtable}


\hyperref[element:9a5cb9e8876844e58e660f909f7df9fc]{AX-09} remains an evidence-selection and retention boundary: observations and progression state are not automatically provenance or audit evidence. \hyperref[element:eaabcdb96aad4f32a050c94f9ecec049]{\texttt{\seqsplit{REQ-FND-002}}} can supply observable operational facts; it is a complementary evidence relationship, not an independence parent or assurance conclusion. \texttt{\seqsplit{REQ-FND-004}} / AX-15 have the complementary uncertainty relationship described above. \hyperref[element:8882573df7ce43c6bc0f66da0a676139]{AX-17} governs this explicit reconciliation and preservation of unresolved downstream derivation; it is not a source of the real-world assurance need.



\par\medskip\textbf{Downstream Derivation Boundary}\par


The requirement supplies upstream motivation from which Strategy responsibility, Business assurance behaviour and Information Architecture may subsequently be derived, followed by Application Architecture and an eventual solution. Those relationships remain unestablished here. No capability allocation, workflow specification, solution construct, outcome values, evidence model, execution mechanism or deployment boundary is defined by this requirement.


{\footnotesize \href{../../../../../../../continuum/motivation/requirement/bbb7438e78604308abba93afc14f4e2e.md}{Canonical record: bbb7438e78604308abba93afc14f4e2e}; full qualifications and provenance.}\par
\section*{I.5.65 SEC-001 — Trusted Authentication}
\phantomsection
\label{navigation:I.5.65}
\addcontentsline{toc}{section}{I.5.65 SEC-001 — Trusted Authentication}

\label{element:e4998b0a132c41afbe610b11869c3297}
\index{SEC-001 — Trusted Authentication (I.5.65)}
{\small Requirement · Historical — current owning-domain standing not inferred · historical\_reference}\par
\href{../../../../../../../continuum/motivation/requirement/e4998b0a132c41afbe610b11869c3297.md}{Historical canonical definition and provenance}.\par
\section*{I.5.66 SEC-002 — Default-Deny Authorisation}
\phantomsection
\label{navigation:I.5.66}
\addcontentsline{toc}{section}{I.5.66 SEC-002 — Default-Deny Authorisation}

\label{element:3bdecb66e3f34110aa0871ba34091bfd}
\index{SEC-002 — Default-Deny Authorisation (I.5.66)}
{\small Requirement · Historical — current owning-domain standing not inferred · historical\_reference}\par
\href{../../../../../../../continuum/motivation/requirement/3bdecb66e3f34110aa0871ba34091bfd.md}{Historical canonical definition and provenance}.\par
\section*{I.5.67 SEC-003 — Authentication and Authorisation Separation}
\phantomsection
\label{navigation:I.5.67}
\addcontentsline{toc}{section}{I.5.67 SEC-003 — Authentication and Authorisation Separation}

\label{element:88f2ace9a07d4ed7b5a54a4af93af3c9}
\index{SEC-003 — Authentication and Authorisation Separation (I.5.67)}
{\small Requirement · Historical — current owning-domain standing not inferred · historical\_reference}\par
\href{../../../../../../../continuum/motivation/requirement/88f2ace9a07d4ed7b5a54a4af93af3c9.md}{Historical canonical definition and provenance}.\par
\section*{I.5.68 SEC-004 — Contextual Resource Authorisation}
\phantomsection
\label{navigation:I.5.68}
\addcontentsline{toc}{section}{I.5.68 SEC-004 — Contextual Resource Authorisation}

\label{element:1f6822aef4bc44548c6708fcc018e4ba}
\index{SEC-004 — Contextual Resource Authorisation (I.5.68)}
{\small Requirement · Historical — current owning-domain standing not inferred · historical\_reference}\par
\href{../../../../../../../continuum/motivation/requirement/1f6822aef4bc44548c6708fcc018e4ba.md}{Historical canonical definition and provenance}.\par
\section*{I.5.69 SEC-005 — End-to-End Security Context}
\phantomsection
\label{navigation:I.5.69}
\addcontentsline{toc}{section}{I.5.69 SEC-005 — End-to-End Security Context}

\label{element:d70e0a51440640b1bd17562b6ea4ff98}
\index{SEC-005 — End-to-End Security Context (I.5.69)}
{\small Requirement · Historical — current owning-domain standing not inferred · historical\_reference}\par
\href{../../../../../../../continuum/motivation/requirement/d70e0a51440640b1bd17562b6ea4ff98.md}{Historical canonical definition and provenance}.\par
\section*{I.5.70 SEC-006 — Initiating and Executing Identity}
\phantomsection
\label{navigation:I.5.70}
\addcontentsline{toc}{section}{I.5.70 SEC-006 — Initiating and Executing Identity}

\label{element:5eaac47ebef94bbb818e4ba751601ab1}
\index{SEC-006 — Initiating and Executing Identity (I.5.70)}
{\small Requirement · Historical — current owning-domain standing not inferred · historical\_reference}\par
\href{../../../../../../../continuum/motivation/requirement/5eaac47ebef94bbb818e4ba751601ab1.md}{Historical canonical definition and provenance}.\par
\section*{I.5.71 SEC-007 — Anti-Spoofing}
\phantomsection
\label{navigation:I.5.71}
\addcontentsline{toc}{section}{I.5.71 SEC-007 — Anti-Spoofing}

\label{element:8996bdd1790e4efeb78bedb85269c916}
\index{SEC-007 — Anti-Spoofing (I.5.71)}
{\small Requirement · Historical — current owning-domain standing not inferred · historical\_reference}\par
\href{../../../../../../../continuum/motivation/requirement/8996bdd1790e4efeb78bedb85269c916.md}{Historical canonical definition and provenance}.\par
\section*{I.5.72 SEC-008 — Security Metadata Preservation}
\phantomsection
\label{navigation:I.5.72}
\addcontentsline{toc}{section}{I.5.72 SEC-008 — Security Metadata Preservation}

\label{element:c037846f015645728b4afe4d011ae0cd}
\index{SEC-008 — Security Metadata Preservation (I.5.72)}
{\small Requirement · Historical — current owning-domain standing not inferred · historical\_reference}\par
\href{../../../../../../../continuum/motivation/requirement/c037846f015645728b4afe4d011ae0cd.md}{Historical canonical definition and provenance}.\par
\section*{I.5.73 SEC-009 — Least Privilege}
\phantomsection
\label{navigation:I.5.73}
\addcontentsline{toc}{section}{I.5.73 SEC-009 — Least Privilege}

\label{element:f760e0cc08354fc88c8da049958b45ed}
\index{SEC-009 — Least Privilege (I.5.73)}
{\small Requirement · Historical — current owning-domain standing not inferred · historical\_reference}\par
\href{../../../../../../../continuum/motivation/requirement/f760e0cc08354fc88c8da049958b45ed.md}{Historical canonical definition and provenance}.\par
\section*{I.5.74 SEC-010 — PHI-Safe Operational Logging}
\phantomsection
\label{navigation:I.5.74}
\addcontentsline{toc}{section}{I.5.74 SEC-010 — PHI-Safe Operational Logging}

\label{element:4a078bd44e874a48abc650f87d64f16f}
\index{SEC-010 — PHI-Safe Operational Logging (I.5.74)}
{\small Requirement · Historical — current owning-domain standing not inferred · historical\_reference}\par
\href{../../../../../../../continuum/motivation/requirement/4a078bd44e874a48abc650f87d64f16f.md}{Historical canonical definition and provenance}.\par
\section*{I.5.75 UX-001 — Consistent Iris Experience}
\phantomsection
\label{navigation:I.5.75}
\addcontentsline{toc}{section}{I.5.75 UX-001 — Consistent Iris Experience}

\label{element:6bea2694012442a39034a6cbfcc79e90}
\index{UX-001 — Consistent Iris Experience (I.5.75)}
{\small Requirement · Historical — current owning-domain standing not inferred · historical\_reference}\par
\href{../../../../../../../continuum/motivation/requirement/6bea2694012442a39034a6cbfcc79e90.md}{Historical canonical definition and provenance}.\par
\section*{I.5.76 UX-002 — Information-Dense Operational Design}
\phantomsection
\label{navigation:I.5.76}
\addcontentsline{toc}{section}{I.5.76 UX-002 — Information-Dense Operational Design}

\label{element:73354d29efa74895afebec8437c875e8}
\index{UX-002 — Information-Dense Operational Design (I.5.76)}
{\small Requirement · Historical — current owning-domain standing not inferred · historical\_reference}\par
\href{../../../../../../../continuum/motivation/requirement/73354d29efa74895afebec8437c875e8.md}{Historical canonical definition and provenance}.\par
\section*{I.5.77 UX-003 — Honest System State}
\phantomsection
\label{navigation:I.5.77}
\addcontentsline{toc}{section}{I.5.77 UX-003 — Honest System State}

\label{element:752452c7367d4d4dbc1cb9ffe0bfd78f}
\index{UX-003 — Honest System State (I.5.77)}
{\small Requirement · Historical — current owning-domain standing not inferred · historical\_reference}\par
\href{../../../../../../../continuum/motivation/requirement/752452c7367d4d4dbc1cb9ffe0bfd78f.md}{Historical canonical definition and provenance}.\par
\section*{I.5.78 UX-004 — Explicit Availability States}
\phantomsection
\label{navigation:I.5.78}
\addcontentsline{toc}{section}{I.5.78 UX-004 — Explicit Availability States}

\label{element:0e4b5219260e419d987195c9832a0bbe}
\index{UX-004 — Explicit Availability States (I.5.78)}
{\small Requirement · Historical — current owning-domain standing not inferred · historical\_reference}\par
\href{../../../../../../../continuum/motivation/requirement/0e4b5219260e419d987195c9832a0bbe.md}{Historical canonical definition and provenance}.\par
\section*{I.5.79 UX-005 — Patient Search and Workspace}
\phantomsection
\label{navigation:I.5.79}
\addcontentsline{toc}{section}{I.5.79 UX-005 — Patient Search and Workspace}

\label{element:32c55c722ae147329784976793abd458}
\index{UX-005 — Patient Search and Workspace (I.5.79)}
{\small Requirement · Historical — current owning-domain standing not inferred · historical\_reference}\par
\href{../../../../../../../continuum/motivation/requirement/32c55c722ae147329784976793abd458.md}{Historical canonical definition and provenance}.\par
\section*{I.5.80 UX-006 — Longitudinal Clinical Navigation}
\phantomsection
\label{navigation:I.5.80}
\addcontentsline{toc}{section}{I.5.80 UX-006 — Longitudinal Clinical Navigation}

\label{element:09318a29f8b747c8ac6b0e276e2febb4}
\index{UX-006 — Longitudinal Clinical Navigation (I.5.80)}
{\small Requirement · Historical — current owning-domain standing not inferred · historical\_reference}\par
\href{../../../../../../../continuum/motivation/requirement/09318a29f8b747c8ac6b0e276e2febb4.md}{Historical canonical definition and provenance}.\par
\section*{I.5.81 UX-007 — Persistent Patient Context}
\phantomsection
\label{navigation:I.5.81}
\addcontentsline{toc}{section}{I.5.81 UX-007 — Persistent Patient Context}

\label{element:873b8b2c9a514ab387833beb5552feda}
\index{UX-007 — Persistent Patient Context (I.5.81)}
{\small Requirement · Historical — current owning-domain standing not inferred · historical\_reference}\par
\href{../../../../../../../continuum/motivation/requirement/873b8b2c9a514ab387833beb5552feda.md}{Historical canonical definition and provenance}.\par
\section*{I.5.82 UX-008 — Provenance-Aware Clinical Presentation}
\phantomsection
\label{navigation:I.5.82}
\addcontentsline{toc}{section}{I.5.82 UX-008 — Provenance-Aware Clinical Presentation}

\label{element:bb298bafe5894d9b9c057b896f669c06}
\index{UX-008 — Provenance-Aware Clinical Presentation (I.5.82)}
{\small Requirement · Historical — current owning-domain standing not inferred · historical\_reference}\par
\href{../../../../../../../continuum/motivation/requirement/bb298bafe5894d9b9c057b896f669c06.md}{Historical canonical definition and provenance}.\par
\section*{I.5.83 UX-009 — Initial Clinical Scope}
\phantomsection
\label{navigation:I.5.83}
\addcontentsline{toc}{section}{I.5.83 UX-009 — Initial Clinical Scope}

\label{element:807861850fa54d6c849e3e26d7172b9c}
\index{UX-009 — Initial Clinical Scope (I.5.83)}
{\small Requirement · Historical — current owning-domain standing not inferred · historical\_reference}\par
\href{../../../../../../../continuum/motivation/requirement/807861850fa54d6c849e3e26d7172b9c.md}{Historical canonical definition and provenance}.\par
\section*{I.5.84 [AssuranceRequirement] Execution Outcome Attribution}
\phantomsection
\label{navigation:I.5.84}
\addcontentsline{toc}{section}{I.5.84 [AssuranceRequirement] Execution Outcome Attribution}

\label{element:79177908b67243a4a5f86fa180f5fd22}
\index{[AssuranceRequirement] Execution Outcome Attribution (I.5.84)}
{\small Requirement / AssuranceRequirement · Solution Architect-reviewed Iteration 001; authorised ADD; resulting increment review pending · current}\par

Attribute outcomes to originating execution contexts



Source entry: \textbf{NFR-026 — Execution Outcome Attribution}. Primary subject: \texttt{\seqsplit{G06.06}}.



The source classification is Requirement «AssuranceRequirement». The source identifier is a scoped register reference and does not determine architectural type. Applicability, release scope and unresolved thresholds remain as stated above.


{\footnotesize \href{../../../../../../../continuum/motivation/requirement/assurance-requirement/79177908b67243a4a5f86fa180f5fd22.md}{Canonical record: 79177908b67243a4a5f86fa180f5fd22}; full qualifications and provenance.}\par
\section*{I.5.85 [AssuranceRequirement] Auditable Security Decisions}
\phantomsection
\label{navigation:I.5.85}
\addcontentsline{toc}{section}{I.5.85 [AssuranceRequirement] Auditable Security Decisions}

\label{element:92ef574eb9b945dfaf6cf2b1b670595b}
\index{[AssuranceRequirement] Auditable Security Decisions (I.5.85)}
{\small Requirement / AssuranceRequirement · Solution Architect-reviewed Iteration 001; authorised ADD; resulting increment review pending · current}\par

Make security decisions independently auditable



Source entry: \textbf{NFR-039 — Auditable Security Decisions}. Primary subject: \texttt{\seqsplit{G03.06}}.



The source classification is Requirement «AssuranceRequirement». The source identifier is a scoped register reference and does not determine architectural type. Applicability, release scope and unresolved thresholds remain as stated above.


{\footnotesize \href{../../../../../../../continuum/motivation/requirement/assurance-requirement/92ef574eb9b945dfaf6cf2b1b670595b.md}{Canonical record: 92ef574eb9b945dfaf6cf2b1b670595b}; full qualifications and provenance.}\par
\section*{I.5.86 [AssuranceRequirement] Audit Context Identification}
\phantomsection
\label{navigation:I.5.86}
\addcontentsline{toc}{section}{I.5.86 [AssuranceRequirement] Audit Context Identification}

\label{element:de466c998be34fae87a2be41983cbcdd}
\index{[AssuranceRequirement] Audit Context Identification (I.5.86)}
{\small Requirement / AssuranceRequirement · Solution Architect-reviewed Iteration 001; authorised ADD; resulting increment review pending · current}\par

Identify relevant operation, subject and execution context in audit records



Source entry: \textbf{NFR-044 — Audit Context Identification}. Primary subject: \texttt{\seqsplit{G04.02}}.



The source classification is Requirement «AssuranceRequirement». The source identifier is a scoped register reference and does not determine architectural type. Applicability, release scope and unresolved thresholds remain as stated above.


{\footnotesize \href{../../../../../../../continuum/motivation/requirement/assurance-requirement/de466c998be34fae87a2be41983cbcdd.md}{Canonical record: de466c998be34fae87a2be41983cbcdd}; full qualifications and provenance.}\par
\section*{I.5.87 [AssuranceRequirement] Audit Retention}
\phantomsection
\label{navigation:I.5.87}
\addcontentsline{toc}{section}{I.5.87 [AssuranceRequirement] Audit Retention}

\label{element:44b1a897fd1a41c097a1138bfee9cea6}
\index{[AssuranceRequirement] Audit Retention (I.5.87)}
{\small Requirement / AssuranceRequirement · Solution Architect-reviewed Iteration 001; authorised ADD; resulting increment review pending · current}\par

Retain audit information for its defined retention period



Source entry: \textbf{NFR-046 — Audit Retention}. Primary subject: \texttt{\seqsplit{G04.04}}.



The source classification is Requirement «AssuranceRequirement». The source identifier is a scoped register reference and does not determine architectural type. Applicability, release scope and unresolved thresholds remain as stated above.


{\footnotesize \href{../../../../../../../continuum/motivation/requirement/assurance-requirement/44b1a897fd1a41c097a1138bfee9cea6.md}{Canonical record: 44b1a897fd1a41c097a1138bfee9cea6}; full qualifications and provenance.}\par
\section*{I.5.88 [AssuranceRequirement] Cross-Component Execution Traceability}
\phantomsection
\label{navigation:I.5.88}
\addcontentsline{toc}{section}{I.5.88 [AssuranceRequirement] Cross-Component Execution Traceability}

\label{element:8014fbaa48b5435785fced72aa748da6}
\index{[AssuranceRequirement] Cross-Component Execution Traceability (I.5.88)}
{\small Requirement / AssuranceRequirement · Solution Architect-reviewed Iteration 001; authorised ADD; resulting increment review pending · current}\par

Trace execution across participating components



Source entry: \textbf{NFR-047 — Cross-Component Execution Traceability}. Primary subject: \texttt{\seqsplit{G04.04}}.



The source classification is Requirement «AssuranceRequirement». The source identifier is a scoped register reference and does not determine architectural type. Applicability, release scope and unresolved thresholds remain as stated above.


{\footnotesize \href{../../../../../../../continuum/motivation/requirement/assurance-requirement/8014fbaa48b5435785fced72aa748da6.md}{Canonical record: 8014fbaa48b5435785fced72aa748da6}; full qualifications and provenance.}\par
\section*{I.5.89 [AssuranceRequirement] Audit Noninterference}
\phantomsection
\label{navigation:I.5.89}
\addcontentsline{toc}{section}{I.5.89 [AssuranceRequirement] Audit Noninterference}

\label{element:ccea3a4dbd494cbe99c525ff2e87db26}
\index{[AssuranceRequirement] Audit Noninterference (I.5.89)}
{\small Requirement / AssuranceRequirement · Solution Architect-reviewed Iteration 001; authorised ADD; resulting increment review pending · current}\par

Ensure audit recording does not compromise audited operation integrity



Source entry: \textbf{NFR-048 — Audit Noninterference}. Primary subject: \texttt{\seqsplit{G04.01}}.



The source classification is Requirement «AssuranceRequirement». The source identifier is a scoped register reference and does not determine architectural type. Applicability, release scope and unresolved thresholds remain as stated above.


{\footnotesize \href{../../../../../../../continuum/motivation/requirement/assurance-requirement/ccea3a4dbd494cbe99c525ff2e87db26.md}{Canonical record: ccea3a4dbd494cbe99c525ff2e87db26}; full qualifications and provenance.}\par
\section*{I.5.90 [AssuranceRequirement] Exchange Identity and Provenance}
\phantomsection
\label{navigation:I.5.90}
\addcontentsline{toc}{section}{I.5.90 [AssuranceRequirement] Exchange Identity and Provenance}

\label{element:5579b7ce9df540359b7ba36f807b52c5}
\index{[AssuranceRequirement] Exchange Identity and Provenance (I.5.90)}
{\small Requirement / AssuranceRequirement · Solution Architect-reviewed Iteration 001; authorised ADD; resulting increment review pending · current}\par

Preserve required identity and provenance in exchange



Source entry: \textbf{NFR-053 — Exchange Identity and Provenance}. Primary subject: \texttt{\seqsplit{G05.04}}.



The source classification is Requirement «AssuranceRequirement». The source identifier is a scoped register reference and does not determine architectural type. Applicability, release scope and unresolved thresholds remain as stated above.


{\footnotesize \href{../../../../../../../continuum/motivation/requirement/assurance-requirement/5579b7ce9df540359b7ba36f807b52c5.md}{Canonical record: 5579b7ce9df540359b7ba36f807b52c5}; full qualifications and provenance.}\par
\section*{I.5.91 [NonFunctionalRequirement] Event Ingestion Capacity}
\phantomsection
\label{navigation:I.5.91}
\addcontentsline{toc}{section}{I.5.91 [NonFunctionalRequirement] Event Ingestion Capacity}

\label{element:75420022c1a7463e91ca48f229a78194}
\index{[NonFunctionalRequirement] Event Ingestion Capacity (I.5.91)}
{\small Requirement / NonFunctionalRequirement · Solution Architect-reviewed Iteration 001; authorised ADD; resulting increment review pending · current}\par

R1.x accepts 100,000 incoming messages/PUTs/POSTs over 8 hours (12,500/hour; approximately 3.47/second average). This workload is processed together with the NFR-066, NFR-068 and NFR-069 obligations; the 8-hour incoming workload is fully processed by hour 24 while 100 practitioners and 10 OA\&M users remain active.



Source entry: \textbf{NFR-063 — Event Ingestion Capacity}. Primary subject: \texttt{\seqsplit{G08.01}}.



The source classification is Requirement «NonFunctionalRequirement». The source identifier is a scoped register reference and does not determine architectural type. Applicability, release scope and unresolved thresholds remain as stated above.


{\footnotesize \href{../../../../../../../continuum/motivation/requirement/non-functional-requirement/75420022c1a7463e91ca48f229a78194.md}{Canonical record: 75420022c1a7463e91ca48f229a78194}; full qualifications and provenance.}\par
\section*{I.5.92 [NonFunctionalRequirement] Governed Search Performance}
\phantomsection
\label{navigation:I.5.92}
\addcontentsline{toc}{section}{I.5.92 [NonFunctionalRequirement] Governed Search Performance}

\label{element:b29075c844cf403fb1322b4d1aba8803}
\index{[NonFunctionalRequirement] Governed Search Performance (I.5.92)}
{\small Requirement / NonFunctionalRequirement · Solution Architect-reviewed Iteration 001; authorised ADD; resulting increment review pending · current}\par

Harmonia's governed search capabilities shall satisfy defined performance requirements independently of create, read, update and delete operations



Source entry: \textbf{NFR-065 — Governed Search Performance}. Primary subject: \texttt{\seqsplit{G08.02}}.



The source classification is Requirement «NonFunctionalRequirement». The source identifier is a scoped register reference and does not determine architectural type. Applicability, release scope and unresolved thresholds remain as stated above.


{\small Recorded unresolved matters: searchPerformanceTargets; exact scope remains in the canonical record.}\par
{\footnotesize \href{../../../../../../../continuum/motivation/requirement/non-functional-requirement/b29075c844cf403fb1322b4d1aba8803.md}{Canonical record: b29075c844cf403fb1322b4d1aba8803}; full qualifications and provenance.}\par
\section*{I.5.93 [NonFunctionalRequirement] Event Processing Completion}
\phantomsection
\label{navigation:I.5.93}
\addcontentsline{toc}{section}{I.5.93 [NonFunctionalRequirement] Event Processing Completion}

\label{element:d1380eec700a492b810354e0798c8146}
\index{[NonFunctionalRequirement] Event Processing Completion (I.5.93)}
{\small Requirement / NonFunctionalRequirement · Solution Architect-reviewed Iteration 001; authorised ADD; resulting increment review pending · current}\par

Complete processing of the 100,000 events received during an 8-hour window within 24 hours of the start of that window (16 additional hours).



Source entry: \textbf{NFR-066 — Event Processing Completion}. Primary subject: \texttt{\seqsplit{G08.01}}.



The source classification is Requirement «NonFunctionalRequirement». The source identifier is a scoped register reference and does not determine architectural type. Applicability, release scope and unresolved thresholds remain as stated above.


{\footnotesize \href{../../../../../../../continuum/motivation/requirement/non-functional-requirement/d1380eec700a492b810354e0798c8146.md}{Canonical record: d1380eec700a492b810354e0798c8146}; full qualifications and provenance.}\par
\section*{I.5.94 [NonFunctionalRequirement] Concurrent Practitioner Workload}
\phantomsection
\label{navigation:I.5.94}
\addcontentsline{toc}{section}{I.5.94 [NonFunctionalRequirement] Concurrent Practitioner Workload}

\label{element:e77f2e780a974dde99815ff63158941c}
\index{[NonFunctionalRequirement] Concurrent Practitioner Workload (I.5.94)}
{\small Requirement / NonFunctionalRequirement · Solution Architect-reviewed Iteration 001; authorised ADD; resulting increment review pending · current}\par

Support 100 concurrent practitioners across all non-OA\&M user interfaces while ingesting and processing the specified 100,000-event workload.



Source entry: \textbf{NFR-068 — Concurrent Practitioner Workload}. Primary subject: \texttt{\seqsplit{G08.04}}.



The source classification is Requirement «NonFunctionalRequirement». The source identifier is a scoped register reference and does not determine architectural type. Applicability, release scope and unresolved thresholds remain as stated above.


{\footnotesize \href{../../../../../../../continuum/motivation/requirement/non-functional-requirement/e77f2e780a974dde99815ff63158941c.md}{Canonical record: e77f2e780a974dde99815ff63158941c}; full qualifications and provenance.}\par
\section*{I.5.95 [NonFunctionalRequirement] Concurrent OA\&M Workload}
\phantomsection
\label{navigation:I.5.95}
\addcontentsline{toc}{section}{I.5.95 [NonFunctionalRequirement] Concurrent OA\&M Workload}

\label{element:5220295ae5d144468066ae70683e0f19}
\index{[NonFunctionalRequirement] Concurrent OA\&M Workload (I.5.95)}
{\small Requirement / NonFunctionalRequirement · Solution Architect-reviewed Iteration 001; authorised ADD; resulting increment review pending · current}\par

Support 10 concurrent OA\&M users simultaneously with the practitioner and event-processing workloads.



Source entry: \textbf{NFR-069 — Concurrent OA\&M Workload}. Primary subject: \texttt{\seqsplit{G08.04}}.



The source classification is Requirement «NonFunctionalRequirement». The source identifier is a scoped register reference and does not determine architectural type. Applicability, release scope and unresolved thresholds remain as stated above.


{\footnotesize \href{../../../../../../../continuum/motivation/requirement/non-functional-requirement/5220295ae5d144468066ae70683e0f19.md}{Canonical record: 5220295ae5d144468066ae70683e0f19}; full qualifications and provenance.}\par
\section*{I.5.96 [FunctionalRequirement] R1.x Accessibility Scope}
\phantomsection
\label{navigation:I.5.96}
\addcontentsline{toc}{section}{I.5.96 [FunctionalRequirement] R1.x Accessibility Scope}

\label{element:8369e22bc3414ae49a1987095549ef03}
\index{[FunctionalRequirement] R1.x Accessibility Scope (I.5.96)}
{\small Requirement / FunctionalRequirement · Solution Architect-reviewed Iteration 001; authorised ADD; resulting increment review pending · current}\par

R1.x user interfaces are not required to conform to WCAG accessibility standards.



Source entry: \textbf{NFR-070 — R1.x Accessibility Scope}. Primary subject: \texttt{\seqsplit{G12.02}}.



The source classification is Requirement «FunctionalRequirement». The source identifier is a scoped register reference and does not determine architectural type. Applicability, release scope and unresolved thresholds remain as stated above.


{\footnotesize \href{../../../../../../../continuum/motivation/requirement/functional-requirement/8369e22bc3414ae49a1987095549ef03.md}{Canonical record: 8369e22bc3414ae49a1987095549ef03}; full qualifications and provenance.}\par
\section*{I.5.97 [FunctionalRequirement] R1.x User Authentication}
\phantomsection
\label{navigation:I.5.97}
\addcontentsline{toc}{section}{I.5.97 [FunctionalRequirement] R1.x User Authentication}

\label{element:9683182663cd40428d806ed2d8685511}
\index{[FunctionalRequirement] R1.x User Authentication (I.5.97)}
{\small Requirement / FunctionalRequirement · Solution Architect-reviewed Iteration 001; authorised ADD; resulting increment review pending · current}\par

R1.x: unauthenticated users shall not access Harmonia functionality; system-to-system identities remain a separate consideration.



Source entry: \textbf{NFR-071 — R1.x User Authentication}. Primary subject: \texttt{\seqsplit{G03.01}}.



The source classification is Requirement «FunctionalRequirement». The source identifier is a scoped register reference and does not determine architectural type. Applicability, release scope and unresolved thresholds remain as stated above.


{\footnotesize \href{../../../../../../../continuum/motivation/requirement/functional-requirement/9683182663cd40428d806ed2d8685511.md}{Canonical record: 9683182663cd40428d806ed2d8685511}; full qualifications and provenance.}\par
\section*{I.5.98 [NonFunctionalRequirement] User Interaction Responsiveness}
\phantomsection
\label{navigation:I.5.98}
\addcontentsline{toc}{section}{I.5.98 [NonFunctionalRequirement] User Interaction Responsiveness}

\label{element:0f07f89a3a854bfdac4d4fc423feefca}
\index{[NonFunctionalRequirement] User Interaction Responsiveness (I.5.98)}
{\small Requirement / NonFunctionalRequirement · Solution Architect-reviewed Iteration 001; authorised ADD; resulting increment review pending · current}\par

Provide prompt interaction responses and appropriate feedback for operations that cannot complete immediately; response-time thresholds TBD.



Source entry: \textbf{NFR-081 — User Interaction Responsiveness}. Primary subject: \texttt{\seqsplit{G12.03}}.



The source classification is Requirement «NonFunctionalRequirement». The source identifier is a scoped register reference and does not determine architectural type. Applicability, release scope and unresolved thresholds remain as stated above.


{\small Recorded unresolved matters: responseTimeThresholds; exact scope remains in the canonical record.}\par
{\footnotesize \href{../../../../../../../continuum/motivation/requirement/non-functional-requirement/0f07f89a3a854bfdac4d4fc423feefca.md}{Canonical record: 0f07f89a3a854bfdac4d4fc423feefca}; full qualifications and provenance.}\par
\section*{I.5.99 [AssuranceRequirement] Single-Machine Failure Recovery}
\phantomsection
\label{navigation:I.5.99}
\addcontentsline{toc}{section}{I.5.99 [AssuranceRequirement] Single-Machine Failure Recovery}

\label{element:0d13da9390a3438ca887a4293bb034b9}
\index{[AssuranceRequirement] Single-Machine Failure Recovery (I.5.99)}
{\small Requirement / AssuranceRequirement · Solution Architect-reviewed Iteration 001; authorised ADD; resulting increment review pending · current}\par

Harmonia shall demonstrate recovery from the failure of any single physical or virtual machine without loss of accepted governed work or the need for manual intervention.



Source entry: \textbf{NFR-094 — Single-Machine Failure Recovery}. Primary subject: \texttt{\seqsplit{G07.03}}.



The source classification is Requirement «AssuranceRequirement». The source identifier is a scoped register reference and does not determine architectural type. Applicability, release scope and unresolved thresholds remain as stated above.


{\footnotesize \href{../../../../../../../continuum/motivation/requirement/assurance-requirement/0d13da9390a3438ca887a4293bb034b9.md}{Canonical record: 0d13da9390a3438ca887a4293bb034b9}; full qualifications and provenance.}\par
\section*{I.5.100 [NonFunctionalRequirement] Operational Availability}
\phantomsection
\label{navigation:I.5.100}
\addcontentsline{toc}{section}{I.5.100 [NonFunctionalRequirement] Operational Availability}

\label{element:65ac0367b7f241149ad1839032004039}
\index{[NonFunctionalRequirement] Operational Availability (I.5.100)}
{\small Requirement / NonFunctionalRequirement · Solution Architect-reviewed Iteration 001; authorised ADD; resulting increment review pending · current}\par

Meet a defined operational availability target



Source entry: \textbf{NFR-105 — Operational Availability}. Primary subject: \texttt{\seqsplit{G07.05}}.



The source classification is Requirement «NonFunctionalRequirement». The source identifier is a scoped register reference and does not determine architectural type. Applicability, release scope and unresolved thresholds remain as stated above.


{\small Recorded unresolved matters: operationalAvailabilityTarget; exact scope remains in the canonical record.}\par
{\footnotesize \href{../../../../../../../continuum/motivation/requirement/non-functional-requirement/65ac0367b7f241149ad1839032004039.md}{Canonical record: 65ac0367b7f241149ad1839032004039}; full qualifications and provenance.}\par
\section*{I.5.101 [FunctionalRequirement] Duplicate Message Handling}
\phantomsection
\label{navigation:I.5.101}
\addcontentsline{toc}{section}{I.5.101 [FunctionalRequirement] Duplicate Message Handling}

\label{element:76a67526fcdf40999f375a8cac54d5ef}
\index{[FunctionalRequirement] Duplicate Message Handling (I.5.101)}
{\small Requirement / FunctionalRequirement · Solution Architect-reviewed Iteration 001; authorised ADD; resulting increment review pending · current}\par

Harmonia's message-processing capabilities shall support the identification and governed handling of duplicate incoming messages, consistent with the applicable message-processing policies.



Source entry: \textbf{NFR-136 — Duplicate Message Handling}. Primary subject: \texttt{\seqsplit{G06.04}}.



The source classification is Requirement «FunctionalRequirement». The source identifier is a scoped register reference and does not determine architectural type. Applicability, release scope and unresolved thresholds remain as stated above.


{\footnotesize \href{../../../../../../../continuum/motivation/requirement/functional-requirement/76a67526fcdf40999f375a8cac54d5ef.md}{Canonical record: 76a67526fcdf40999f375a8cac54d5ef}; full qualifications and provenance.}\par
\section*{I.5.102 [FunctionalRequirement] Message Processing Ordering}
\phantomsection
\label{navigation:I.5.102}
\addcontentsline{toc}{section}{I.5.102 [FunctionalRequirement] Message Processing Ordering}

\label{element:3c11adccd7194276b787416368c13858}
\index{[FunctionalRequirement] Message Processing Ordering (I.5.102)}
{\small Requirement / FunctionalRequirement · Solution Architect-reviewed Iteration 001; authorised ADD; resulting increment review pending · current}\par

Harmonia's message-processing capabilities shall support the preservation and enforcement of message-processing order where required by the applicable message-processing policies and interaction contracts.



Source entry: \textbf{NFR-137 — Message Processing Ordering}. Primary subject: \texttt{\seqsplit{G13.06}}.



The source classification is Requirement «FunctionalRequirement». The source identifier is a scoped register reference and does not determine architectural type. Applicability, release scope and unresolved thresholds remain as stated above.


{\footnotesize \href{../../../../../../../continuum/motivation/requirement/functional-requirement/3c11adccd7194276b787416368c13858.md}{Canonical record: 3c11adccd7194276b787416368c13858}; full qualifications and provenance.}\par
\section*{I.5.103 [FunctionalRequirement] Manual Processing Intervention}
\phantomsection
\label{navigation:I.5.103}
\addcontentsline{toc}{section}{I.5.103 [FunctionalRequirement] Manual Processing Intervention}

\label{element:911511665a8a4135b7f2a37f5f79ad36}
\index{[FunctionalRequirement] Manual Processing Intervention (I.5.103)}
{\small Requirement / FunctionalRequirement · Solution Architect-reviewed Iteration 001; authorised ADD; resulting increment review pending · current}\par

Harmonia shall support authorised manual intervention in stalled, failed or otherwise exceptional processing activities, consistent with the applicable workflow governance and operational control policies.



Source entry: \textbf{NFR-143 — Manual Processing Intervention}. Primary subject: \texttt{\seqsplit{G06.02}}.



The source classification is Requirement «FunctionalRequirement». The source identifier is a scoped register reference and does not determine architectural type. Applicability, release scope and unresolved thresholds remain as stated above.


{\footnotesize \href{../../../../../../../continuum/motivation/requirement/functional-requirement/911511665a8a4135b7f2a37f5f79ad36.md}{Canonical record: 911511665a8a4135b7f2a37f5f79ad36}; full qualifications and provenance.}\par
\section*{I.5.104 [AssuranceRequirement] Processing Outcome Evidence}
\phantomsection
\label{navigation:I.5.104}
\addcontentsline{toc}{section}{I.5.104 [AssuranceRequirement] Processing Outcome Evidence}

\label{element:8ac5880576ad4f8fa9612f231178a06c}
\index{[AssuranceRequirement] Processing Outcome Evidence (I.5.104)}
{\small Requirement / AssuranceRequirement · Solution Architect-reviewed Iteration 001; authorised ADD; resulting increment review pending · current}\par

Harmonia shall provide sufficient diagnostic and execution evidence to enable verification of processing outcomes against their defined expectations and applicable processing obligations.



Source entry: \textbf{NFR-151 — Processing Outcome Evidence}. Primary subject: \texttt{\seqsplit{G06.06}}.



The source classification is Requirement «AssuranceRequirement». The source identifier is a scoped register reference and does not determine architectural type. Applicability, release scope and unresolved thresholds remain as stated above.


{\footnotesize \href{../../../../../../../continuum/motivation/requirement/assurance-requirement/8ac5880576ad4f8fa9612f231178a06c.md}{Canonical record: 8ac5880576ad4f8fa9612f231178a06c}; full qualifications and provenance.}\par
\chapter*{I.6 Constraints}
\phantomsection
\label{navigation:I.6}
\addcontentsline{toc}{chapter}{I.6 Constraints}

\section*{I.6.1 AuthorisationRole Criteria Are Additive}
\phantomsection
\label{navigation:I.6.1}
\addcontentsline{toc}{section}{I.6.1 AuthorisationRole Criteria Are Additive}

\label{element:d5cfaf6f8c224d1b9769b1edd7c27591}
\index{AuthorisationRole Criteria Are Additive (I.6.1)}
{\small Constraint · Adjudicated Domain05 direction — U12/U15/U16; holistic result acceptance pending · current}\par

AuthorisationRole containment SHALL NOT be interpreted as requiring all contained Criteria to be simultaneously satisfied.



This assertion is normalised from the governed DesignNote under U12/U15/U16 and the ARB release Praxis. Original candidate/release qualifications and separately deferred implementation details remain in force.



Architectural subject: \hyperref[element:899d34ec491948b68aad09e507dea6d1]{Themis}.


{\footnotesize \href{../../../../../../../continuum/motivation/constraint/d5cfaf6f8c224d1b9769b1edd7c27591.md}{Canonical record: d5cfaf6f8c224d1b9769b1edd7c27591}; full qualifications and provenance.}\par
\section*{I.6.2 Authoritative AuditEvents Are Immutable}
\phantomsection
\label{navigation:I.6.2}
\addcontentsline{toc}{section}{I.6.2 Authoritative AuditEvents Are Immutable}

\label{element:21eb20cdc600411cb71fbfe2abce4dbe}
\index{Authoritative AuditEvents Are Immutable (I.6.2)}
{\small Constraint · Adjudicated Domain05 direction — U07/U15; holistic result acceptance pending · current}\par

\begin{quote}


\textbf{An existing authoritative AuditEvent SHALL NOT be corrected, amended, replaced or overwritten.}



\end{quote}


This assertion is normalised from the governed DesignNote under U07/U15 and the ARB release Praxis. Original candidate/release qualifications and separately deferred implementation details remain in force.



Architectural subject: \hyperref[element:fefda110e2274884ad58906733606efd]{Kleio}.


{\footnotesize \href{../../../../../../../continuum/motivation/constraint/21eb20cdc600411cb71fbfe2abce4dbe.md}{Canonical record: 21eb20cdc600411cb71fbfe2abce4dbe}; full qualifications and provenance.}\par
\section*{I.6.3 Authority Forwarding Requires FORWARD Authorisation}
\phantomsection
\label{navigation:I.6.3}
\addcontentsline{toc}{section}{I.6.3 Authority Forwarding Requires FORWARD Authorisation}

\label{element:5aed12f6540445ed85c1866adbef93af}
\index{Authority Forwarding Requires FORWARD Authorisation (I.6.3)}
{\small Constraint · Adjudicated Domain05 direction — U12/U15/U16; holistic result acceptance pending · current}\par

Forwarding SHALL itself be subject to applicable FORWARD authorisation.



This assertion is normalised from the governed DesignNote under U12/U15/U16 and the ARB release Praxis. Original candidate/release qualifications and separately deferred implementation details remain in force.



Architectural subject: \hyperref[element:899d34ec491948b68aad09e507dea6d1]{Themis}.


{\footnotesize \href{../../../../../../../continuum/motivation/constraint/5aed12f6540445ed85c1866adbef93af.md}{Canonical record: 5aed12f6540445ed85c1866adbef93af}; full qualifications and provenance.}\par
\section*{I.6.4 Configuration and Governance Has Coherent Authority}
\phantomsection
\label{navigation:I.6.4}
\addcontentsline{toc}{section}{I.6.4 Configuration and Governance Has Coherent Authority}

\label{element:7cdbc752ca2e4ebf9ee5acff84fd7400}
\index{Configuration and Governance Has Coherent Authority (I.6.4)}
{\small Constraint · Adjudicated Domain05 direction — U11; holistic result acceptance pending · current}\par

\textbf{Harmonia Configuration \& Governance (C\&G)} SHALL remain a single coherent governed information space and SHALL NOT be partitioned in a manner that creates independent or divergent governance authority.



This assertion is normalised from the governed DesignNote under U11 and the ARB release Praxis. Original candidate/release qualifications and separately deferred implementation details remain in force.



Architectural subject: \hyperref[element:7985e647bf9e4ca3ba70faf9e3e4ae02]{Mnemosyne}.


{\footnotesize \href{../../../../../../../continuum/motivation/constraint/7cdbc752ca2e4ebf9ee5acff84fd7400.md}{Canonical record: 7cdbc752ca2e4ebf9ee5acff84fd7400}; full qualifications and provenance.}\par
\section*{I.6.5 CST-EXT-001: Applicable Health-Information Privacy and Data-Protection Obligations}
\phantomsection
\label{navigation:I.6.5}
\addcontentsline{toc}{section}{I.6.5 CST-EXT-001: Applicable Health-Information Privacy and Data-Protection Obligations}

\label{element:0ed6fb9fcc9547659ce76b917881274e}
\index{CST-EXT-001: Applicable Health-Information Privacy and Data-Protection Obligations (I.6.5)}
{\small Constraint · Accepted · current}\par

\par\medskip\textbf{Categorical Definition}\par


\begin{quote}


\textbf{Harmonia must operate within the health-information privacy, confidentiality, disclosure, retention, and data-protection obligations applicable to the jurisdiction and deployment context.}



\end{quote}


\par\medskip\textbf{Context \& Architectural Boundary}\par


Health data is categorized as hypersensitive personal data across all legal jurisdictions. Wherever Harmonia is deployed, it is bound by the statutory data protection regime governing that region. The platform must be architecturally designed to satisfy these obligations without requiring bespoke rewrites of core domain logic.



\par\medskip\textbf{Representative Instances \& Examples}\par


\begin{itemize}


\item \textbf{Australia}: The \emph{Privacy Act 1988 (Cth)}, the Australian Privacy Principles (APPs), the \emph{My Health Records Act 2012}, and state-based health records legislation (e.g., \emph{Health Records Act 2001 (Vic)}).

\item \textbf{United States}: The Health Insurance Portability and Accountability Act (HIPAA) Security and Privacy Rules.

\item \textbf{European Union}: The General Data Protection Regulation (GDPR) and European Health Data Space (EHDS) regulations.

\end{itemize}


\par\medskip\textbf{Direct Architectural Impacts}\par


\begin{itemize}


\item Requires downstream security architecture to govern access to protected health information and enforce applicable access controls (\href{\detokenize{../../../../../../../../docs/markdown/01-motivation/requirements-constraints/master-requirements-catalogue.md}}{\texttt{\seqsplit{SEC-001 / SEC-002}}}).

\item Requires downstream operational logging policies and guardrails to protect health information from unauthorized exposure in diagnostic telemetry (\href{\detokenize{../../../../../../../../docs/markdown/01-motivation/requirements-constraints/master-requirements-catalogue.md}}{\texttt{\seqsplit{SEC-010}}}).

\item Requires platform communications and stored data representations to apply appropriate protection mechanisms according to deployment risk and statutory requirements.

\item Informs audit and provenance requirements for demonstrating authorized handling of health information (\href{\detokenize{../../../../../../../../docs/markdown/01-motivation/requirements-constraints/master-requirements-catalogue.md}}{\texttt{\seqsplit{PROV-001}}}).

\item Bounds Platform Driver: \href{\detokenize{../../../../../../../../docs/markdown/01-motivation/drivers-assessments/drivers.md}}{Statutory Health Information Privacy \& Protection}.

\end{itemize}

{\footnotesize \href{../../../../../../../continuum/motivation/constraint/0ed6fb9fcc9547659ce76b917881274e.md}{Canonical record: 0ed6fb9fcc9547659ce76b917881274e}; full qualifications and provenance.}\par
\section*{I.6.6 CST-EXT-002: Applicable National / Jurisdictional Healthcare Identifier Obligations}
\phantomsection
\label{navigation:I.6.6}
\addcontentsline{toc}{section}{I.6.6 CST-EXT-002: Applicable National / Jurisdictional Healthcare Identifier Obligations}

\label{element:edb1db6b18d84487ab48b9416f26e980}
\index{CST-EXT-002: Applicable National / Jurisdictional Healthcare Identifier Obligations (I.6.6)}
{\small Constraint · Accepted · current}\par

\par\medskip\textbf{Categorical Definition}\par


\begin{quote}


\textbf{Where a deployment participates in a national or jurisdictional healthcare identifier ecosystem, Harmonia must conform to applicable identifier representation, validation, authority, and usage obligations.}



\end{quote}


\par\medskip\textbf{Context \& Architectural Boundary}\par


Many modern nations operate statutory healthcare identifier services to uniquely identify patients, clinicians, and health organizations. Participation in these ecosystems is legally regulated: the law dictates how identifiers may be retrieved, verified, stored, cross-referenced, and disclosed. Harmonia must conform strictly to these statutory rules while maintaining internal referential integrity.



\par\medskip\textbf{Representative Instances \& Examples}\par


\begin{itemize}


\item \textbf{Australia}: The \emph{Healthcare Identifiers Act 2010 (Cth)} and HI Service Operating Rules:

\item \textbf{IHI} (Individual Healthcare Identifier): Strict prohibitions against using the IHI as a general-purpose identity token or publishing it outside authorized clinical communication channels.

\item \textbf{HPI-I} (Healthcare Provider Identifier - Individual): Validation against national registers and credentialing rules.

\item \textbf{HPI-O} (Healthcare Provider Identifier - Organisation): Association with authorized electronic service delivery endpoints.

\item \textbf{United Kingdom}: NHS Number validation and Spine integration rules.

\item \textbf{United States}: National Provider Identifier (NPI) specifications.

\end{itemize}


\par\medskip\textbf{Direct Architectural Impacts}\par


\begin{itemize}


\item Requires strict validation of identifier check digits, namespaces, and authority types before accepting clinical records (\href{\detokenize{../../../../../../../../docs/markdown/01-motivation/requirements-constraints/foundational-requirements.md}}{\texttt{\seqsplit{REQ-FND-003}}}).

\item Dictates that external national identifiers must not be conflated with internal synthetic database keys or process identifiers.

\item Requires integration boundaries interacting with national directory or identifier services to adhere to governing access, query, and communication rules.

\item Bounds Platform Driver: \href{\detokenize{../../../../../../../../docs/markdown/01-motivation/drivers-assessments/drivers.md}}{Patient Safety \& Identity Integrity}.

\end{itemize}

{\footnotesize \href{../../../../../../../continuum/motivation/constraint/edb1db6b18d84487ab48b9416f26e980.md}{Canonical record: edb1db6b18d84487ab48b9416f26e980}; full qualifications and provenance.}\par
\section*{I.6.7 CST-EXT-003: Mandated External Interoperability Contracts}
\phantomsection
\label{navigation:I.6.7}
\addcontentsline{toc}{section}{I.6.7 CST-EXT-003: Mandated External Interoperability Contracts}

\label{element:cad6d41364ef4bb1ae0a2a519db11da4}
\index{CST-EXT-003: Mandated External Interoperability Contracts (I.6.7)}
{\small Constraint · Accepted · current}\par

\par\medskip\textbf{Categorical Definition}\par


\begin{quote}


\textbf{Where an external system or ecosystem mandates a protocol, information standard, transport contract, or serialization format, Harmonia must conform to that contract at the applicable interoperability boundary.}



\end{quote}


\par\medskip\textbf{Context \& Architectural Boundary}\par


Harmonia operates as an integration hub connecting hundreds of disparate health systems. The platform cannot dictate the internal wire protocols of external hospitals, laboratory analyzers, or national repositories. At its ingress and egress boundaries, Harmonia must conform strictly to the wire formats, message structures, and transport handshakes mandated by external systems.



However, in accordance with \hyperref[element:eee8c3783d964d18a7f9dcea85fcd4cf]{AX-02 (Standards at Boundary; Sovereignty Within)}, external contracts govern only the external membrane. Once an event crosses into Harmonia, the platform enforces its own canonical domain semantics.



\par\medskip\textbf{Representative Instances \& Examples}\par


\begin{itemize}


\item \textbf{HL7 v2.x over MLLP}: Minimal Lower Layer Protocol (MLLP) framing over raw TCP, utilizing pipe-and-hat (\texttt{\seqsplit{ER7}}) segment syntax for admissions (ADT), orders (ORM), and observations (ORU).

\item \textbf{HL7 FHIR RESTful APIs}: HTTP/JSON or HTTP/XML operations conforming to specific regional profiles (e.g., Australian Core AU-Base / AU-Core, US Core).

\item \textbf{Secure Messaging Delivery (SMD)}: Encrypted Web Services (WS-*) protocols mandated for inter-clinic electronic correspondence.

\item \textbf{DICOM}: Digital Imaging and Communications in Medicine protocol for radiological studies.

\end{itemize}


\par\medskip\textbf{Direct Architectural Impacts}\par


\begin{itemize}


\item Demarcates the boundary between Pylai (the protocol membrane gateways) and core internal platform engines.

\item Governs wire-level error signaling (e.g., mapping internal durable persistence failures to protocol-appropriate error responses).

\item Prevents external wire formats from polluting internal domain models.

\item Imposes on Platform Driver: \href{\detokenize{../../../../../../../../docs/markdown/01-motivation/drivers-assessments/drivers.md}}{Clinical Interoperability Across Heterogeneous Standards}.

\end{itemize}

{\footnotesize \href{../../../../../../../continuum/motivation/constraint/cad6d41364ef4bb1ae0a2a519db11da4.md}{Canonical record: cad6d41364ef4bb1ae0a2a519db11da4}; full qualifications and provenance.}\par
\section*{I.6.8 Dataset Access Has No Runtime Bootstrap Cycle}
\phantomsection
\label{navigation:I.6.8}
\addcontentsline{toc}{section}{I.6.8 Dataset Access Has No Runtime Bootstrap Cycle}

\label{element:c8c479dc9aad450ea58d07dec9aea1a6}
\index{Dataset Access Has No Runtime Bootstrap Cycle (I.6.8)}
{\small Constraint · Adjudicated Domain05 direction — U13/U15; holistic result acceptance pending · current}\par

\begin{quote}


\textbf{The structural contracts necessary to access and manage
Calliope-DataSets SHALL be established through Calliope-Models and the
software build/deployment lifecycle. Runtime Calliope-DataSets SHALL
provide the authoritative mutable content governed by those contracts
and SHALL NOT be required to establish the structural contracts
necessary for its own access.}



\end{quote}


This assertion is normalised from the governed DesignNote under U13/U15 and the ARB release Praxis. Original candidate/release qualifications and separately deferred implementation details remain in force.



Architectural subject: \hyperref[element:466ad8f75a954c0581b5bba090443849]{Calliope}.


{\footnotesize \href{../../../../../../../continuum/motivation/constraint/c8c479dc9aad450ea58d07dec9aea1a6.md}{Canonical record: c8c479dc9aad450ea58d07dec9aea1a6}; full qualifications and provenance.}\par
\section*{I.6.9 Dokimasia Evidence Replaces a Separate Self AuditEvent}
\phantomsection
\label{navigation:I.6.9}
\addcontentsline{toc}{section}{I.6.9 Dokimasia Evidence Replaces a Separate Self AuditEvent}

\label{element:e1ce978457e543d5894de24e9be9ac2c}
\index{Dokimasia Evidence Replaces a Separate Self AuditEvent (I.6.9)}
{\small Constraint · Adjudicated Domain05 direction — U02/U06/U19; holistic result acceptance pending · current}\par

A Dokimasia execution SHALL NOT generate a separate AuditEvent for its own execution. The governed Dokimasia assurance record/result provides the authoritative evidence of that assurance activity.



This assertion is normalised from the governed DesignNote under U02/U06/U19 and the ARB release Praxis. Original candidate/release qualifications and separately deferred implementation details remain in force.



Architectural subject: \hyperref[element:fb3f37a6ac614415ac486a6de9d35994]{Energeia}.


{\footnotesize \href{../../../../../../../continuum/motivation/constraint/e1ce978457e543d5894de24e9be9ac2c.md}{Canonical record: e1ce978457e543d5894de24e9be9ac2c}; full qualifications and provenance.}\par
\section*{I.6.10 Dokimasia Has a Finite Assurance Boundary}
\phantomsection
\label{navigation:I.6.10}
\addcontentsline{toc}{section}{I.6.10 Dokimasia Has a Finite Assurance Boundary}

\label{element:3ba73277abf146728ca8c1f5846227d7}
\index{Dokimasia Has a Finite Assurance Boundary (I.6.10)}
{\small Constraint · Adjudicated Domain05 direction — U02/U06/U19; holistic result acceptance pending · current}\par

A Dokimasia execution SHALL NOT itself be subject to monitoring, assurance or assessment by another Dokimasia execution. This establishes a finite assurance boundary.



This assertion is normalised from the governed DesignNote under U02/U06/U19 and the ARB release Praxis. Original candidate/release qualifications and separately deferred implementation details remain in force.



Architectural subject: \hyperref[element:fb3f37a6ac614415ac486a6de9d35994]{Energeia}.


{\footnotesize \href{../../../../../../../continuum/motivation/constraint/3ba73277abf146728ca8c1f5846227d7.md}{Canonical record: 3ba73277abf146728ca8c1f5846227d7}; full qualifications and provenance.}\par
\section*{I.6.11 Every Petasos Queue Has a Defined Functional Purpose}
\phantomsection
\label{navigation:I.6.11}
\addcontentsline{toc}{section}{I.6.11 Every Petasos Queue Has a Defined Functional Purpose}

\label{element:6fa0f390bfa140aa8bb6acc56f1a5ec6}
\index{Every Petasos Queue Has a Defined Functional Purpose (I.6.11)}
{\small Constraint · Adjudicated Domain05 direction — U09; holistic result acceptance pending · current}\par

A Petasos queue SHALL exist for a defined Harmonia functional purpose.



This assertion is normalised from the governed DesignNote under U09 and the ARB release Praxis. Original candidate/release qualifications and separately deferred implementation details remain in force.



Architectural subject: \hyperref[element:195e5f10db8e419391354197b5d7a8cc]{Petasos}.


{\footnotesize \href{../../../../../../../continuum/motivation/constraint/6fa0f390bfa140aa8bb6acc56f1a5ec6.md}{Canonical record: 6fa0f390bfa140aa8bb6acc56f1a5ec6}; full qualifications and provenance.}\par
\section*{I.6.12 Harmonisation Does Not Resume the Originating Ergon}
\phantomsection
\label{navigation:I.6.12}
\addcontentsline{toc}{section}{I.6.12 Harmonisation Does Not Resume the Originating Ergon}

\label{element:1739ad181a754d01ae07f94de49cd65e}
\index{Harmonisation Does Not Resume the Originating Ergon (I.6.12)}
{\small Constraint · Adjudicated Domain05 direction — U10/U18; holistic result acceptance pending · current}\par

Harmonisation SHALL NOT be implemented by resuming inside the originating Ergon.



This assertion is normalised from the governed DesignNote under U10/U18 and the ARB release Praxis. Original candidate/release qualifications and separately deferred implementation details remain in force.



Architectural subject: \hyperref[element:1af8a56358f7429fabfb57be66542a99]{Mneme}.


{\footnotesize \href{../../../../../../../continuum/motivation/constraint/1739ad181a754d01ae07f94de49cd65e.md}{Canonical record: 1739ad181a754d01ae07f94de49cd65e}; full qualifications and provenance.}\par
\section*{I.6.13 Harmonisation Reuses the Common Execution Framework}
\phantomsection
\label{navigation:I.6.13}
\addcontentsline{toc}{section}{I.6.13 Harmonisation Reuses the Common Execution Framework}

\label{element:59e1a236500240609a78b6529fe8e358}
\index{Harmonisation Reuses the Common Execution Framework (I.6.13)}
{\small Constraint · Adjudicated Domain05 direction — U10/U18; holistic result acceptance pending · current}\par

The RHM SHALL NOT duplicate or embed a second execution framework within Mneme.



This assertion is normalised from the governed DesignNote under U10/U18 and the ARB release Praxis. Original candidate/release qualifications and separately deferred implementation details remain in force.



Architectural subject: \hyperref[element:1af8a56358f7429fabfb57be66542a99]{Mneme}.


{\footnotesize \href{../../../../../../../continuum/motivation/constraint/59e1a236500240609a78b6529fe8e358.md}{Canonical record: 59e1a236500240609a78b6529fe8e358}; full qualifications and provenance.}\par
\section*{I.6.14 Host Components Do Not Reinterpret Themis Policy}
\phantomsection
\label{navigation:I.6.14}
\addcontentsline{toc}{section}{I.6.14 Host Components Do Not Reinterpret Themis Policy}

\label{element:80e5d019948b4c16acb510547d7a3ac6}
\index{Host Components Do Not Reinterpret Themis Policy (I.6.14)}
{\small Constraint · Adjudicated Domain05 direction — U08/U12/U15; holistic result acceptance pending · current}\par

A Harmonia component SHALL NOT implement an independent interpretation of Themis policy in place of the Themis PEP.



This assertion is normalised from the governed DesignNote under U08/U12/U15 and the ARB release Praxis. Original candidate/release qualifications and separately deferred implementation details remain in force.



Architectural subject: \hyperref[element:899d34ec491948b68aad09e507dea6d1]{Themis}.


{\footnotesize \href{../../../../../../../continuum/motivation/constraint/80e5d019948b4c16acb510547d7a3ac6.md}{Canonical record: 80e5d019948b4c16acb510547d7a3ac6}; full qualifications and provenance.}\par
\section*{I.6.15 Independent Audit Targets Are Not Assumed Replicas}
\phantomsection
\label{navigation:I.6.15}
\addcontentsline{toc}{section}{I.6.15 Independent Audit Targets Are Not Assumed Replicas}

\label{element:e84fc7e020474eda99a0e32def0a6371}
\index{Independent Audit Targets Are Not Assumed Replicas (I.6.15)}
{\small Constraint · Adjudicated Domain05 direction — U07/U15; holistic result acceptance pending · current}\par

These targets SHALL NOT be described as replicas unless and until convergence/replication semantics are actually established. For R1.x/R2.x they are \textbf{independent authoritative audit persistence targets}.



This assertion is normalised from the governed DesignNote under U07/U15 and the ARB release Praxis. Original candidate/release qualifications and separately deferred implementation details remain in force.



Architectural subject: \hyperref[element:fefda110e2274884ad58906733606efd]{Kleio}.


{\footnotesize \href{../../../../../../../continuum/motivation/constraint/e84fc7e020474eda99a0e32def0a6371.md}{Canonical record: e84fc7e020474eda99a0e32def0a6371}; full qualifications and provenance.}\par
\section*{I.6.16 Information Topic and Messaging Destination Remain Distinct}
\phantomsection
\label{navigation:I.6.16}
\addcontentsline{toc}{section}{I.6.16 Information Topic and Messaging Destination Remain Distinct}

\label{element:ba909d70784349cd94b8adb4d6a9ab69}
\index{Information Topic and Messaging Destination Remain Distinct (I.6.16)}
{\small Constraint · Adjudicated Domain05 direction — U05/U09/U18; holistic result acceptance pending · current}\par

Care must be taken not to confuse Petasos queue/destination classification with InformationUnit semantic Topic.



This assertion is normalised from the governed DesignNote under U05/U09/U18 and the ARB release Praxis. Original candidate/release qualifications and separately deferred implementation details remain in force.



Architectural subject: \hyperref[element:1af8a56358f7429fabfb57be66542a99]{Mneme}.


{\footnotesize \href{../../../../../../../continuum/motivation/constraint/ba909d70784349cd94b8adb4d6a9ab69.md}{Canonical record: ba909d70784349cd94b8adb4d6a9ab69}; full qualifications and provenance.}\par
\section*{I.6.17 Ingress Context Loading Does Not Interpret Raw Payload Semantics}
\phantomsection
\label{navigation:I.6.17}
\addcontentsline{toc}{section}{I.6.17 Ingress Context Loading Does Not Interpret Raw Payload Semantics}

\label{element:f5af90e09e9a422d9b1c9a7f3c44090f}
\index{Ingress Context Loading Does Not Interpret Raw Payload Semantics (I.6.17)}
{\small Constraint · Adjudicated Domain05 direction — U14/U15; holistic result acceptance pending · current}\par

It must \textbf{not} perform interface-specific semantic interpretation or transformation of the raw payload. That belongs to the paired developer-written \texttt{\seqsplit{Ergo.BusinessLogic}}.



This assertion is normalised from the governed DesignNote under U14/U15 and the ARB release Praxis. Original candidate/release qualifications and separately deferred implementation details remain in force.



Architectural subject: \hyperref[element:e3e4147fea3640198dc251ed41ee854d]{Pylai}.


{\footnotesize \href{../../../../../../../continuum/motivation/constraint/f5af90e09e9a422d9b1c9a7f3c44090f.md}{Canonical record: f5af90e09e9a422d9b1c9a7f3c44090f}; full qualifications and provenance.}\par
\section*{I.6.18 Internal Mneme Responsibilities Are Not Automatically Exposed Services}
\phantomsection
\label{navigation:I.6.18}
\addcontentsline{toc}{section}{I.6.18 Internal Mneme Responsibilities Are Not Automatically Exposed Services}

\label{element:f0a85fcf515946e2b7c88543061ebac1}
\index{Internal Mneme Responsibilities Are Not Automatically Exposed Services (I.6.18)}
{\small Constraint · Adjudicated Domain05 direction — U10/U18; holistic result acceptance pending · current}\par

Other Mneme responsibilities described in this note are supporting/internal responsibilities and SHALL NOT automatically be interpreted as additional exposed services.



This assertion is normalised from the governed DesignNote under U10/U18 and the ARB release Praxis. Original candidate/release qualifications and separately deferred implementation details remain in force.



Architectural subject: \hyperref[element:1af8a56358f7429fabfb57be66542a99]{Mneme}.


{\footnotesize \href{../../../../../../../continuum/motivation/constraint/f0a85fcf515946e2b7c88543061ebac1.md}{Canonical record: f0a85fcf515946e2b7c88543061ebac1}; full qualifications and provenance.}\par
\section*{I.6.19 Journey Observations Remain Distinct from Assurance Inference}
\phantomsection
\label{navigation:I.6.19}
\addcontentsline{toc}{section}{I.6.19 Journey Observations Remain Distinct from Assurance Inference}

\label{element:ae2fa71039034290b019eb1cc6eb1996}
\index{Journey Observations Remain Distinct from Assurance Inference (I.6.19)}
{\small Constraint · Adjudicated Domain05 direction — U06/U09/U19; holistic result acceptance pending · current}\par

Journey SHALL remain a representation of what Harmonia observed. Assurance interpretation, including AI-assisted interpretation, SHALL NOT rewrite Journey into a representation of what an assurance process inferred happened.



This assertion is normalised from the governed DesignNote under U06/U09/U19 and the ARB release Praxis. Original candidate/release qualifications and separately deferred implementation details remain in force.



Architectural subject: \hyperref[element:fb3f37a6ac614415ac486a6de9d35994]{Energeia}.


{\footnotesize \href{../../../../../../../continuum/motivation/constraint/ae2fa71039034290b019eb1cc6eb1996.md}{Canonical record: ae2fa71039034290b019eb1cc6eb1996}; full qualifications and provenance.}\par
\section*{I.6.20 Kleio Constructs AuditEvent Representations for Clients}
\phantomsection
\label{navigation:I.6.20}
\addcontentsline{toc}{section}{I.6.20 Kleio Constructs AuditEvent Representations for Clients}

\label{element:1aa42a9aa5a644bc9072522e0a123352}
\index{Kleio Constructs AuditEvent Representations for Clients (I.6.20)}
{\small Constraint · Adjudicated Domain05 direction — U07/U15; holistic result acceptance pending · current}\par

\begin{quote}


\textbf{Clients SHALL NOT be required to construct FHIR AuditEvent resources. Clients submit auditable occurrences through the Kleio Audit Capture contract. Kleio is responsible for constructing and managing the authoritative FHIR AuditEvent representation.}



\end{quote}


This assertion is normalised from the governed DesignNote under U07/U15 and the ARB release Praxis. Original candidate/release qualifications and separately deferred implementation details remain in force.



Architectural subject: \hyperref[element:fefda110e2274884ad58906733606efd]{Kleio}.


{\footnotesize \href{../../../../../../../continuum/motivation/constraint/1aa42a9aa5a644bc9072522e0a123352.md}{Canonical record: 1aa42a9aa5a644bc9072522e0a123352}; full qualifications and provenance.}\par
\section*{I.6.21 Local PEP Enforcement Is Constrained and Deterministic}
\phantomsection
\label{navigation:I.6.21}
\addcontentsline{toc}{section}{I.6.21 Local PEP Enforcement Is Constrained and Deterministic}

\label{element:8c9ef077cb5d4cc9aacb55d743b1af92}
\index{Local PEP Enforcement Is Constrained and Deterministic (I.6.21)}
{\small Constraint · Adjudicated Domain05 direction — U12/U15/U16; holistic result acceptance pending · current}\par

The PEP is intentionally constrained and deterministic. It SHALL NOT become a general-purpose policy-language interpreter.



This assertion is normalised from the governed DesignNote under U12/U15/U16 and the ARB release Praxis. Original candidate/release qualifications and separately deferred implementation details remain in force.



Architectural subject: \hyperref[element:899d34ec491948b68aad09e507dea6d1]{Themis}.


{\footnotesize \href{../../../../../../../continuum/motivation/constraint/8c9ef077cb5d4cc9aacb55d743b1af92.md}{Canonical record: 8c9ef077cb5d4cc9aacb55d743b1af92}; full qualifications and provenance.}\par
\section*{I.6.22 Mneme Is Insulated from Ponos Implementation Internals}
\phantomsection
\label{navigation:I.6.22}
\addcontentsline{toc}{section}{I.6.22 Mneme Is Insulated from Ponos Implementation Internals}

\label{element:c30ae90fc19b465982b28d61ac1edca5}
\index{Mneme Is Insulated from Ponos Implementation Internals (I.6.22)}
{\small Constraint · Adjudicated Domain05 direction — U10/U18; holistic result acceptance pending · current}\par

The cross-component dependency SHALL remain narrow and architectural rather than coupling Mneme to Ponos implementation internals.



This assertion is normalised from the governed DesignNote under U10/U18 and the ARB release Praxis. Original candidate/release qualifications and separately deferred implementation details remain in force.



Architectural subject: \hyperref[element:1af8a56358f7429fabfb57be66542a99]{Mneme}.


{\footnotesize \href{../../../../../../../continuum/motivation/constraint/c30ae90fc19b465982b28d61ac1edca5.md}{Canonical record: c30ae90fc19b465982b28d61ac1edca5}; full qualifications and provenance.}\par
\section*{I.6.23 Negative Criteria Require a Demonstrated Business Requirement}
\phantomsection
\label{navigation:I.6.23}
\addcontentsline{toc}{section}{I.6.23 Negative Criteria Require a Demonstrated Business Requirement}

\label{element:820e647c43614052aa25e385d6fe1dac}
\index{Negative Criteria Require a Demonstrated Business Requirement (I.6.23)}
{\small Constraint · Adjudicated Domain05 direction — U12/U15/U16; holistic result acceptance pending · current}\par

The Candidate Architecture does not introduce explicit negative/DENY Criteria. Such a mechanism SHALL NOT be introduced unless a demonstrated business requirement justifies the additional policy semantics.



This assertion is normalised from the governed DesignNote under U12/U15/U16 and the ARB release Praxis. Original candidate/release qualifications and separately deferred implementation details remain in force.



Architectural subject: \hyperref[element:899d34ec491948b68aad09e507dea6d1]{Themis}.


{\footnotesize \href{../../../../../../../continuum/motivation/constraint/820e647c43614052aa25e385d6fe1dac.md}{Canonical record: 820e647c43614052aa25e385d6fe1dac}; full qualifications and provenance.}\par
\section*{I.6.24 New Policy Context Requires Authoritative PDP Processing}
\phantomsection
\label{navigation:I.6.24}
\addcontentsline{toc}{section}{I.6.24 New Policy Context Requires Authoritative PDP Processing}

\label{element:80629d4c229e4b7095d011bf77d704aa}
\index{New Policy Context Requires Authoritative PDP Processing (I.6.24)}
{\small Constraint · Adjudicated Domain05 direction — U08/U12/U15; holistic result acceptance pending · current}\par

If the PDP cannot establish the required new PEP-Context, the activity SHALL NOT proceed beyond the point requiring that context and must wait until authoritative policy context can be established.



This assertion is normalised from the governed DesignNote under U08/U12/U15 and the ARB release Praxis. Original candidate/release qualifications and separately deferred implementation details remain in force.



Architectural subject: \hyperref[element:899d34ec491948b68aad09e507dea6d1]{Themis}.


{\footnotesize \href{../../../../../../../continuum/motivation/constraint/80629d4c229e4b7095d011bf77d704aa.md}{Canonical record: 80629d4c229e4b7095d011bf77d704aa}; full qualifications and provenance.}\par
\section*{I.6.25 PDP Unavailability Is Not Policy Denial}
\phantomsection
\label{navigation:I.6.25}
\addcontentsline{toc}{section}{I.6.25 PDP Unavailability Is Not Policy Denial}

\label{element:8d0e87312d1f4d6f9298843ce840e654}
\index{PDP Unavailability Is Not Policy Denial (I.6.25)}
{\small Constraint · Adjudicated Domain05 direction — U08/U12/U15; holistic result acceptance pending · current}\par

\begin{enumerate}


\item \textbf{PDP unavailability SHALL NOT be represented as a policy denial.}

\end{enumerate}


This assertion is normalised from the governed DesignNote under U08/U12/U15 and the ARB release Praxis. Original candidate/release qualifications and separately deferred implementation details remain in force.



Architectural subject: \hyperref[element:899d34ec491948b68aad09e507dea6d1]{Themis}.


{\footnotesize \href{../../../../../../../continuum/motivation/constraint/8d0e87312d1f4d6f9298843ce840e654.md}{Canonical record: 8d0e87312d1f4d6f9298843ce840e654}; full qualifications and provenance.}\par
\section*{I.6.26 Persistence Preserves Applicable Compartments}
\phantomsection
\label{navigation:I.6.26}
\addcontentsline{toc}{section}{I.6.26 Persistence Preserves Applicable Compartments}

\label{element:6e963678cf6f4c3488c798d6afc3b910}
\index{Persistence Preserves Applicable Compartments (I.6.26)}
{\small Constraint · Adjudicated Domain05 direction — U11; holistic result acceptance pending · current}\par

Persistence SHALL NOT accidentally strip, weaken or lose the applicable compartment classification.



This assertion is normalised from the governed DesignNote under U11 and the ARB release Praxis. Original candidate/release qualifications and separately deferred implementation details remain in force.



Architectural subject: \hyperref[element:7985e647bf9e4ca3ba70faf9e3e4ae02]{Mnemosyne}.


{\footnotesize \href{../../../../../../../continuum/motivation/constraint/6e963678cf6f4c3488c798d6afc3b910.md}{Canonical record: 6e963678cf6f4c3488c798d6afc3b910}; full qualifications and provenance.}\par
\section*{I.6.27 PetasosParcel Is Transported by Value}
\phantomsection
\label{navigation:I.6.27}
\addcontentsline{toc}{section}{I.6.27 PetasosParcel Is Transported by Value}

\label{element:b599ebffaf2f4e9e8deaec2ed98f2cd7}
\index{PetasosParcel Is Transported by Value (I.6.27)}
{\small Constraint · Adjudicated Domain05 direction — U09; holistic result acceptance pending · current}\par

Petasos SHALL NOT be described as transporting only references. It transports PetasosParcel instances by value.



This assertion is normalised from the governed DesignNote under U09 and the ARB release Praxis. Original candidate/release qualifications and separately deferred implementation details remain in force.



Architectural subject: \hyperref[element:195e5f10db8e419391354197b5d7a8cc]{Petasos}.


{\footnotesize \href{../../../../../../../continuum/motivation/constraint/b599ebffaf2f4e9e8deaec2ed98f2cd7.md}{Canonical record: b599ebffaf2f4e9e8deaec2ed98f2cd7}; full qualifications and provenance.}\par
\section*{I.6.28 R1 and R2 Observation Does Not Authorise Autonomous Adaptation}
\phantomsection
\label{navigation:I.6.28}
\addcontentsline{toc}{section}{I.6.28 R1 and R2 Observation Does Not Authorise Autonomous Adaptation}

\label{element:7ca8481aa04a4dfdb34766511e002a5b}
\index{R1 and R2 Observation Does Not Authorise Autonomous Adaptation (I.6.28)}
{\small Constraint · Adjudicated Domain05 direction — U06/U09/U19; holistic result acceptance pending · current}\par

R1.x/R2.x SHALL NOT introduce autonomous adaptation merely because Energeia can observe Harmonia operation.



This assertion is normalised from the governed DesignNote under U06/U09/U19 and the ARB release Praxis. Original candidate/release qualifications and separately deferred implementation details remain in force.



Architectural subject: \hyperref[element:fb3f37a6ac614415ac486a6de9d35994]{Energeia}.


{\footnotesize \href{../../../../../../../continuum/motivation/constraint/7ca8481aa04a4dfdb34766511e002a5b.md}{Canonical record: 7ca8481aa04a4dfdb34766511e002a5b}; full qualifications and provenance.}\par
\section*{I.6.29 Revocation Is Distinct from Temporal Expiry}
\phantomsection
\label{navigation:I.6.29}
\addcontentsline{toc}{section}{I.6.29 Revocation Is Distinct from Temporal Expiry}

\label{element:462d94a9a14d4eb8a74710fff4b12402}
\index{Revocation Is Distinct from Temporal Expiry (I.6.29)}
{\small Constraint · Adjudicated Domain05 direction — U12/U15/U16; holistic result acceptance pending · current}\par

Revocation SHALL NOT be treated merely as temporal expiry.



This assertion is normalised from the governed DesignNote under U12/U15/U16 and the ARB release Praxis. Original candidate/release qualifications and separately deferred implementation details remain in force.



Architectural subject: \hyperref[element:899d34ec491948b68aad09e507dea6d1]{Themis}.


{\footnotesize \href{../../../../../../../continuum/motivation/constraint/462d94a9a14d4eb8a74710fff4b12402.md}{Canonical record: 462d94a9a14d4eb8a74710fff4b12402}; full qualifications and provenance.}\par
\section*{I.6.30 Runtime Vocabulary Membership Is Authoritative}
\phantomsection
\label{navigation:I.6.30}
\addcontentsline{toc}{section}{I.6.30 Runtime Vocabulary Membership Is Authoritative}

\label{element:872a5f1aa5a44fd4ba68d793c022adf2}
\index{Runtime Vocabulary Membership Is Authoritative (I.6.30)}
{\small Constraint · Adjudicated Domain05 direction — U13/U15; holistic result acceptance pending · current}\par

Software built against an earlier representation SHALL NOT automatically
treat a runtime value as invalid merely because that value was unknown
at build time.



This assertion is normalised from the governed DesignNote under U13/U15 and the ARB release Praxis. Original candidate/release qualifications and separately deferred implementation details remain in force.



Architectural subject: \hyperref[element:466ad8f75a954c0581b5bba090443849]{Calliope}.


{\footnotesize \href{../../../../../../../continuum/motivation/constraint/872a5f1aa5a44fd4ba68d793c022adf2.md}{Canonical record: 872a5f1aa5a44fd4ba68d793c022adf2}; full qualifications and provenance.}\par
\section*{I.6.31 Semantic Persistence Domains Do Not Prescribe Physical Topology}
\phantomsection
\label{navigation:I.6.31}
\addcontentsline{toc}{section}{I.6.31 Semantic Persistence Domains Do Not Prescribe Physical Topology}

\label{element:eca9a99c2ad44a1ea7833d2b0c5a3e93}
\index{Semantic Persistence Domains Do Not Prescribe Physical Topology (I.6.31)}
{\small Constraint · Adjudicated Domain05 direction — U11; holistic result acceptance pending · current}\par

They SHALL NOT automatically be interpreted as three databases, schemas, services, workloads or physical deployments.



This assertion is normalised from the governed DesignNote under U11 and the ARB release Praxis. Original candidate/release qualifications and separately deferred implementation details remain in force.



Architectural subject: \hyperref[element:7985e647bf9e4ca3ba70faf9e3e4ae02]{Mnemosyne}.


{\footnotesize \href{../../../../../../../continuum/motivation/constraint/eca9a99c2ad44a1ea7833d2b0c5a3e93.md}{Canonical record: eca9a99c2ad44a1ea7833d2b0c5a3e93}; full qualifications and provenance.}\par
\section*{I.6.32 Separate Models and DataSets Lifecycles}
\phantomsection
\label{navigation:I.6.32}
\addcontentsline{toc}{section}{I.6.32 Separate Models and DataSets Lifecycles}

\label{element:5c15fa0761da496b8194f3a6f742ecb5}
\index{Separate Models and DataSets Lifecycles (I.6.32)}
{\small Constraint · Adjudicated Domain05 direction — U13/U15; holistic result acceptance pending · current}\par

The two modules collaborate but SHALL NOT be treated as having a single
lifecycle.



This assertion is normalised from the governed DesignNote under U13/U15 and the ARB release Praxis. Original candidate/release qualifications and separately deferred implementation details remain in force.



Architectural subject: \hyperref[element:466ad8f75a954c0581b5bba090443849]{Calliope}.


{\footnotesize \href{../../../../../../../continuum/motivation/constraint/5c15fa0761da496b8194f3a6f742ecb5.md}{Canonical record: 5c15fa0761da496b8194f3a6f742ecb5}; full qualifications and provenance.}\par
\section*{I.6.33 Specialised Execution Reuses Common Platform Capabilities}
\phantomsection
\label{navigation:I.6.33}
\addcontentsline{toc}{section}{I.6.33 Specialised Execution Reuses Common Platform Capabilities}

\label{element:968e208d336c48b781a4eabd76e6337b}
\index{Specialised Execution Reuses Common Platform Capabilities (I.6.33)}
{\small Constraint · Adjudicated Domain05 direction — U02/U06/U19; holistic result acceptance pending · current}\par

Specialised Harmonia execution constructs SHALL reuse common platform management, assurance, reporting and integration capabilities inherited through their base architectural constructs. Specialisation introduces only behaviour and information semantics necessary for the specialised role.



This assertion is normalised from the governed DesignNote under U02/U06/U19 and the ARB release Praxis. Original candidate/release qualifications and separately deferred implementation details remain in force.



Architectural subject: \hyperref[element:fb3f37a6ac614415ac486a6de9d35994]{Energeia}.


{\footnotesize \href{../../../../../../../continuum/motivation/constraint/968e208d336c48b781a4eabd76e6337b.md}{Canonical record: 968e208d336c48b781a4eabd76e6337b}; full qualifications and provenance.}\par
\section*{I.6.34 Unavailable Authoritative Audit Capture Fails Closed}
\phantomsection
\label{navigation:I.6.34}
\addcontentsline{toc}{section}{I.6.34 Unavailable Authoritative Audit Capture Fails Closed}

\label{element:6c5d5f1d872642f2b7770248d4e1e1f4}
\index{Unavailable Authoritative Audit Capture Fails Closed (I.6.34)}
{\small Constraint · Adjudicated Domain05 direction — U07/U15; holistic result acceptance pending · current}\par

The inability to establish required audit evidence is therefore a platform-level loss of essential Harmonia service and SHALL result in fail-closed behaviour.



This assertion is normalised from the governed DesignNote under U07/U15 and the ARB release Praxis. Original candidate/release qualifications and separately deferred implementation details remain in force.



Architectural subject: \hyperref[element:fefda110e2274884ad58906733606efd]{Kleio}.


{\footnotesize \href{../../../../../../../continuum/motivation/constraint/6c5d5f1d872642f2b7770248d4e1e1f4.md}{Canonical record: 6c5d5f1d872642f2b7770248d4e1e1f4}; full qualifications and provenance.}\par
\section*{I.6.35 Unknown Required Security Attributes Fail Closed}
\phantomsection
\label{navigation:I.6.35}
\addcontentsline{toc}{section}{I.6.35 Unknown Required Security Attributes Fail Closed}

\label{element:7ea0d91ed25848f4add672ddf08f42a2}
\index{Unknown Required Security Attributes Fail Closed (I.6.35)}
{\small Constraint · Adjudicated Domain05 direction — U12/U15/U16; holistic result acceptance pending · current}\par

If required effective security attributes cannot be established, enforcement SHALL fail closed.



This assertion is normalised from the governed DesignNote under U12/U15/U16 and the ARB release Praxis. Original candidate/release qualifications and separately deferred implementation details remain in force.



Architectural subject: \hyperref[element:899d34ec491948b68aad09e507dea6d1]{Themis}.


{\footnotesize \href{../../../../../../../continuum/motivation/constraint/7ea0d91ed25848f4add672ddf08f42a2.md}{Canonical record: 7ea0d91ed25848f4add672ddf08f42a2}; full qualifications and provenance.}\par
\section*{I.6.36 Valid Context Permits Enforcement During PDP Unavailability}
\phantomsection
\label{navigation:I.6.36}
\addcontentsline{toc}{section}{I.6.36 Valid Context Permits Enforcement During PDP Unavailability}

\label{element:5a6285eb1a1646e1badd693b8c175f01}
\index{Valid Context Permits Enforcement During PDP Unavailability (I.6.36)}
{\small Constraint · Adjudicated Domain05 direction — U08/U12/U15; holistic result acceptance pending · current}\par

Themis PDP availability SHALL NOT, by itself, interrupt an in-flight activity that already possesses a valid PEP-Context sufficient for the policy enforcement decisions required by that activity.



This assertion is normalised from the governed DesignNote under U08/U12/U15 and the ARB release Praxis. Original candidate/release qualifications and separately deferred implementation details remain in force.



Architectural subject: \hyperref[element:899d34ec491948b68aad09e507dea6d1]{Themis}.


{\footnotesize \href{../../../../../../../continuum/motivation/constraint/5a6285eb1a1646e1badd693b8c175f01.md}{Canonical record: 5a6285eb1a1646e1badd693b8c175f01}; full qualifications and provenance.}\par
\section*{I.6.37 [Mandate] Architectural Decision Traceability}
\phantomsection
\label{navigation:I.6.37}
\addcontentsline{toc}{section}{I.6.37 [Mandate] Architectural Decision Traceability}

\label{element:ae28fe841f6b476da4393c5e3c6a2eab}
\index{[Mandate] Architectural Decision Traceability (I.6.37)}
{\small Constraint / Mandate · Solution Architect-reviewed Iteration 001; authorised ADD; resulting increment review pending · current}\par

Trace architectural decisions to governing principles, constraints and requirements



Source entry: \textbf{NFR-002 — Architectural Decision Traceability}. Primary subject: \texttt{\seqsplit{G01.01}}.



The source classification is Constraint «Mandate». The source identifier is a scoped register reference and does not determine architectural type. Applicability, release scope and unresolved thresholds remain as stated above.


{\footnotesize \href{../../../../../../../continuum/motivation/constraint/mandate/ae28fe841f6b476da4393c5e3c6a2eab.md}{Canonical record: ae28fe841f6b476da4393c5e3c6a2eab}; full qualifications and provenance.}\par
\section*{I.6.38 [Mandate] Implementation Reconciliation}
\phantomsection
\label{navigation:I.6.38}
\addcontentsline{toc}{section}{I.6.38 [Mandate] Implementation Reconciliation}

\label{element:5acccfa8cf634669a26bb31010926720}
\index{[Mandate] Implementation Reconciliation (I.6.38)}
{\small Constraint / Mandate · Solution Architect-reviewed Iteration 001; authorised ADD; resulting increment review pending · current}\par

Reconcile implementation changes against governed architecture



Source entry: \textbf{NFR-003 — Implementation Reconciliation}. Primary subject: \texttt{\seqsplit{G01.01}}.



The source classification is Constraint «Mandate». The source identifier is a scoped register reference and does not determine architectural type. Applicability, release scope and unresolved thresholds remain as stated above.


{\footnotesize \href{../../../../../../../continuum/motivation/constraint/mandate/5acccfa8cf634669a26bb31010926720.md}{Canonical record: 5acccfa8cf634669a26bb31010926720}; full qualifications and provenance.}\par
\section*{I.6.39 [Mandate] Governed Architecture Review}
\phantomsection
\label{navigation:I.6.39}
\addcontentsline{toc}{section}{I.6.39 [Mandate] Governed Architecture Review}

\label{element:d30a62f256ae44599841362118a84ec8}
\index{[Mandate] Governed Architecture Review (I.6.39)}
{\small Constraint / Mandate · Solution Architect-reviewed Iteration 001; authorised ADD; resulting increment review pending · current}\par

Subject architectural changes to governed review and acceptance



Source entry: \textbf{NFR-006 — Governed Architecture Review}. Primary subject: \texttt{\seqsplit{G01.04}}.



The source classification is Constraint «Mandate». The source identifier is a scoped register reference and does not determine architectural type. Applicability, release scope and unresolved thresholds remain as stated above.


{\footnotesize \href{../../../../../../../continuum/motivation/constraint/mandate/d30a62f256ae44599841362118a84ec8.md}{Canonical record: d30a62f256ae44599841362118a84ec8}; full qualifications and provenance.}\par
\section*{I.6.40 [Mandate] Immutable Information Protection}
\phantomsection
\label{navigation:I.6.40}
\addcontentsline{toc}{section}{I.6.40 [Mandate] Immutable Information Protection}

\label{element:4a25d9c8aaa9418fbddb5bff7302c46d}
\index{[Mandate] Immutable Information Protection (I.6.40)}
{\small Constraint / Mandate · Solution Architect-reviewed Iteration 001; authorised ADD; resulting increment review pending · current}\par

Information designated as immutable shall not be modified under any circumstances within the scope of its immutability designation.



Source entry: \textbf{NFR-023 — Immutable Information Protection}. Primary subject: \texttt{\seqsplit{G02.02}}.



The source classification is Constraint «Mandate». The source identifier is a scoped register reference and does not determine architectural type. Applicability, release scope and unresolved thresholds remain as stated above.


{\footnotesize \href{../../../../../../../continuum/motivation/constraint/mandate/4a25d9c8aaa9418fbddb5bff7302c46d.md}{Canonical record: 4a25d9c8aaa9418fbddb5bff7302c46d}; full qualifications and provenance.}\par
\section*{I.6.41 [Mandate] Accepted Work Durability}
\phantomsection
\label{navigation:I.6.41}
\addcontentsline{toc}{section}{I.6.41 [Mandate] Accepted Work Durability}

\label{element:eabd38371b6940289d9a4bf8e9773942}
\index{[Mandate] Accepted Work Durability (I.6.41)}
{\small Constraint / Mandate · Solution Architect-reviewed Iteration 001; authorised ADD; resulting increment review pending · current}\par

Recoverable processing failures shall not result in the loss of accepted governed work.



Source entry: \textbf{NFR-029 — Accepted Work Durability}. Primary subject: \texttt{\seqsplit{G07.05}}.



The source classification is Constraint «Mandate». The source identifier is a scoped register reference and does not determine architectural type. Applicability, release scope and unresolved thresholds remain as stated above.


{\footnotesize \href{../../../../../../../continuum/motivation/constraint/mandate/eabd38371b6940289d9a4bf8e9773942.md}{Canonical record: eabd38371b6940289d9a4bf8e9773942}; full qualifications and provenance.}\par
\section*{I.6.42 [Mandate] Controlled Processing Recovery}
\phantomsection
\label{navigation:I.6.42}
\addcontentsline{toc}{section}{I.6.42 [Mandate] Controlled Processing Recovery}

\label{element:014a4549398d49ee81eadb7108598f48}
\index{[Mandate] Controlled Processing Recovery (I.6.42)}
{\small Constraint / Mandate · Solution Architect-reviewed Iteration 001; authorised ADD; resulting increment review pending · current}\par

Support controlled recovery or retry of failed operations



Source entry: \textbf{NFR-030 — Controlled Processing Recovery}. Primary subject: \texttt{\seqsplit{G07.03}}.



The source classification is Constraint «Mandate». The source identifier is a scoped register reference and does not determine architectural type. Applicability, release scope and unresolved thresholds remain as stated above.


{\footnotesize \href{../../../../../../../continuum/motivation/constraint/mandate/014a4549398d49ee81eadb7108598f48.md}{Canonical record: 014a4549398d49ee81eadb7108598f48}; full qualifications and provenance.}\par
\section*{I.6.43 [Mandate] Governed Information Authorisation}
\phantomsection
\label{navigation:I.6.43}
\addcontentsline{toc}{section}{I.6.43 [Mandate] Governed Information Authorisation}

\label{element:8812fcd52d3f4afcbeef41b60bda5874}
\index{[Mandate] Governed Information Authorisation (I.6.43)}
{\small Constraint / Mandate · Solution Architect-reviewed Iteration 001; authorised ADD; resulting increment review pending · current}\par

Access to governed information shall be subject to authorisation.



Source entry: \textbf{NFR-035 — Governed Information Authorisation}. Primary subject: \texttt{\seqsplit{G03.03}}.



The source classification is Constraint «Mandate». The source identifier is a scoped register reference and does not determine architectural type. Applicability, release scope and unresolved thresholds remain as stated above.


{\footnotesize \href{../../../../../../../continuum/motivation/constraint/mandate/8812fcd52d3f4afcbeef41b60bda5874.md}{Canonical record: 8812fcd52d3f4afcbeef41b60bda5874}; full qualifications and provenance.}\par
\section*{I.6.44 [Mandate] Authenticated Identity}
\phantomsection
\label{navigation:I.6.44}
\addcontentsline{toc}{section}{I.6.44 [Mandate] Authenticated Identity}

\label{element:72b6d13c85944e2d93efc6c5f5a53d82}
\index{[Mandate] Authenticated Identity (I.6.44)}
{\small Constraint / Mandate · Solution Architect-reviewed Iteration 001; authorised ADD; resulting increment review pending · current}\par

Establish authenticated identity suitable for authorisation



Source entry: \textbf{NFR-037 — Authenticated Identity}. Primary subject: \texttt{\seqsplit{G03.03}}.



The source classification is Constraint «Mandate». The source identifier is a scoped register reference and does not determine architectural type. Applicability, release scope and unresolved thresholds remain as stated above.


{\footnotesize \href{../../../../../../../continuum/motivation/constraint/mandate/72b6d13c85944e2d93efc6c5f5a53d82.md}{Canonical record: 72b6d13c85944e2d93efc6c5f5a53d82}; full qualifications and provenance.}\par
\section*{I.6.45 [Mandate] Authorisation Boundary Enforcement}
\phantomsection
\label{navigation:I.6.45}
\addcontentsline{toc}{section}{I.6.45 [Mandate] Authorisation Boundary Enforcement}

\label{element:e04672ef7f2341b68c09a0348742a861}
\index{[Mandate] Authorisation Boundary Enforcement (I.6.45)}
{\small Constraint / Mandate · Solution Architect-reviewed Iteration 001; authorised ADD; resulting increment review pending · current}\par

Enforce authorisation at defined architectural boundaries



Source entry: \textbf{NFR-038 — Authorisation Boundary Enforcement}. Primary subject: \texttt{\seqsplit{G03.04}}.



The source classification is Constraint «Mandate». The source identifier is a scoped register reference and does not determine architectural type. Applicability, release scope and unresolved thresholds remain as stated above.


{\footnotesize \href{../../../../../../../continuum/motivation/constraint/mandate/e04672ef7f2341b68c09a0348742a861.md}{Canonical record: e04672ef7f2341b68c09a0348742a861}; full qualifications and provenance.}\par
\section*{I.6.46 [Mandate] Provider-Side Security Enforcement}
\phantomsection
\label{navigation:I.6.46}
\addcontentsline{toc}{section}{I.6.46 [Mandate] Provider-Side Security Enforcement}

\label{element:d93b50d0abe04c47bf09759627785745}
\index{[Mandate] Provider-Side Security Enforcement (I.6.46)}
{\small Constraint / Mandate · Solution Architect-reviewed Iteration 001; authorised ADD; resulting increment review pending · current}\par

Do not rely solely on consumer-side security enforcement



Source entry: \textbf{NFR-040 — Provider-Side Security Enforcement}. Primary subject: \texttt{\seqsplit{G03.04}}.



The source classification is Constraint «Mandate». The source identifier is a scoped register reference and does not determine architectural type. Applicability, release scope and unresolved thresholds remain as stated above.


{\footnotesize \href{../../../../../../../continuum/motivation/constraint/mandate/d93b50d0abe04c47bf09759627785745.md}{Canonical record: d93b50d0abe04c47bf09759627785745}; full qualifications and provenance.}\par
\section*{I.6.47 [Mandate] Governed Audit Generation}
\phantomsection
\label{navigation:I.6.47}
\addcontentsline{toc}{section}{I.6.47 [Mandate] Governed Audit Generation}

\label{element:d388ee6fb1d24ca886b2802cfb42c30e}
\index{[Mandate] Governed Audit Generation (I.6.47)}
{\small Constraint / Mandate · Solution Architect-reviewed Iteration 001; authorised ADD; resulting increment review pending · current}\par

All activities subject to Harmonia's governed audit obligations shall generate appropriate audit records.



Source entry: \textbf{NFR-043 — Governed Audit Generation}. Primary subject: \texttt{\seqsplit{G04.01}}.



The source classification is Constraint «Mandate». The source identifier is a scoped register reference and does not determine architectural type. Applicability, release scope and unresolved thresholds remain as stated above.


{\footnotesize \href{../../../../../../../continuum/motivation/constraint/mandate/d388ee6fb1d24ca886b2802cfb42c30e.md}{Canonical record: d388ee6fb1d24ca886b2802cfb42c30e}; full qualifications and provenance.}\par
\section*{I.6.48 [Mandate] Audit Record Protection}
\phantomsection
\label{navigation:I.6.48}
\addcontentsline{toc}{section}{I.6.48 [Mandate] Audit Record Protection}

\label{element:de15359269cf45adb646a435c18306f8}
\index{[Mandate] Audit Record Protection (I.6.48)}
{\small Constraint / Mandate · Solution Architect-reviewed Iteration 001; authorised ADD; resulting increment review pending · current}\par

Protect audit records against unauthorised modification



Source entry: \textbf{NFR-045 — Audit Record Protection}. Primary subject: \texttt{\seqsplit{G04.03}}.



The source classification is Constraint «Mandate». The source identifier is a scoped register reference and does not determine architectural type. Applicability, release scope and unresolved thresholds remain as stated above.


{\footnotesize \href{../../../../../../../continuum/motivation/constraint/mandate/de15359269cf45adb646a435c18306f8.md}{Canonical record: de15359269cf45adb646a435c18306f8}; full qualifications and provenance.}\par
\section*{I.6.49 [Mandate] Governed Exchange Representation}
\phantomsection
\label{navigation:I.6.49}
\addcontentsline{toc}{section}{I.6.49 [Mandate] Governed Exchange Representation}

\label{element:e9d06b3fe6b744cc86bd2a136e02165a}
\index{[Mandate] Governed Exchange Representation (I.6.49)}
{\small Constraint / Mandate · Solution Architect-reviewed Iteration 001; authorised ADD; resulting increment review pending · current}\par

Information exchanged with Harmonia shall either be encapsulated within a governed information structure or mapped/transformed into a governed information representation.



Source entry: \textbf{NFR-052 — Governed Exchange Representation}. Primary subject: \texttt{\seqsplit{G05.03}}.



The source classification is Constraint «Mandate». The source identifier is a scoped register reference and does not determine architectural type. Applicability, release scope and unresolved thresholds remain as stated above.


{\footnotesize \href{../../../../../../../continuum/motivation/constraint/mandate/e9d06b3fe6b744cc86bd2a136e02165a.md}{Canonical record: e9d06b3fe6b744cc86bd2a136e02165a}; full qualifications and provenance.}\par
\section*{I.6.50 [Mandate] Native Operational Observability}
\phantomsection
\label{navigation:I.6.50}
\addcontentsline{toc}{section}{I.6.50 [Mandate] Native Operational Observability}

\label{element:151906f8acad4519b7f9204cce9fc758}
\index{[Mandate] Native Operational Observability (I.6.50)}
{\small Constraint / Mandate · Solution Architect-reviewed Iteration 001; authorised ADD; resulting increment review pending · current}\par

Observability is a fundamental Harmonia capability. Every component shall expose sufficient operational information to enable effective monitoring, diagnosis and assessment of its operational state.



Source entry: \textbf{NFR-055 — Native Operational Observability}. Primary subject: \texttt{\seqsplit{G09.01}}.



The source classification is Constraint «Mandate». The source identifier is a scoped register reference and does not determine architectural type. Applicability, release scope and unresolved thresholds remain as stated above.


{\footnotesize \href{../../../../../../../continuum/motivation/constraint/mandate/151906f8acad4519b7f9204cce9fc758.md}{Canonical record: 151906f8acad4519b7f9204cce9fc758}; full qualifications and provenance.}\par
\section*{I.6.51 [Mandate] Dependency Health Monitoring}
\phantomsection
\label{navigation:I.6.51}
\addcontentsline{toc}{section}{I.6.51 [Mandate] Dependency Health Monitoring}

\label{element:c5c8a76fbcbb466bb94f27b4863fb84c}
\index{[Mandate] Dependency Health Monitoring (I.6.51)}
{\small Constraint / Mandate · Solution Architect-reviewed Iteration 001; authorised ADD; resulting increment review pending · current}\par

Component deployment and execution dependencies shall be explicitly identified and monitored, with their operational state incorporated into solution operational readiness reporting.



Source entry: \textbf{NFR-056 — Dependency Health Monitoring}. Primary subject: \texttt{\seqsplit{G09.01}}.



The source classification is Constraint «Mandate». The source identifier is a scoped register reference and does not determine architectural type. Applicability, release scope and unresolved thresholds remain as stated above.


{\footnotesize \href{../../../../../../../continuum/motivation/constraint/mandate/c5c8a76fbcbb466bb94f27b4863fb84c.md}{Canonical record: c5c8a76fbcbb466bb94f27b4863fb84c}; full qualifications and provenance.}\par
\section*{I.6.52 [Mandate] Concurrent Execution Units}
\phantomsection
\label{navigation:I.6.52}
\addcontentsline{toc}{section}{I.6.52 [Mandate] Concurrent Execution Units}

\label{element:12466ce2701146c1b52a9c68c2fb426c}
\index{[Mandate] Concurrent Execution Units (I.6.52)}
{\small Constraint / Mandate · Solution Architect-reviewed Iteration 001; authorised ADD; resulting increment review pending · current}\par

Where practical, Harmonia processing activities shall be encapsulated within independently managed execution units that support concurrent operation.



Source entry: \textbf{NFR-060 — Concurrent Execution Units}. Primary subject: \texttt{\seqsplit{G06.03}}.



The source classification is Constraint «Mandate». The source identifier is a scoped register reference and does not determine architectural type. Applicability, release scope and unresolved thresholds remain as stated above.


{\footnotesize \href{../../../../../../../continuum/motivation/constraint/mandate/12466ce2701146c1b52a9c68c2fb426c.md}{Canonical record: 12466ce2701146c1b52a9c68c2fb426c}; full qualifications and provenance.}\par
\section*{I.6.53 [Mandate] Elastic Processing Capacity}
\phantomsection
\label{navigation:I.6.53}
\addcontentsline{toc}{section}{I.6.53 [Mandate] Elastic Processing Capacity}

\label{element:ea33d8db942f4b5b9e322d74cbfe975b}
\index{[Mandate] Elastic Processing Capacity (I.6.53)}
{\small Constraint / Mandate · Solution Architect-reviewed Iteration 001; authorised ADD; resulting increment review pending · current}\par

Harmonia's workflow, message-processing and Digital Twin capabilities shall support the addition and removal of processing resources, where practical, without requiring fundamental changes to their operational behaviour.



Source entry: \textbf{NFR-061 — Elastic Processing Capacity}. Primary subject: \texttt{\seqsplit{G08.03}}.



The source classification is Constraint «Mandate». The source identifier is a scoped register reference and does not determine architectural type. Applicability, release scope and unresolved thresholds remain as stated above.


{\footnotesize \href{../../../../../../../continuum/motivation/constraint/mandate/ea33d8db942f4b5b9e322d74cbfe975b.md}{Canonical record: ea33d8db942f4b5b9e322d74cbfe975b}; full qualifications and provenance.}\par
\section*{I.6.54 [Mandate] Concurrent Service Processing}
\phantomsection
\label{navigation:I.6.54}
\addcontentsline{toc}{section}{I.6.54 [Mandate] Concurrent Service Processing}

\label{element:338a4cb1c2444ec2aa7703fcd9f28247}
\index{[Mandate] Concurrent Service Processing (I.6.54)}
{\small Constraint / Mandate · Solution Architect-reviewed Iteration 001; authorised ADD; resulting increment review pending · current}\par

Harmonia services responsible for receiving, processing or routing information, or modifying persisted information, shall support appropriate concurrent execution mechanisms to minimise processing bottlenecks and avoid single points of failure.



Source entry: \textbf{NFR-062 — Concurrent Service Processing}. Primary subject: \texttt{\seqsplit{G06.03}}.



The source classification is Constraint «Mandate». The source identifier is a scoped register reference and does not determine architectural type. Applicability, release scope and unresolved thresholds remain as stated above.


{\footnotesize \href{../../../../../../../continuum/motivation/constraint/mandate/338a4cb1c2444ec2aa7703fcd9f28247.md}{Canonical record: 338a4cb1c2444ec2aa7703fcd9f28247}; full qualifications and provenance.}\par
\section*{I.6.55 [Mandate] Operational Log PHI Protection}
\phantomsection
\label{navigation:I.6.55}
\addcontentsline{toc}{section}{I.6.55 [Mandate] Operational Log PHI Protection}

\label{element:ab99a48c29c14a1883e820baf98b643a}
\index{[Mandate] Operational Log PHI Protection (I.6.55)}
{\small Constraint / Mandate · Solution Architect-reviewed Iteration 001; authorised ADD; resulting increment review pending · current}\par

Exclude PHI from operational logs unless its inclusion complies with the governed Harmonia Logging Strategy.



Source entry: \textbf{NFR-076 — Operational Log PHI Protection}. Primary subject: \texttt{\seqsplit{G03.04}}.



The source classification is Constraint «Mandate». The source identifier is a scoped register reference and does not determine architectural type. Applicability, release scope and unresolved thresholds remain as stated above.


{\footnotesize \href{../../../../../../../continuum/motivation/constraint/mandate/ab99a48c29c14a1883e820baf98b643a.md}{Canonical record: ab99a48c29c14a1883e820baf98b643a}; full qualifications and provenance.}\par
\section*{I.6.56 [Mandate] Encryption at Rest}
\phantomsection
\label{navigation:I.6.56}
\addcontentsline{toc}{section}{I.6.56 [Mandate] Encryption at Rest}

\label{element:4522bf0326334607a4a6177c819e61b6}
\index{[Mandate] Encryption at Rest (I.6.56)}
{\small Constraint / Mandate · Solution Architect-reviewed Iteration 001; authorised ADD; resulting increment review pending · current}\par

Encrypt all Harmonia-managed data at rest.



Source entry: \textbf{NFR-077 — Encryption at Rest}. Primary subject: \texttt{\seqsplit{G03.04}}.



The source classification is Constraint «Mandate». The source identifier is a scoped register reference and does not determine architectural type. Applicability, release scope and unresolved thresholds remain as stated above.


{\footnotesize \href{../../../../../../../continuum/motivation/constraint/mandate/4522bf0326334607a4a6177c819e61b6.md}{Canonical record: 4522bf0326334607a4a6177c819e61b6}; full qualifications and provenance.}\par
\section*{I.6.57 [Mandate] Internal Encryption in Transit}
\phantomsection
\label{navigation:I.6.57}
\addcontentsline{toc}{section}{I.6.57 [Mandate] Internal Encryption in Transit}

\label{element:4b540b55db1b4c1197482d8f29d70e45}
\index{[Mandate] Internal Encryption in Transit (I.6.57)}
{\small Constraint / Mandate · Solution Architect-reviewed Iteration 001; authorised ADD; resulting increment review pending · current}\par

Encrypt data in transit between internal Harmonia components.



Source entry: \textbf{NFR-078 — Internal Encryption in Transit}. Primary subject: \texttt{\seqsplit{G03.04}}.



The source classification is Constraint «Mandate». The source identifier is a scoped register reference and does not determine architectural type. Applicability, release scope and unresolved thresholds remain as stated above.


{\footnotesize \href{../../../../../../../continuum/motivation/constraint/mandate/4b540b55db1b4c1197482d8f29d70e45.md}{Canonical record: 4b540b55db1b4c1197482d8f29d70e45}; full qualifications and provenance.}\par
\section*{I.6.58 [Mandate] Hard Shell Hard Body Security}
\phantomsection
\label{navigation:I.6.58}
\addcontentsline{toc}{section}{I.6.58 [Mandate] Hard Shell Hard Body Security}

\label{element:fb7783088c894106a017912467a9e40d}
\index{[Mandate] Hard Shell Hard Body Security (I.6.58)}
{\small Constraint / Mandate · Solution Architect-reviewed Iteration 001; authorised ADD; resulting increment review pending · current}\par

Apply a 'Hard Shell, Hard Body' security profile: enforce security at the perimeter and throughout internal components and interactions.



Source entry: \textbf{NFR-083 — Hard Shell Hard Body Security}. Primary subject: \texttt{\seqsplit{G03.04}}.



The source classification is Constraint «Mandate». The source identifier is a scoped register reference and does not determine architectural type. Applicability, release scope and unresolved thresholds remain as stated above.


{\footnotesize \href{../../../../../../../continuum/motivation/constraint/mandate/fb7783088c894106a017912467a9e40d.md}{Canonical record: fb7783088c894106a017912467a9e40d}; full qualifications and provenance.}\par
\section*{I.6.59 [Mandate] Harmonia FOSS Licensing}
\phantomsection
\label{navigation:I.6.59}
\addcontentsline{toc}{section}{I.6.59 [Mandate] Harmonia FOSS Licensing}

\label{element:40880a1432ab4848ab5278839620a406}
\index{[Mandate] Harmonia FOSS Licensing (I.6.59)}
{\small Constraint / Mandate · Solution Architect-reviewed Iteration 001; authorised ADD; resulting increment review pending · current}\par

Harmonia itself shall be Free and Open Source Software (FOSS).



Source entry: \textbf{NFR-100 — Harmonia FOSS Licensing}. Primary subject: \texttt{\seqsplit{G01.03}}.



The source classification is Constraint «Mandate». The source identifier is a scoped register reference and does not determine architectural type. Applicability, release scope and unresolved thresholds remain as stated above.


{\footnotesize \href{../../../../../../../continuum/motivation/constraint/mandate/40880a1432ab4848ab5278839620a406.md}{Canonical record: 40880a1432ab4848ab5278839620a406}; full qualifications and provenance.}\par
\section*{I.6.60 [Mandate] FOSS Technology Dependencies}
\phantomsection
\label{navigation:I.6.60}
\addcontentsline{toc}{section}{I.6.60 [Mandate] FOSS Technology Dependencies}

\label{element:13c56fcda53c406494561c883f32fc1c}
\index{[Mandate] FOSS Technology Dependencies (I.6.60)}
{\small Constraint / Mandate · Solution Architect-reviewed Iteration 001; authorised ADD; resulting increment review pending · current}\par

Use FOSS technologies/dependencies; proprietary or closed-source software requires explicit HCRL approval.



Source entry: \textbf{NFR-101 — FOSS Technology Dependencies}. Primary subject: \texttt{\seqsplit{G01.03}}.



The source classification is Constraint «Mandate». The source identifier is a scoped register reference and does not determine architectural type. Applicability, release scope and unresolved thresholds remain as stated above.


{\footnotesize \href{../../../../../../../continuum/motivation/constraint/mandate/13c56fcda53c406494561c883f32fc1c.md}{Canonical record: 13c56fcda53c406494561c883f32fc1c}; full qualifications and provenance.}\par
\section*{I.6.61 [Mandate] Disaster Recovery Consistency}
\phantomsection
\label{navigation:I.6.61}
\addcontentsline{toc}{section}{I.6.61 [Mandate] Disaster Recovery Consistency}

\label{element:81fdb5611a0f4338ae85478781032dd7}
\index{[Mandate] Disaster Recovery Consistency (I.6.61)}
{\small Constraint / Mandate · Solution Architect-reviewed Iteration 001; authorised ADD; resulting increment review pending · current}\par

Harmonia shall support recovery to a known, operationally consistent state following catastrophic infrastructure failure or information corruption.



Source entry: \textbf{NFR-106 — Disaster Recovery Consistency}. Primary subject: \texttt{\seqsplit{G07.04}}.



The source classification is Constraint «Mandate». The source identifier is a scoped register reference and does not determine architectural type. Applicability, release scope and unresolved thresholds remain as stated above.


{\footnotesize \href{../../../../../../../continuum/motivation/constraint/mandate/81fdb5611a0f4338ae85478781032dd7.md}{Canonical record: 81fdb5611a0f4338ae85478781032dd7}; full qualifications and provenance.}\par
\section*{I.6.62 [Mandate] Backup and Replication}
\phantomsection
\label{navigation:I.6.62}
\addcontentsline{toc}{section}{I.6.62 [Mandate] Backup and Replication}

\label{element:8731f88e202a4f5dacb408e5a1d474a8}
\index{[Mandate] Backup and Replication (I.6.62)}
{\small Constraint / Mandate · Solution Architect-reviewed Iteration 001; authorised ADD; resulting increment review pending · current}\par

Governed information and recoverable operational state shall be protected through appropriate backup and/or replication mechanisms, consistent with applicable information retention, integrity and recovery obligations.



Source entry: \textbf{NFR-108 — Backup and Replication}. Primary subject: \texttt{\seqsplit{G07.04}}.



The source classification is Constraint «Mandate». The source identifier is a scoped register reference and does not determine architectural type. Applicability, release scope and unresolved thresholds remain as stated above.


{\footnotesize \href{../../../../../../../continuum/motivation/constraint/mandate/8731f88e202a4f5dacb408e5a1d474a8.md}{Canonical record: 8731f88e202a4f5dacb408e5a1d474a8}; full qualifications and provenance.}\par
\section*{I.6.63 [Mandate] Controlled Capacity Admission}
\phantomsection
\label{navigation:I.6.63}
\addcontentsline{toc}{section}{I.6.63 [Mandate] Controlled Capacity Admission}

\label{element:fb30eb9b8a3c41eda3a8b5edfee95ac1}
\index{[Mandate] Controlled Capacity Admission (I.6.63)}
{\small Constraint / Mandate · Solution Architect-reviewed Iteration 001; authorised ADD; resulting increment review pending · current}\par

Harmonia shall manage increasing workloads within its operational capacity and, where defined capacity thresholds are reached, shall apply controlled admission or rejection of incoming activity and report the associated capacity or performance conditions.



Source entry: \textbf{NFR-111 — Controlled Capacity Admission}. Primary subject: \texttt{\seqsplit{G08.04}}.



The source classification is Constraint «Mandate». The source identifier is a scoped register reference and does not determine architectural type. Applicability, release scope and unresolved thresholds remain as stated above.


{\footnotesize \href{../../../../../../../continuum/motivation/constraint/mandate/fb30eb9b8a3c41eda3a8b5edfee95ac1.md}{Canonical record: fb30eb9b8a3c41eda3a8b5edfee95ac1}; full qualifications and provenance.}\par
\section*{I.6.64 [Mandate] Information Integrity Assessment}
\phantomsection
\label{navigation:I.6.64}
\addcontentsline{toc}{section}{I.6.64 [Mandate] Information Integrity Assessment}

\label{element:6316761bfd414032a8d889c7ba34131d}
\index{[Mandate] Information Integrity Assessment (I.6.64)}
{\small Constraint / Mandate · Solution Architect-reviewed Iteration 001; authorised ADD; resulting increment review pending · current}\par

Harmonia's governed information framework shall support the assessment and reporting of the integrity of individual governed information items.



Source entry: \textbf{NFR-119 — Information Integrity Assessment}. Primary subject: \texttt{\seqsplit{G02.02}}.



The source classification is Constraint «Mandate». The source identifier is a scoped register reference and does not determine architectural type. Applicability, release scope and unresolved thresholds remain as stated above.


{\footnotesize \href{../../../../../../../continuum/motivation/constraint/mandate/6316761bfd414032a8d889c7ba34131d.md}{Canonical record: 6316761bfd414032a8d889c7ba34131d}; full qualifications and provenance.}\par
\section*{I.6.65 [Mandate] Migration Information Integrity}
\phantomsection
\label{navigation:I.6.65}
\addcontentsline{toc}{section}{I.6.65 [Mandate] Migration Information Integrity}

\label{element:22583ccb89dc4cf683c710760bc4b1f2}
\index{[Mandate] Migration Information Integrity (I.6.65)}
{\small Constraint / Mandate · Solution Architect-reviewed Iteration 001; authorised ADD; resulting increment review pending · current}\par

Any migration of governed information shall preserve information integrity, including the identities of information items and their associated references, across the migration boundary.



Source entry: \textbf{NFR-121 — Migration Information Integrity}. Primary subject: \texttt{\seqsplit{G02.02}}.



The source classification is Constraint «Mandate». The source identifier is a scoped register reference and does not determine architectural type. Applicability, release scope and unresolved thresholds remain as stated above.


{\footnotesize \href{../../../../../../../continuum/motivation/constraint/mandate/22583ccb89dc4cf683c710760bc4b1f2.md}{Canonical record: 22583ccb89dc4cf683c710760bc4b1f2}; full qualifications and provenance.}\par
\section*{I.6.66 [Mandate] Secrets Protection}
\phantomsection
\label{navigation:I.6.66}
\addcontentsline{toc}{section}{I.6.66 [Mandate] Secrets Protection}

\label{element:4547bd339ff143c988391159e3407f20}
\index{[Mandate] Secrets Protection (I.6.66)}
{\small Constraint / Mandate · Solution Architect-reviewed Iteration 001; authorised ADD; resulting increment review pending · current}\par

Harmonia shall protect security-sensitive configuration information, credentials, cryptographic keys and other secrets against unauthorised access, disclosure and modification.



Source entry: \textbf{NFR-125 — Secrets Protection}. Primary subject: \texttt{\seqsplit{G03.04}}.



The source classification is Constraint «Mandate». The source identifier is a scoped register reference and does not determine architectural type. Applicability, release scope and unresolved thresholds remain as stated above.


{\footnotesize \href{../../../../../../../continuum/motivation/constraint/mandate/4547bd339ff143c988391159e3407f20.md}{Canonical record: 4547bd339ff143c988391159e3407f20}; full qualifications and provenance.}\par
\section*{I.6.67 [Mandate] Governed DevSecOps Deployment}
\phantomsection
\label{navigation:I.6.67}
\addcontentsline{toc}{section}{I.6.67 [Mandate] Governed DevSecOps Deployment}

\label{element:54eb7e671c094ebcbb5503af36c7892d}
\index{[Mandate] Governed DevSecOps Deployment (I.6.67)}
{\small Constraint / Mandate · Solution Architect-reviewed Iteration 001; authorised ADD; resulting increment review pending · current}\par

Harmonia shall adopt a defined DevSecOps and software-defined deployment model, with software builds, deployment definitions, configuration changes and release activities governed, version-controlled and traceable through an established change-management framework.



Source entry: \textbf{NFR-129 — Governed DevSecOps Deployment}. Primary subject: \texttt{\seqsplit{G10.02}}.



The source classification is Constraint «Mandate». The source identifier is a scoped register reference and does not determine architectural type. Applicability, release scope and unresolved thresholds remain as stated above.


{\footnotesize \href{../../../../../../../continuum/motivation/constraint/mandate/54eb7e671c094ebcbb5503af36c7892d.md}{Canonical record: 54eb7e671c094ebcbb5503af36c7892d}; full qualifications and provenance.}\par
\section*{I.6.68 [Mandate] Controlled Component Shutdown and Restart}
\phantomsection
\label{navigation:I.6.68}
\addcontentsline{toc}{section}{I.6.68 [Mandate] Controlled Component Shutdown and Restart}

\label{element:96f7bfa320b049d2b76fcc35a1926883}
\index{[Mandate] Controlled Component Shutdown and Restart (I.6.68)}
{\small Constraint / Mandate · Solution Architect-reviewed Iteration 001; authorised ADD; resulting increment review pending · current}\par

Harmonia shall support controlled shutdown and restart of its operational components, modules and subsystems, preserving accepted governed work and recoverable processing state throughout the transition.



Source entry: \textbf{NFR-146 — Controlled Component Shutdown and Restart}. Primary subject: \texttt{\seqsplit{G09.03}}.



The source classification is Constraint «Mandate». The source identifier is a scoped register reference and does not determine architectural type. Applicability, release scope and unresolved thresholds remain as stated above.


{\footnotesize \href{../../../../../../../continuum/motivation/constraint/mandate/96f7bfa320b049d2b76fcc35a1926883.md}{Canonical record: 96f7bfa320b049d2b76fcc35a1926883}; full qualifications and provenance.}\par
\section*{I.6.69 [Mandate] Governed Reference Test Harness}
\phantomsection
\label{navigation:I.6.69}
\addcontentsline{toc}{section}{I.6.69 [Mandate] Governed Reference Test Harness}

\label{element:5ce1ef8baaf84848979ff6c9ccf6ca91}
\index{[Mandate] Governed Reference Test Harness (I.6.69)}
{\small Constraint / Mandate · Solution Architect-reviewed Iteration 001; authorised ADD; resulting increment review pending · current}\par

Harmonia shall maintain a governed reference test harness that is reviewed, updated and executed as part of each solution release, providing regression verification of the solution's functional and non-functional behaviour against its applicable requirements and architectural contracts.



Source entry: \textbf{NFR-148 — Governed Reference Test Harness}. Primary subject: \texttt{\seqsplit{G11.06}}.



The source classification is Constraint «Mandate». The source identifier is a scoped register reference and does not determine architectural type. Applicability, release scope and unresolved thresholds remain as stated above.


{\footnotesize \href{../../../../../../../continuum/motivation/constraint/mandate/5ce1ef8baaf84848979ff6c9ccf6ca91.md}{Canonical record: 5ce1ef8baaf84848979ff6c9ccf6ca91}; full qualifications and provenance.}\par
\section*{I.6.70 [Mandate] Controlled Resilience Testing}
\phantomsection
\label{navigation:I.6.70}
\addcontentsline{toc}{section}{I.6.70 [Mandate] Controlled Resilience Testing}

\label{element:f2d38e7beef84b61a769de4ca7f48951}
\index{[Mandate] Controlled Resilience Testing (I.6.70)}
{\small Constraint / Mandate · Solution Architect-reviewed Iteration 001; authorised ADD; resulting increment review pending · current}\par

Harmonia's governed reference test harness shall support the controlled simulation of representative infrastructure and component failures to verify solution resilience, failure handling and recovery behaviour.



Source entry: \textbf{NFR-150 — Controlled Resilience Testing}. Primary subject: \texttt{\seqsplit{G11.05}}.



The source classification is Constraint «Mandate». The source identifier is a scoped register reference and does not determine architectural type. Applicability, release scope and unresolved thresholds remain as stated above.


{\footnotesize \href{../../../../../../../continuum/motivation/constraint/mandate/f2d38e7beef84b61a769de4ca7f48951.md}{Canonical record: f2d38e7beef84b61a769de4ca7f48951}; full qualifications and provenance.}\par
\section*{I.6.71 [Mandate] Exclusive PetasosParcel Routing}
\phantomsection
\label{navigation:I.6.71}
\addcontentsline{toc}{section}{I.6.71 [Mandate] Exclusive PetasosParcel Routing}

\label{element:e024703b32c5437481ce290e8e50b079}
\index{[Mandate] Exclusive PetasosParcel Routing (I.6.71)}
{\small Constraint / Mandate · Solution Architect explicitly approved inclusion; incorporated-result/visual review pending · current}\par

\par\medskip\textbf{Approved definition}\par


Harmonia SHALL route all PetasosParcel instances exclusively through Apache ActiveMQ Artemis, irrespective of deployment location or runtime proximity of participating components. Petasos SHALL utilise Artemis native messaging capabilities for routing, concurrency, availability and delivery management.



The Solution Architect's complete Subject/Views Task explicitly approves this
native Constraint / Mandate for the governed Continuum.
The incoming candidate/pending-ADD caption records its earlier source standing;
it does not override this subsequent inclusion approval. Original source wording,
native type, specialisation and Exchange identity evidence remain recoverable in
the retained immutable 004 export. This incorporation supplies neither HCRL closure
nor actual Archi visual acceptance.



This element retains its own bounded meaning. Broader existing security,
authorisation, messaging or stakeholder concerns do not establish identity
equivalence. No unrelated incoming properties, relationships, classifications,
Capability allocations or OpenSpec candidates are incorporated.


{\footnotesize \href{../../../../../../../continuum/motivation/constraint/mandate/e024703b32c5437481ce290e8e50b079.md}{Canonical record: e024703b32c5437481ce290e8e50b079}; full qualifications and provenance.}\par
\chapter*{I.7 Architecture Axioms}
\phantomsection
\label{navigation:I.7}
\addcontentsline{toc}{chapter}{I.7 Architecture Axioms}

\section*{I.7.1 AX-01 Harmonia Is Health-Information Centric}
\phantomsection
\label{navigation:I.7.1}
\addcontentsline{toc}{section}{I.7.1 AX-01 Harmonia Is Health-Information Centric}

\label{element:e9bd61acaaea4577bc4086a5639e52c3}
\index{AX-01 Harmonia Is Health-Information Centric (I.7.1)}
{\small Principle · Canonical · current}\par

\par\medskip\textbf{Axiom}\par


Harmonia is a distributed health-information management and
interoperability framework.



Its primary purpose is the ingestion, validation, transformation,
governance, processing, persistence, distribution and exposure of
health-related information.



Harmonia SHALL support applicable health-information standards and
models while retaining the ability to manage information according to
Harmonia's own internal operational requirements.



\par\medskip\textbf{Rationale}\par


Harmonia is not intended to become a generic distributed application
framework to which healthcare happens to be attached.



Its architecture exists to solve health-information problems:
interoperability, information authority, provenance, security,
terminology, workflow, persistence, distribution, resilience and
longitudinal information management.



\par\medskip\textbf{Consequences}\par


Health-information semantics take precedence over generic infrastructure
convenience.



Standards such as HL7 FHIR and HL7 v2 are first-class architectural
concerns.



Generic infrastructure capabilities SHOULD be used where useful, but
SHALL NOT dictate Harmonia's health-information semantics.



\textbf{This does not mean}



\begin{itemize}


\item Harmonia is a FHIR server.

\item Harmonia is restricted to FHIR information.

\item Every internal Harmonia object must correspond to a healthcare interoperability standard.

\end{itemize}

{\footnotesize \href{../../../../../../../continuum/motivation/principle/e9bd61acaaea4577bc4086a5639e52c3.md}{Canonical record: e9bd61acaaea4577bc4086a5639e52c3}; full qualifications and provenance.}\par
\section*{I.7.2 AX-02 Standards at the Boundary; Harmonia Within the Boundary}
\phantomsection
\label{navigation:I.7.2}
\addcontentsline{toc}{section}{I.7.2 AX-02 Standards at the Boundary; Harmonia Within the Boundary}

\label{element:eee8c3783d964d18a7f9dcea85fcd4cf}
\index{AX-02 Standards at the Boundary; Harmonia Within the Boundary (I.7.2)}
{\small Principle · Canonical · current}\par

\par\medskip\textbf{Axiom}\par


Harmonia SHALL ingest and expose health information using the standards,
information models and interoperability protocols applicable to an
external interface contract.



Within Harmonia, information MAY be augmented, represented, indexed,
cached, distributed, governed, versioned, correlated and persisted using
Harmonia-specific semantics and metadata necessary to satisfy Harmonia's
operational requirements.



Harmonia's internal operational semantics are private to Harmonia.
External systems SHALL NOT be required to understand, preserve,
reproduce or participate in them.



\par\medskip\textbf{Rationale}\par


Interoperability requires agreement at system boundaries. It does not
require participating systems to share an internal architecture.



Allowing external contracts to dictate Harmonia's internal
implementation would couple Harmonia to the assumptions and limitations
of external systems.



\par\medskip\textbf{Consequences}\par


Pylai acts as an interoperability membrane, not merely a protocol
adapter.



External representations SHALL conform to their applicable interface
contracts.



Harmonia-private operational metadata SHALL NOT leak into an external
interoperability contract.



\par\medskip\textbf{Maxim}\par


Standards at the boundary. Harmonia within the boundary.


{\footnotesize \href{../../../../../../../continuum/motivation/principle/eee8c3783d964d18a7f9dcea85fcd4cf.md}{Canonical record: eee8c3783d964d18a7f9dcea85fcd4cf}; full qualifications and provenance.}\par
\section*{I.7.3 AX-03 Native Standards Models Remain Native}
\phantomsection
\label{navigation:I.7.3}
\addcontentsline{toc}{section}{I.7.3 AX-03 Native Standards Models Remain Native}

\label{element:e99a4e2bf99f476d876394300d6b9c3d}
\index{AX-03 Native Standards Models Remain Native (I.7.3)}
{\small Principle · Canonical · current}\par

\par\medskip\textbf{Axiom}\par


Standards-defined information objects SHALL retain their
standards-defined representation within Harmonia wherever practicable.



Harmonia SHALL NOT introduce parallel or derivative representations
solely to accommodate internal management concerns where the standard
representation provides an appropriate extensibility mechanism.



\par\medskip\textbf{Rationale}\par


Parallel representations introduce mapping complexity, semantic drift
and unnecessary ownership questions.



For FHIR information, the FHIR resource is already an extensible
information model.



\par\medskip\textbf{Consequences}\par


FHIR R5 and FHIR R4 AU resources SHOULD remain native FHIR resources
represented using the applicable HAPI FHIR models.



Where appropriate, Harmonia-specific durable metadata MAY be represented
using private FHIR extensions.



Jurisdictional and implementation-guide extensions, such as extensions
required for Australian FHIR conformance, remain part of the
standards-based representation and SHALL NOT be confused with
Harmonia-private operational extensions.



Non-FHIR managed objects MAY use representation-appropriate Harmonia
metadata mechanisms.



\textbf{This does not mean:}



\begin{itemize}


\item All Harmonia metadata must be represented as FHIR extensions.

\item Cache-specific state must be embedded into FHIR resources.

\item Non-FHIR Harmonia objects must be converted into FHIR.

\end{itemize}

{\footnotesize \href{../../../../../../../continuum/motivation/principle/e99a4e2bf99f476d876394300d6b9c3d.md}{Canonical record: e99a4e2bf99f476d876394300d6b9c3d}; full qualifications and provenance.}\par
\section*{I.7.4 AX-04 Harmonia Owns the Semantics; Engines Provide the Machinery}
\phantomsection
\label{navigation:I.7.4}
\addcontentsline{toc}{section}{I.7.4 AX-04 Harmonia Owns the Semantics; Engines Provide the Machinery}

\label{element:54508cda80484a6caf95b901780eaf7d}
\index{AX-04 Harmonia Owns the Semantics; Engines Provide the Machinery (I.7.4)}
{\small Principle · Canonical · current}\par

\par\medskip\textbf{Axiom}\par


Harmonia owns the architectural semantics governing information
authority, security, lifecycle, concurrency, provenance, auditability,
resilience and operational integrity.



Harmonia SHALL preferentially use capabilities supplied by underlying
technology engines where those capabilities satisfy Harmonia's
architectural invariants rather than reproduce equivalent functionality.



\par\medskip\textbf{Rationale}\par


Harmonia should own its meaning without unnecessarily owning commodity
infrastructure machinery.



\par\medskip\textbf{Consequences}\par


HAPI FHIR MAY provide FHIR persistence, indexing, versioning and search.



Infinispan MAY provide distributed caching, entry metadata and atomic
coordination primitives.



ActiveMQ Artemis MAY provide durable messaging.



PostgreSQL MAY provide transactional relational persistence.



These technologies implement Harmonia capabilities. They do not define
Harmonia's architectural semantics.



\par\medskip\textbf{Maxim}\par


Use the machinery. Own the semantics.



\textbf{This does not mean:}



\begin{itemize}


\item Harmonia should wrap every engine capability behind a bespoke implementation.

\item An engine's native behaviour automatically satisfies a Harmonia invariant.

\item Using HAPI FHIR makes Harmonia a HAPI FHIR application.

\end{itemize}

{\footnotesize \href{../../../../../../../continuum/motivation/principle/54508cda80484a6caf95b901780eaf7d.md}{Canonical record: 54508cda80484a6caf95b901780eaf7d}; full qualifications and provenance.}\par
\section*{I.7.5 AX-05 Active State and Authoritative Durable State Are Distinct}
\phantomsection
\label{navigation:I.7.5}
\addcontentsline{toc}{section}{I.7.5 AX-05 Active State and Authoritative Durable State Are Distinct}

\label{element:0412a2084f0145babc4f022c71fea65f}
\index{AX-05 Active State and Authoritative Durable State Are Distinct (I.7.5)}
{\small Principle · Canonical · current}\par

\par\medskip\textbf{Axiom}\par


Mneme owns Harmonia's application-facing access to managed information
and the distributed active-state representation, observation and
coordination required to use that information safely.



Mnemosyne owns Harmonia's authoritative durable representation of
managed information. It atomically establishes authoritative state and
authoritative version progression and persists the durable management
metadata required to interpret that state.



Mneme manages active use; Mnemosyne establishes durable truth.



\par\medskip\textbf{Rationale}\par


Distributed availability and authoritative persistence are different
responsibilities.



Conflating them makes cache state authoritative, persistence state
application-facing, or both.



\par\medskip\textbf{Consequences}\par


Mneme MAY reject or coordinate a proposed state progression before
persistence.



Only Mnemosyne can establish a new authoritative durable state.



Following authoritative commit, Mneme SHALL converge its active
representation toward the authoritative state.



Loss of Mneme SHALL NOT cause process-local or cached state to become
authoritative.



Loss of Mnemosyne SHALL NOT cause Mneme to promote its active
representation into authoritative durable state.



Mneme active state SHALL be reconstructable from authoritative state
where applicable.



\textbf{This does not mean:}



\begin{itemize}


\item Every Mneme operation requires a database operation.

\item Mnemosyne is exposed as an application database API.

\item Mneme and Mnemosyne must use the same version or concurrency mechanism.

\end{itemize}


Mneme active-state generation SHALL represent successfully established active state. A valid active-state token SHALL identify an observed generation of that state.



Failed or degraded convergence following authoritative state progression SHALL NOT establish or advance a valid active-state generation. Where the active representation cannot be successfully converged, its coordination state SHALL be treated as untrusted until reconciled with authoritative state.



Mneme active-state generation and Mnemosyne authoritative version SHALL remain distinct concurrency domains. Neither SHALL be inferred from, substituted for, or treated as an alias of the other.


{\footnotesize \href{../../../../../../../continuum/motivation/principle/0412a2084f0145babc4f022c71fea65f.md}{Canonical record: 0412a2084f0145babc4f022c71fea65f}; full qualifications and provenance.}\par
\section*{I.7.6 AX-06 Information Authority Is Explicit}
\phantomsection
\label{navigation:I.7.6}
\addcontentsline{toc}{section}{I.7.6 AX-06 Information Authority Is Explicit}

\label{element:952c3b6148ce4fce94b69f99a0c2fec7}
\index{AX-06 Information Authority Is Explicit (I.7.6)}
{\small Principle · Canonical · current}\par

\par\medskip\textbf{Axiom}\par


Harmonia SHALL represent the authority and credibility of managed
information independently of its technical transport, cache state,
persistence state or concurrency state.



Harmonia SHALL support governed information-authority concepts including:



\begin{itemize}


\item Authoritative,

\item Informational, and

\item Anecdotal.

\end{itemize}


\par\medskip\textbf{Rationale}\par


Technical correctness and information credibility answer different
questions.



A successfully persisted assertion is not necessarily authoritative. A
stale assertion is not necessarily false. A technically reliable source
is not necessarily the authoritative source for every fact it supplies.



\par\medskip\textbf{Consequences}\par


Concurrency determines whether information has changed.



Semantic analysis determines what changed.



Information authority and assurance describe the standing of competing
assertions.



Policy determines what Harmonia does about them.



Calliope governs the semantics of information authority. Appropriate
Harmonia processes and policies apply those semantics.



\textbf{This does not mean:}



\begin{itemize}


\item Authoritative information automatically wins every conflict.

\item Anecdotal information is discarded.

\item Persistence success establishes information authority.

\end{itemize}

{\footnotesize \href{../../../../../../../continuum/motivation/principle/952c3b6148ce4fce94b69f99a0c2fec7.md}{Canonical record: 952c3b6148ce4fce94b69f99a0c2fec7}; full qualifications and provenance.}\par
\section*{I.7.7 AX-07 Security Is Intrinsic to Managed Operations}
\phantomsection
\label{navigation:I.7.7}
\addcontentsline{toc}{section}{I.7.7 AX-07 Security Is Intrinsic to Managed Operations}

\label{element:5de2ae2c3b794c2d9f3a14efd51d564f}
\index{AX-07 Security Is Intrinsic to Managed Operations (I.7.7)}
{\small Principle · Canonical · current}\par

\par\medskip\textbf{Axiom}\par


Every governed operation SHALL execute within an established security
context and SHALL be subject to Harmonia security policy.



Security SHALL be enforced by framework and platform boundaries where
practicable and SHALL NOT depend solely upon developer knowledge, coding
convention or voluntary caller behaviour.



\par\medskip\textbf{Rationale}\par


Security is part of the semantics of performing an operation, not an
optional service invoked by well-behaved code.



\par\medskip\textbf{Consequences}\par


Themis owns Harmonia security policy evaluation.



Trusted identity and security context SHALL be established before
governed actions are performed.



Security context is operational context and SHALL NOT automatically
become persistent information content.



Security-significant events MAY become durable audit evidence according
to policy.







\par\medskip\textbf{Kleio mandatory authoritative audit capture}\par


Kleio SHALL perform mandatory authoritative audit capture under intrinsic
platform authority associated with its established audit responsibility.
That capture SHALL execute within an established security context and remain
subject to Harmonia security policy; Themis retains policy evaluation ownership.
This is governed platform authority, not a security bypass or policy exception.



The Identity associated with the audited activity SHALL remain audit evidence.
It is not the authority under which the mandatory AuditEvent is persisted.
Kleio SHALL NOT borrow the originating Identity's PEP-Criteria merely to persist
evidence of that Identity's activity. Identity-derived authorisation resolved
and enforced through Themis remains distinct from Kleio's capture authority.



Where authoritative audit capture is required for a governed operation and
Kleio cannot exercise that platform authority sufficiently to capture the
required evidence, the governed operation SHALL fail closed.



This authority is established specifically for Kleio mandatory authoritative
audit capture. It SHALL NOT be inferred for other Kleio activities or generalised
to other Harmonia components. Any proposed equivalent authority elsewhere is
a separate architecture question under AX-17. Detailed authority representation
and enforcement mechanisms remain deferred.



\textbf{This does not mean:}



\begin{itemize}


\item Every security decision belongs in a FHIR resource.

\item Every authorization evaluation belongs in the audit repository.

\item Internal security context should be exposed to external systems.

\end{itemize}

{\footnotesize \href{../../../../../../../continuum/motivation/principle/5de2ae2c3b794c2d9f3a14efd51d564f.md}{Canonical record: 5de2ae2c3b794c2d9f3a14efd51d564f}; full qualifications and provenance.}\par
\section*{I.7.8 AX-08 Evidence Records Meaning, Not Machinery}
\phantomsection
\label{navigation:I.7.8}
\addcontentsline{toc}{section}{I.7.8 AX-08 Evidence Records Meaning, Not Machinery}

\label{element:21d6864657964ba5b03e3f6b5eb29cd5}
\index{AX-08 Evidence Records Meaning, Not Machinery (I.7.8)}
{\small Principle · Canonical · current}\par

\par\medskip\textbf{Axiom}\par


Provenance and audit evidence SHALL describe information-significant,
security-significant and business-significant events, assertions and
decisions.



Internal implementation mechanics SHALL NOT ordinarily become provenance
or audit evidence merely because they occurred.



\par\medskip\textbf{Rationale}\par


A durable evidence repository should explain what happened to
information and why, rather than become a permanent distributed debug
log.



\par\medskip\textbf{Consequences}\par


Authoritative state acceptance, significant information transformations,
stewardship decisions and security-significant actions MAY require
durable evidence.



Cache misses, cache rebalancing, CAS retries, connection retries,
reconstruction and ordinary infrastructure mechanics normally belong in
logging, metrics or transient diagnostics.



Kleio preserves durable evidence.



\par\medskip\textbf{Maxim}\par


Preserve the meaningful fact, not every mechanism that produced it.


{\footnotesize \href{../../../../../../../continuum/motivation/principle/21d6864657964ba5b03e3f6b5eb29cd5.md}{Canonical record: 21d6864657964ba5b03e3f6b5eb29cd5}; full qualifications and provenance.}\par
\section*{I.7.9 AX-09 Transient Operational State Is Ephemeral by Default}
\phantomsection
\label{navigation:I.7.9}
\addcontentsline{toc}{section}{I.7.9 AX-09 Transient Operational State Is Ephemeral by Default}

\label{element:9a5cb9e8876844e58e660f909f7df9fc}
\index{AX-09 Transient Operational State Is Ephemeral by Default (I.7.9)}
{\small Principle · Canonical · current}\par

\par\medskip\textbf{Axiom}\par


Harmonia MAY maintain transient operational state associated with
managed information where required for caching, persistence,
concurrency, security, routing, resilience, diagnostics and processing
integrity.



Such state SHALL be retained only for as long as required by its
operational purpose and SHALL normally disappear when that purpose has
been satisfied.



Transient operational state SHALL NOT automatically become provenance or
audit evidence.



\par\medskip\textbf{Diagnostic observability}\par


Harmonia SHALL provide controlled mechanisms by which designated
transient operational state MAY be observed or reported when diagnostics
are explicitly enabled at framework, subsystem or component scope.



Where transient state directly contributes to an information-significant
decision, assertion or evidentiary fact, the meaningful information
SHALL be promoted into the appropriate durable provenance or audit
representation.



\par\medskip\textbf{Rationale}\par


Distributed systems require substantial operational state. Persisting
all of it indefinitely would confuse implementation mechanics with
information history.



Conversely, making operational state completely invisible would make the
framework extremely difficult to diagnose.



\textbf{This does not mean:}



\begin{itemize}


\item Transient state is unimportant.

\item Transient state cannot be logged or inspected.

\item Diagnostic traces should automatically become permanent evidence.

\end{itemize}

{\footnotesize \href{../../../../../../../continuum/motivation/principle/9a5cb9e8876844e58e660f909f7df9fc.md}{Canonical record: 9a5cb9e8876844e58e660f909f7df9fc}; full qualifications and provenance.}\par
\section*{I.7.10 AX-10 Distribution, Load and Failure Are Normal Operating Conditions}
\phantomsection
\label{navigation:I.7.10}
\addcontentsline{toc}{section}{I.7.10 AX-10 Distribution, Load and Failure Are Normal Operating Conditions}

\label{element:0d9673833eb840eab03804aa076c1b0c}
\index{AX-10 Distribution, Load and Failure Are Normal Operating Conditions (I.7.10)}
{\small Principle · Canonical · current}\par

\par\medskip\textbf{Axiom}\par


Harmonia SHALL be designed on the assumption of distributed deployment,
sustained processing load and failure or degradation of individual
runtime components.



Loss or degradation of implementation machinery SHALL NOT silently alter
the authoritative, security, governance or information-authority
semantics of managed information.



\par\medskip\textbf{Rationale}\par


Distribution and failure are not exceptional architectural cases for
Harmonia. They are part of its normal operating environment.



\par\medskip\textbf{Consequences}\par


Work that has crossed an explicit durable processing boundary SHALL be
recoverable according to that boundary's contract.



Processing pressure SHOULD preferentially be represented as bounded,
durable backlog rather than unbounded growth of JVM heap, threads,
database connections or volatile work.



Failure SHALL be visible rather than silently replaced by semantically
weaker local behaviour.



\textbf{This does not mean:}



\begin{itemize}


\item Every component must be independently distributed.

\item Every operation must be asynchronous.

\item Every internal Java boundary should become a network boundary.

\end{itemize}

{\footnotesize \href{../../../../../../../continuum/motivation/principle/0d9673833eb840eab03804aa076c1b0c.md}{Canonical record: 0d9673833eb840eab03804aa076c1b0c}; full qualifications and provenance.}\par
\section*{I.7.11 AX-11 Managed Information Access Is Highly Available and Responsive}
\phantomsection
\label{navigation:I.7.11}
\addcontentsline{toc}{section}{I.7.11 AX-11 Managed Information Access Is Highly Available and Responsive}

\label{element:8a1e995f16954484a78ae1c5c114d546}
\index{AX-11 Managed Information Access Is Highly Available and Responsive (I.7.11)}
{\small Principle · Canonical · current}\par

\par\medskip\textbf{Axiom}\par


Mneme SHALL provide highly available, responsive and load-tolerant
access to Harmonia-managed active information.



Its architecture SHALL favour bounded resource consumption, distributed
workload, reconstructable active state and graceful expression of
processing pressure.



Availability SHALL NOT be achieved by weakening authoritative-state,
concurrency, security, governance or information-authority guarantees.



\par\medskip\textbf{Rationale}\par


Mneme exists partly to prevent every application interaction from
becoming an authoritative persistence interaction while still preserving
Harmonia's correctness guarantees.



\par\medskip\textbf{Consequences}\par


Different categories of Mneme state MAY use different distribution and
consistency strategies.



Small coordination state need not have the same topology as large
resource caches.



Cache topology SHALL be selected according to the semantics and
operational requirements of the managed state rather than applied
universally.



\textbf{This does not mean:}



\begin{itemize}


\item Every cache should use REPL\_SYNC.

\item Every managed resource must always reside in memory.

\item Cache availability is equivalent to authoritative information availability.

\end{itemize}

{\footnotesize \href{../../../../../../../continuum/motivation/principle/8a1e995f16954484a78ae1c5c114d546.md}{Canonical record: 8a1e995f16954484a78ae1c5c114d546}; full qualifications and provenance.}\par
\section*{I.7.12 AX-12 Hide Plumbing, Not Information}
\phantomsection
\label{navigation:I.7.12}
\addcontentsline{toc}{section}{I.7.12 AX-12 Hide Plumbing, Not Information}

\label{element:ce85889bd9b34f79a53d4c907ceb19d4}
\index{AX-12 Hide Plumbing, Not Information (I.7.12)}
{\small Principle · Canonical · current}\par

\par\medskip\textbf{Axiom}\par


Developers implementing Harmonia work units SHALL have direct and fluent
access to applicable standards-defined health-information models and
Harmonia's managed-information capabilities.



Framework plumbing SHALL provide security, governance, authority,
persistence, concurrency, provenance, messaging and operational services
without unnecessarily obscuring or replacing the underlying information
models.



\par\medskip\textbf{Rationale}\par


Ergon and other Harmonia developers should
implement health-information behaviour, not repeatedly reconstruct the
framework's infrastructure contracts.



At the same time, hiding the actual information model behind generic
wrappers makes health semantics harder rather than easier to understand.



\par\medskip\textbf{Consequences}\par


A developer working with a FHIR Practitioner, Observation,
DiagnosticReport or other resource SHOULD be able to work with the real
standards model.



Harmonia SHALL make the safe and governed path the normal and convenient
path.



Raw infrastructure mechanisms SHOULD NOT become the ordinary application
programming model.



\par\medskip\textbf{Maxim}\par


Hide plumbing, not information.



\textbf{This does not mean:}



\begin{itemize}


\item Every capability must be exposed through one universal API.

\item Standards models should contain Harmonia infrastructure logic.

\item Developers should directly manipulate RemoteCache, EntityManager, JMS or other infrastructure merely because Harmonia uses those technologies internally.

\end{itemize}

{\footnotesize \href{../../../../../../../continuum/motivation/principle/ce85889bd9b34f79a53d4c907ceb19d4.md}{Canonical record: ce85889bd9b34f79a53d4c907ceb19d4}; full qualifications and provenance.}\par
\section*{I.7.13 AX-13 Harmonia Management Has an Explicit Boundary}
\phantomsection
\label{navigation:I.7.13}
\addcontentsline{toc}{section}{I.7.13 AX-13 Harmonia Management Has an Explicit Boundary}

\label{element:b6a55decfeec4f1bb3d7f18534a2ffea}
\index{AX-13 Harmonia Management Has an Explicit Boundary (I.7.13)}
{\small Principle · Canonical · current}\par

\par\medskip\textbf{Axiom}\par


Ingress establishes Harmonia management. Internal processing preserves
Harmonia management. Egress terminates Harmonia management of the
emitted representation.



Harmonia MAY retain evidence describing the information from which an
external representation was derived, the transformation applied, the
transmission performed and the externally observable outcome.



Harmonia SHALL NOT attribute its internal authority, governance,
concurrency, security, availability or operational guarantees to an
emitted representation after that representation crosses an external
egress boundary.



\par\medskip\textbf{Rationale}\par


Harmonia can govern information while it is under Harmonia's control. It
cannot govern a copy after that copy has been transferred to another
system.



Provenance can describe what Harmonia emitted; it cannot extend
Harmonia's operational control into another system.



\par\medskip\textbf{Consequences}\par


An HL7 v2 message ceases to be a Harmonia-managed representation when
emitted through its external egress boundary.



A FHIR resource returned through REST similarly becomes an externally
held representation after transmission.



A subsequent FHIR PUT, HL7 v2 message or other submission represents a
new ingress interaction and is assessed accordingly.



An HTTP ETag, FHIR meta.versionId, HL7 acknowledgement or equivalent
interoperability mechanism MAY communicate externally meaningful
transaction or version semantics. It SHALL NOT be interpreted as
extending Harmonia's internal operational control beyond its boundary.



\par\medskip\textbf{Maxim}\par


Provenance may describe or cross the boundary; operational control does
not.



\textbf{This does not mean:}



\begin{itemize}


\item The underlying information inside Harmonia ceases to be managed when a representation is published.

\item Harmonia forgets what it transmitted.

\item Acknowledgements and transaction outcomes are irrelevant to audit or provenance.

\end{itemize}

{\footnotesize \href{../../../../../../../continuum/motivation/principle/b6a55decfeec4f1bb3d7f18534a2ffea.md}{Canonical record: b6a55decfeec4f1bb3d7f18534a2ffea}; full qualifications and provenance.}\par
\section*{I.7.14 AX-14 Semantic Distinctions Are Preserved}
\phantomsection
\label{navigation:I.7.14}
\addcontentsline{toc}{section}{I.7.14 AX-14 Semantic Distinctions Are Preserved}

\label{element:d3fc6025dc7646d1b2020b51f7a44bfa}
\index{AX-14 Semantic Distinctions Are Preserved (I.7.14)}
{\small Principle · Canonical · current}\par

\par\medskip\textbf{Axiom}\par


Harmonia SHALL preserve meaningful distinctions between states, outcomes and concepts throughout its internal processing and across subsystem boundaries.



A distinction that is significant to information meaning, authority, security, concurrency, persistence, processing outcome or operational correctness SHALL NOT be collapsed merely because an underlying technology, transport, API or implementation abstraction does not represent that distinction directly.



Where an external interface contract intentionally presents a simpler or different semantic model, Harmonia MAY project its internal semantics into that contract explicitly at the applicable boundary.



\par\medskip\textbf{Rationale}\par


Loss of semantic distinction causes information about what Harmonia knows, what occurred and what may safely happen next to be lost.



Two conditions may produce similar implementation behaviour while having materially different meanings. Resource absence, access denial, persistence failure and indeterminate outcome may all prevent information from being returned, but they do not describe the same state.



Likewise, persistence version, active-state generation, interoperability version and information authority describe different properties even where an implementation happens to correlate them.



Implementation convenience SHALL NOT redefine those meanings.



\par\medskip\textbf{Consequences}\par


Semantically distinct outcomes SHALL remain distinguishable through Harmonia's internal contracts.



Security denial SHALL NOT be represented as resource absence merely because both prevent access to information.



A failure to establish authoritative state SHALL NOT be represented as authoritative absence.



An indeterminate outcome SHALL NOT be represented as either success or failure unless subsequent processing establishes that outcome.



ActiveStateToken, AuthoritativeVersion, externally meaningful version identifiers and standards-defined version metadata SHALL remain distinct where they represent different semantic domains.



Where an underlying engine or protocol lacks a distinction required by Harmonia, Harmonia SHALL introduce an appropriate semantic representation rather than infer meaning from incidental implementation details such as exception classes, message text or correlated version values.



Boundary projection MAY deliberately reduce or transform semantic distinctions where required by an external contract or policy, but such projection SHALL be explicit and SHALL NOT redefine Harmonia's internal semantics.



\par\medskip\textbf{Maxim}\par


Preserve semantics internally; project deliberately at boundaries.



This does not mean: Harmonia must expose every internal distinction externally. Similar outcomes can never share implementation machinery. Every semantic distinction requires a separate Java type. External standards or interface contracts may be ignored in favour of Harmonia's internal model.


{\footnotesize \href{../../../../../../../continuum/motivation/principle/d3fc6025dc7646d1b2020b51f7a44bfa.md}{Canonical record: d3fc6025dc7646d1b2020b51f7a44bfa}; full qualifications and provenance.}\par
\section*{I.7.15 AX-15 Uncertainty Is Preserved Until Resolved}
\phantomsection
\label{navigation:I.7.15}
\addcontentsline{toc}{section}{I.7.15 AX-15 Uncertainty Is Preserved Until Resolved}

\label{element:ef8abbaec37e4890ad78b9f263e67906}
\index{AX-15 Uncertainty Is Preserved Until Resolved (I.7.15)}
{\small Principle · Canonical · current}\par

\par\medskip\textbf{Axiom}\par


Where Harmonia cannot establish the state, outcome or effect of an operation with the certainty required by its governing contract, that uncertainty SHALL be represented explicitly.



Harmonia SHALL NOT infer, manufacture or assume a more certain outcome merely to simplify processing, recovery or application behaviour.



Uncertainty MAY be resolved through authoritative observation, reconciliation or other governed evidence capable of establishing the required state.



\par\medskip\textbf{Rationale}\par


Distributed operations can fail at points where the observable failure does not establish whether the requested operation occurred.



For example, loss of a connection after transmitting a request does not establish whether the receiving subsystem processed or committed that request. Treating such an outcome as either success or failure introduces information that Harmonia does not possess.



Preserving uncertainty allows subsequent processing to determine what actually occurred without compounding an unknown outcome through unsafe assumptions or repeated actions.



\par\medskip\textbf{Consequences}\par


An indeterminate operation outcome SHALL remain distinguishable from both confirmed success and confirmed non-occurrence.



Harmonia SHALL NOT automatically retry a state-changing operation where the previous attempt may have succeeded unless the operation's contract makes that retry demonstrably safe.



Recovery from an uncertain outcome SHOULD preferentially establish the current authoritative state and reconcile against it before further state-changing action is attempted.



Timeouts, connection failures and similar technical events SHALL NOT, by themselves, determine semantic outcome. Their interpretation SHALL depend upon what can be established about the operation's progress and the governing boundary contract.



Where subsequent authoritative observation or other governed evidence resolves an uncertainty, Harmonia MAY continue processing from the newly established state.



Uncertainty SHALL NOT cause transient, cached or process-local state to be promoted into authoritative state.



\par\medskip\textbf{Maxim}\par


\textbf{Do not turn “unknown” into “yes” or “no”.}



\textbf{This does not mean:} Harmonia must retain uncertainty indefinitely. Every technical failure produces an uncertain outcome. Operations known not to have crossed the relevant state-changing boundary cannot be safely retried. Reconciliation must always require human intervention.


{\footnotesize \href{../../../../../../../continuum/motivation/principle/ef8abbaec37e4890ad78b9f263e67906.md}{Canonical record: ef8abbaec37e4890ad78b9f263e67906}; full qualifications and provenance.}\par
\section*{I.7.16 AX-16 Operational Activity \& Entity State Progress Together}
\phantomsection
\label{navigation:I.7.16}
\addcontentsline{toc}{section}{I.7.16 AX-16 Operational Activity \& Entity State Progress Together}

\label{element:de7e90f46127455a9ec9ad3b9b48ba54}
\index{AX-16 Operational Activity \& Entity State Progress Together (I.7.16)}
{\small Principle · Canonical · current}\par

This axiom is maintained in this register. Current Harmonia realisation describes downstream machinery and is non-normative.



\par\medskip\textbf{Axiom}\par


\begin{quote}


\textbf{Operational activity associated with a real-world entity progresses through explicit and observable state transitions coordinated with the governed state of that entity. Operational activity is not treated merely as a sequence of disconnected message transfers.}



\end{quote}


\par\medskip\textbf{Architectural Implications}\par


\begin{itemize}


\item Health integration involves multi-stage workflows (e.g., ADT distribution, lab specimen processing).

\item Operational tasks must be tracked in lockstep with the real-world state of the clinical entity.

\item Prevents silent destination divergence during fan-out processing: every destination hop is an observable sub-state of the parent operational activity.

\end{itemize}


\par\medskip\textbf{Current Harmonia Realisation (Non-Normative)}\par


\begin{itemize}


\item Implemented in \texttt{\seqsplit{energeia-erga}} (e.g., \texttt{\seqsplit{AdtDistributionErgon}} tracking destination fan-out checkpoints).

\item Multi-destination statuses (\texttt{\seqsplit{FANOUT\_DISPATCH\_INITIATED}}, \texttt{\seqsplit{QUEUED}}, \texttt{\seqsplit{DELIVERED}}) are recorded in Pragma envelopes and FHIR \texttt{\seqsplit{Task.output}} extensions.

\end{itemize}


\par\medskip\textbf{Approved interpretation and related authority}\par


The 2026-10-09 migration adjudication preserves coordinated, explicit and
observable progression. The "lockstep" implication above does not mandate
synchronous execution, atomic entity/activity updates or a literal lockstep
implementation. AX-16 does not require universal Digital Twin routing, a Twin
for population/search activity, or clinical-work authority.



The \href{\detokenize{../../../../../../../../docs/markdown/02-strategy/strategic-views/logical-component-responsibilities.md#entity-specific-operational-coordination}}{Strategy entity-specific Digital Twin coordination rule}
remains the controlling Strategy-level specialisation of AX-16.
\hyperref[element:b8cafe3cf5904d2bb9a9767260fddf78]{AX-18} constrains inference of domain authority from coordination.


{\footnotesize \href{../../../../../../../continuum/motivation/principle/de7e90f46127455a9ec9ad3b9b48ba54.md}{Canonical record: de7e90f46127455a9ec9ad3b9b48ba54}; full qualifications and provenance.}\par
\section*{I.7.17 AX-17 Architectural Authority and Explicit Uncertainty}
\phantomsection
\label{navigation:I.7.17}
\addcontentsline{toc}{section}{I.7.17 AX-17 Architectural Authority and Explicit Uncertainty}

\label{element:8882573df7ce43c6bc0f66da0a676139}
\index{AX-17 Architectural Authority and Explicit Uncertainty (I.7.17)}
{\small Principle · Canonical · current}\par

\par\medskip\textbf{Axiom}\par


Documented architectural decisions, definitions, relationships, boundaries
and constraints are authoritative and SHALL NOT be silently reinterpreted,
replaced, bypassed or contradicted by downstream architecture or
implementation.



Where the architecture does not establish a fact, relationship,
responsibility or decision, that absence SHALL be preserved explicitly
rather than completed through inference, convention, structural
convenience, lexical similarity or anticipated implementation.



Unresolved architecture MAY be accompanied by clearly identified
candidates, suggestions or recommendations, but these SHALL remain
explicitly non-authoritative until explicitly accepted into the
architecture.



Where evidence indicates that authoritative architecture may be incorrect,
incomplete or internally inconsistent, the conflict SHALL be raised
explicitly for architectural review. The existing architecture remains
authoritative until an approved architectural change is made. That
approved change becomes the new authoritative architecture.



\par\medskip\textbf{Rationale}\par


Downstream derivation must distinguish established architectural knowledge
from gaps and proposed resolutions. A plausible relationship supplies no
authority merely because it makes a model or implementation appear
complete. Absence of an established relationship is architecturally
meaningful information.



Repeated architectural review during Domain04 Package2 G1 confirmed this
need across different traceability relationships:



\begin{itemize}


\item K9: unresolved Capability hierarchy must not be completed through
  invented ancestry.

\item K10: Behaviour / Process traceability must not be completed through
  invented stage equivalence.

\item K11: Function / Service traceability must not be completed through
  unsupported Feature association.

\item K12: cross-capability dependency traceability must not be completed
  through invented Service ownership or consumption.

\item K13: downstream representation traceability must not be completed through
  anticipated FHIR, application or technology mappings.

\end{itemize}


These examples justify an architecture-wide rule; its application does
not depend on Domain04 or G1 terminology.



\par\medskip\textbf{Consequences}\par


Human architects, human developers and AI development agents SHALL apply
this axiom to architecture derivation, documentation generation and
implementation.



Architectural traceability SHALL represent established relationships and
explicitly preserve unresolved or unestablished relationships.
Completeness SHALL NOT be manufactured through:



\begin{itemize}


\item inferred hierarchy;

\item structural convenience;

\item lexical similarity;

\item conventional modelling patterns;

\item assumed equivalence;

\item anticipated implementation;

\item familiar standards;

\item probable application design.

\end{itemize}


Candidates, suggestions and recommendations SHALL be distinguishable from
documented, approved architecture. A proposed resolution or apparently
better alternative SHALL NOT acquire authority through downstream use.



AI development agents SHOULD challenge architecture when evidence
warrants architectural review. They SHALL NOT silently bypass it because
they recognise a familiar destination.



This axiom applies the authority and explicit-review obligations in
\href{\detokenize{../../../../../../../governance/architectural-axioms.md#1-purpose}}{Purpose}, \href{\detokenize{../../../../../../../governance/architectural-axioms.md#6-relationship-to-architecture-decisions}}{Relationship to Architecture Decisions}
and \href{\detokenize{../../../../../../../governance/architectural-axioms.md#7-architectural-review-rule}}{Architectural Review Rule} to downstream
derivation and implementation.



\hyperref[element:d3fc6025dc7646d1b2020b51f7a44bfa]{AX-14} preserves semantic
distinctions; this axiom preserves the distinction between established,
unresolved and proposed architecture.
\hyperref[element:ef8abbaec37e4890ad78b9f263e67906]{AX-15} governs uncertainty
about operational state, outcome or effect; this axiom governs uncertainty
about the architecture itself. Neither form of uncertainty may be silently
resolved by assumption.



\hyperref[element:eee8c3783d964d18a7f9dcea85fcd4cf]{AX-02}
and \hyperref[element:54508cda80484a6caf95b901780eaf7d]{AX-04}
preserve Harmonia's internal semantics from external standards and engine
assumptions. Familiar standards or technology patterns do not establish
otherwise undocumented architectural relationships.
\hyperref[element:952c3b6148ce4fce94b69f99a0c2fec7]{AX-06} governs the authority
of managed information; architectural authority remains a distinct concern.



\par\medskip\textbf{Maxim}\par


\textbf{Traceability must be truthful, not artificially complete.}


{\footnotesize \href{../../../../../../../continuum/motivation/principle/8882573df7ce43c6bc0f66da0a676139.md}{Canonical record: 8882573df7ce43c6bc0f66da0a676139}; full qualifications and provenance.}\par
\section*{I.7.18 AX-18 Integration and Coordination Do Not Confer Domain Authority}
\phantomsection
\label{navigation:I.7.18}
\addcontentsline{toc}{section}{I.7.18 AX-18 Integration and Coordination Do Not Confer Domain Authority}

\label{element:b8cafe3cf5904d2bb9a9767260fddf78}
\index{AX-18 Integration and Coordination Do Not Confer Domain Authority (I.7.18)}
{\small Principle · Canonical · current}\par

\par\medskip\textbf{Axiom}\par


\begin{quote}


Harmonia MAY integrate information concerning, and coordinate operational activity associated with, an externally governed domain activity without thereby acquiring authority to determine or manage that domain activity.



Authority to determine or manage a domain activity SHALL remain with the responsible actor, system or governed capability unless that authority is explicitly assigned to Harmonia by the architecture.



Harmonia's observation, processing, persistence, distribution, presentation or use of information concerning a domain activity SHALL NOT, by itself, establish authority over that activity.



\end{quote}


\par\medskip\textbf{Rationale}\par


Participation in a domain activity and authority over that activity are
different architectural concerns. Information possession, observation,
integration, coordination and processing do not, by themselves, confer
domain authority. Authority requires explicit architectural assignment.



This new human-approved axiom generalises an already-established class of
architectural boundary. It is established at the 2026-10-09 migration
boundary; it is not a claim that AX-01 through AX-17 always contained this
axiom.



\par\medskip\textbf{Consequences and qualifications}\par


AX-18 SHALL NOT be interpreted as preventing Harmonia from performing
operational activity for which architecture has assigned Harmonia
responsibility. Explicit architectural assignment may establish authority
to determine or manage a domain activity.



In particular, AX-18 does not mean:



\begin{itemize}


\item Harmonia cannot execute Praxis/Ergo activity;

\item Harmonia cannot coordinate entity-specific operational activity through a Digital Twin;

\item Harmonia cannot process or transform clinical information;

\item Harmonia cannot perform workflow/activity coordination;

\item Harmonia cannot manage responsibilities explicitly assigned to Harmonia;

\item Harmonia cannot enforce policy, security, information governance or operational controls assigned to it.

\end{itemize}


Participation, information, integration and coordination do not, by
themselves, establish authority to determine or manage the underlying domain
activity. Assigned Harmonia operational responsibility remains applicable.



\par\medskip\textbf{Examples and established specialisations (non-normative)}\par


These examples explain the boundary; they do not expand the axiom or
establish Application Components, Information Concepts, Services, Processes,
AI allocation or implementation mechanisms.



\begin{itemize}


\item \textbf{Clinical work:} Harmonia may integrate and coordinate around clinical work
  without thereby acquiring authority to determine clinical work, manage
  clinical handover, own clinical worklists, allocate clinical work, make
  clinical decisions or become the EMR workflow authority. The detailed
  \href{\detokenize{../../../../../../../../docs/markdown/03-business-architecture/metamodel/business-architecture-metamodel.md#37-clinical-work-integration-boundary}}{Domain03 clinical-work integration boundary}
  remains the Business specialisation.

\item \textbf{Digital Twins:} A Twin may coordinate operational activity associated
  with its real-world entity without thereby acquiring authority over
  originating clinical facts or domain decisions. The
  \href{\detokenize{../../../../../../../../docs/markdown/02-strategy/strategic-views/logical-component-responsibilities.md#entity-specific-operational-coordination}}{Strategy entity-specific coordination rule}
  remains controlling for that coordination.

\item \textbf{Capacity:} Higher-order Service Capacity Management may consume
  context-specific capacity without acquiring the originating responsibility
  for Clinic, Ward, Theatre, ED, Outpatient or Mobile Service Capacity
  Management. The detailed
  \href{\detokenize{../../../../../../../../docs/markdown/03-business-architecture/behaviours/04-health-service-operations.md#capacity-management-responsibility-boundary}}{Business capacity responsibility boundary}
  remains in Business Architecture.

\item \textbf{Assurance:} Observation, access to evidence or establishment of an
  assurance conclusion does not confer management authority over the activity
  being assured. Independent Assurance remains governed by
  \hyperref[element:bbb7438e78604308abba93afc14f4e2e]{REQ-FND-005}
  and its established
  \href{\detokenize{../../../../../../../../docs/markdown/02-strategy/capability-maps/health-service-assurance-derivation.md}}{Strategy}
  and \href{\detokenize{../../../../../../../../docs/markdown/03-business-architecture/behaviours/health-service-assurance.md}}{Business}
  derivation.

\item \textbf{Presentation / UI:} Presentation or collection of information does not,
  by itself, confer originating information or domain-decision authority on
  the presenting component. The existing
  \href{\detokenize{../../../../../../../../docs/markdown/02-strategy/strategic-views/logical-component-responsibilities.md#component-6-iris-contextual-human-interaction}}{Iris contextual human interaction responsibility}
  remains distinct.

\item \textbf{AI:} Assistance with interpretation, pattern recognition, criteria
  evaluation or decision support does not, by itself, confer the authority of
  the governed actor, role or capability using that assistance. This example
  assigns no AI responsibility or capability.

\end{itemize}


\par\medskip\textbf{Related architectural authority}\par


\begin{itemize}


\item \hyperref[element:952c3b6148ce4fce94b69f99a0c2fec7]{AX-06} governs information authority and credibility.

\item \hyperref[element:b6a55decfeec4f1bb3d7f18534a2ffea]{AX-13} governs the Harmonia management boundary at ingress/egress.

\item \hyperref[element:d3fc6025dc7646d1b2020b51f7a44bfa]{AX-14} preserves meaningful semantic distinctions.

\item \hyperref[element:de7e90f46127455a9ec9ad3b9b48ba54]{AX-16} governs coordinated operational activity/entity progression.

\item \hyperref[element:8882573df7ce43c6bc0f66da0a676139]{AX-17} governs architectural authority and explicit uncertainty.

\item AX-18 governs inference of domain authority from participation, information,
  integration or coordination.

\end{itemize}

{\footnotesize \href{../../../../../../../continuum/motivation/principle/b8cafe3cf5904d2bb9a9767260fddf78.md}{Canonical record: b8cafe3cf5904d2bb9a9767260fddf78}; full qualifications and provenance.}\par
\chapter*{I.8 Unclassified Motivation}
\phantomsection
\label{navigation:I.8}
\addcontentsline{toc}{chapter}{I.8 Unclassified Motivation}

\section*{I.8.1 At-Least-Once}
\phantomsection
\label{navigation:I.8.1}
\addcontentsline{toc}{section}{I.8.1 At-Least-Once}

\label{element:bd4c8726125648e68d5746bab3c765cd}
\index{At-Least-Once (I.8.1)}
{\small Principle · Historical / Reclassified — Domain06 · historical\_reference}\par
\href{../../../../../../../continuum/motivation/principle/bd4c8726125648e68d5746bab3c765cd.md}{Historical canonical definition and provenance}.\par
\section*{I.8.2 Canonical Modeling with Standardized FHIR R5 Boundaries}
\phantomsection
\label{navigation:I.8.2}
\addcontentsline{toc}{section}{I.8.2 Canonical Modeling with Standardized FHIR R5 Boundaries}

\label{element:5d97974406c546078b1d64e5ff85bb13}
\index{Canonical Modeling with Standardized FHIR R5 Boundaries (I.8.2)}
{\small Principle · Historical / Reclassified · historical\_reference}\par
\href{../../../../../../../continuum/motivation/principle/5d97974406c546078b1d64e5ff85bb13.md}{Historical canonical definition and provenance}.\par
\section*{I.8.3 Governed Asynchronous Change Management for Directory Entities}
\phantomsection
\label{navigation:I.8.3}
\addcontentsline{toc}{section}{I.8.3 Governed Asynchronous Change Management for Directory Entities}

\label{element:11001a6ea3df45689a27216df641dac2}
\index{Governed Asynchronous Change Management for Directory Entities (I.8.3)}
{\small Principle · Historical — current applicability unresolved · historical\_reference}\par
\href{../../../../../../../continuum/motivation/principle/11001a6ea3df45689a27216df641dac2.md}{Historical canonical definition and provenance}.\par
\section*{I.8.4 Governed PHI Logging \& Dual-Logger Privacy Controls}
\phantomsection
\label{navigation:I.8.4}
\addcontentsline{toc}{section}{I.8.4 Governed PHI Logging \& Dual-Logger Privacy Controls}

\label{element:96ed278624c34ba4bfade0df08d45d74}
\index{Governed PHI Logging \& Dual-Logger Privacy Controls (I.8.4)}
{\small Principle · Historical / Reclassified · historical\_reference}\par
\href{../../../../../../../continuum/motivation/principle/96ed278624c34ba4bfade0df08d45d74.md}{Historical canonical definition and provenance}.\par
\section*{I.8.5 Provenance and Information Authority}
\phantomsection
\label{navigation:I.8.5}
\addcontentsline{toc}{section}{I.8.5 Provenance and Information Authority}

\label{element:638c501ba4df4c67aee92d1f328aad84}
\index{Provenance and Information Authority (I.8.5)}
{\small Principle · Historical — current applicability unresolved · historical\_reference}\par
\href{../../../../../../../continuum/motivation/principle/638c501ba4df4c67aee92d1f328aad84.md}{Historical canonical definition and provenance}.\par
\section*{I.8.6 Separation of Transport from Business Logic}
\phantomsection
\label{navigation:I.8.6}
\addcontentsline{toc}{section}{I.8.6 Separation of Transport from Business Logic}

\label{element:5ae28648334e4f6184044af0462eb9a8}
\index{Separation of Transport from Business Logic (I.8.6)}
{\small Principle · Historical / Reclassified · historical\_reference}\par
\href{../../../../../../../continuum/motivation/principle/5ae28648334e4f6184044af0462eb9a8.md}{Historical canonical definition and provenance}.\par
\section*{I.8.7 Write-Behind}
\phantomsection
\label{navigation:I.8.7}
\addcontentsline{toc}{section}{I.8.7 Write-Behind}

\label{element:a5dffc0325d44915971eea600bba48f1}
\index{Write-Behind (I.8.7)}
{\small Principle · Retired / Superseded · historical\_reference}\par
\href{../../../../../../../continuum/motivation/principle/a5dffc0325d44915971eea600bba48f1.md}{Historical canonical definition and provenance}.\par
\chapter*{I.9 Requirement Groups}
\phantomsection
\label{navigation:I.9}
\addcontentsline{toc}{chapter}{I.9 Requirement Groups}

\section*{I.9.1 [RequirementGroup] Functional Requirements}
\phantomsection
\label{navigation:I.9.1}
\addcontentsline{toc}{section}{I.9.1 [RequirementGroup] Functional Requirements}

\label{element:6acde68c43b14b6d8c7707779abdfd53}
\index{[RequirementGroup] Functional Requirements (I.9.1)}
{\small Grouping / RequirementGroup · Accepted bounded architecture — HCRL 1.1.15-001 DR-1 · current}\par

Functional Requirements grouping hierarchy root accepted under HCRL 1.1.15-001 DR-1.


{\small Recorded unresolved matters: canonicalName, friendlyName; exact scope remains in the canonical record.}\par
{\footnotesize \href{../../../../../../../continuum/grouping/requirement-group/6acde68c43b14b6d8c7707779abdfd53.md}{Canonical record: 6acde68c43b14b6d8c7707779abdfd53}; full qualifications and provenance.}\par
\subsection*{I.9.1.1 [RequirementGroup] Information Management}
\phantomsection
\label{navigation:I.9.1.1}
\addcontentsline{toc}{subsection}{I.9.1.1 [RequirementGroup] Information Management}

\label{element:c34c366824ff452eb575345c16db720b}
\index{[RequirementGroup] Information Management (I.9.1.1)}
{\small Grouping / RequirementGroup · Accepted bounded architecture — HCRL 1.1.15-001 DR-1 · current}\par

Information management requirements grouping nested under Functional Requirements accepted under HCRL 1.1.15-001 DR-1.


{\small Recorded unresolved matters: canonicalName, friendlyName; exact scope remains in the canonical record.}\par
{\footnotesize \href{../../../../../../../continuum/grouping/requirement-group/c34c366824ff452eb575345c16db720b.md}{Canonical record: c34c366824ff452eb575345c16db720b}; full qualifications and provenance.}\par
\subsubsection*{I.9.1.1.1 [Requirement] Information Durability During Platform Failure}
\phantomsection
\label{navigation:I.9.1.1.1}
\addcontentsline{toc}{subsubsection}{I.9.1.1.1 [Requirement] Information Durability During Platform Failure}

\label{element:6a2a3de03beb47b884293797b92a91eb}
\index{[Requirement] Information Durability During Platform Failure (I.9.1.1.1)}
{\small Requirement · Accepted bounded architecture — HCRL 1.1.15-001 DR-1 · current}\par

Intentionally incomplete requirement placeholder / thinking point accepted under HCRL 1.1.15-001 DR-1. Contained beneath [RequirementGroup] Information Management.


{\small Recorded unresolved matters: canonicalName, friendlyName; exact scope remains in the canonical record.}\par
{\footnotesize \href{../../../../../../../continuum/motivation/requirement/6a2a3de03beb47b884293797b92a91eb.md}{Canonical record: 6a2a3de03beb47b884293797b92a91eb}; full qualifications and provenance.}\par
\section*{I.9.2 [RequirementGroup] Non-Functional-Requirements}
\phantomsection
\label{navigation:I.9.2}
\addcontentsline{toc}{section}{I.9.2 [RequirementGroup] Non-Functional-Requirements}

\label{element:f4a3060eba464f4ab89b50aec4628e5b}
\index{[RequirementGroup] Non-Functional-Requirements (I.9.2)}
{\small Grouping / RequirementGroup · Accepted bounded architecture — HCRL 1.1.15-001 DR-1 · current}\par

Non-Functional Requirements grouping hierarchy root accepted under HCRL 1.1.15-001 DR-1.


{\small Recorded unresolved matters: canonicalName, friendlyName; exact scope remains in the canonical record.}\par
{\footnotesize \href{../../../../../../../continuum/grouping/requirement-group/f4a3060eba464f4ab89b50aec4628e5b.md}{Canonical record: f4a3060eba464f4ab89b50aec4628e5b}; full qualifications and provenance.}\par
\subsection*{I.9.2.1 [RequirementGroup] Performance Requirements}
\phantomsection
\label{navigation:I.9.2.1}
\addcontentsline{toc}{subsection}{I.9.2.1 [RequirementGroup] Performance Requirements}

\label{element:dc3f2a927e284da88734df4f7e0f9311}
\index{[RequirementGroup] Performance Requirements (I.9.2.1)}
{\small Grouping / RequirementGroup · Accepted bounded architecture — HCRL 1.1.15-001 DR-1 · current}\par

Performance requirements grouping nested under Non-Functional-Requirements accepted under HCRL 1.1.15-001 DR-1.


{\small Recorded unresolved matters: canonicalName, friendlyName; exact scope remains in the canonical record.}\par
{\footnotesize \href{../../../../../../../continuum/grouping/requirement-group/dc3f2a927e284da88734df4f7e0f9311.md}{Canonical record: dc3f2a927e284da88734df4f7e0f9311}; full qualifications and provenance.}\par
\subsection*{I.9.2.2 [RequirementGroup] Disaster Recovery Requirements}
\phantomsection
\label{navigation:I.9.2.2}
\addcontentsline{toc}{subsection}{I.9.2.2 [RequirementGroup] Disaster Recovery Requirements}

\label{element:cb090e968a25409fb8ceb0fcc9a11ee7}
\index{[RequirementGroup] Disaster Recovery Requirements (I.9.2.2)}
{\small Grouping / RequirementGroup · Accepted bounded architecture — HCRL 1.1.15-001 DR-1 · current}\par

Disaster recovery requirements grouping nested under Non-Functional-Requirements accepted under HCRL 1.1.15-001 DR-1.


{\small Recorded unresolved matters: canonicalName, friendlyName; exact scope remains in the canonical record.}\par
{\footnotesize \href{../../../../../../../continuum/grouping/requirement-group/cb090e968a25409fb8ceb0fcc9a11ee7.md}{Canonical record: cb090e968a25409fb8ceb0fcc9a11ee7}; full qualifications and provenance.}\par
\subsection*{I.9.2.3 [RequirementGroup] Availability}
\phantomsection
\label{navigation:I.9.2.3}
\addcontentsline{toc}{subsection}{I.9.2.3 [RequirementGroup] Availability}

\label{element:abf16578801a466b8f08d80ae292e644}
\index{[RequirementGroup] Availability (I.9.2.3)}
{\small Grouping / RequirementGroup · Accepted bounded architecture — HCRL 1.1.15-001 DR-1 · current}\par

Availability requirements grouping nested under Non-Functional-Requirements accepted under HCRL 1.1.15-001 DR-1.


{\small Recorded unresolved matters: canonicalName, friendlyName; exact scope remains in the canonical record.}\par
{\footnotesize \href{../../../../../../../continuum/grouping/requirement-group/abf16578801a466b8f08d80ae292e644.md}{Canonical record: abf16578801a466b8f08d80ae292e644}; full qualifications and provenance.}\par
\subsubsection*{I.9.2.3.1 [RequirementGroup] Information Availability}
\phantomsection
\label{navigation:I.9.2.3.1}
\addcontentsline{toc}{subsubsection}{I.9.2.3.1 [RequirementGroup] Information Availability}

\label{element:1862d183cf7d45c4972e6aa3411d1193}
\index{[RequirementGroup] Information Availability (I.9.2.3.1)}
{\small Grouping / RequirementGroup · Accepted bounded architecture — HCRL 1.1.15-001 DR-1 · current}\par

Information availability requirements grouping nested under Availability accepted under HCRL 1.1.15-001 DR-1.


{\small Recorded unresolved matters: canonicalName, friendlyName; exact scope remains in the canonical record.}\par
{\footnotesize \href{../../../../../../../continuum/grouping/requirement-group/1862d183cf7d45c4972e6aa3411d1193.md}{Canonical record: 1862d183cf7d45c4972e6aa3411d1193}; full qualifications and provenance.}\par
\chapter*{I.10 Principles}
\phantomsection
\label{navigation:I.10}
\addcontentsline{toc}{chapter}{I.10 Principles}

Iteration 001 native Principles and prescriptive ArchitecturalAxiom specialisations; existing highest-authority AX definitions remain in Architecture Axioms.\par
\href{../../../../../../offline-design/Harmoni-NFR-Derived Architetural Elements - Iteration 001.md}{Governed source: Principles}.\par
\section*{I.10.1 [ArchitecturalAxiom] Conformance to Governed Architecture}
\phantomsection
\label{navigation:I.10.1}
\addcontentsline{toc}{section}{I.10.1 [ArchitecturalAxiom] Conformance to Governed Architecture}

\label{element:45eb197073df406ea4080f9665ca12ab}
\index{[ArchitecturalAxiom] Conformance to Governed Architecture (I.10.1)}
{\small Principle / ArchitecturalAxiom · Solution Architect-reviewed Iteration 001; authorised ADD; resulting increment review pending · current}\par

Conform to governed Harmonia architecture



Source entry: \textbf{NFR-001 — Conformance to Governed Architecture}. Primary subject: \texttt{\seqsplit{G01.02}}.



The source classification is Principle «ArchitecturalAxiom». The source identifier is a scoped register reference and does not determine architectural type. Applicability, release scope and unresolved thresholds remain as stated above.


{\footnotesize \href{../../../../../../../continuum/motivation/principle/architectural-axiom/45eb197073df406ea4080f9665ca12ab.md}{Canonical record: 45eb197073df406ea4080f9665ca12ab}; full qualifications and provenance.}\par
\section*{I.10.2 [ArchitecturalAxiom] Stable Architectural Identity}
\phantomsection
\label{navigation:I.10.2}
\addcontentsline{toc}{section}{I.10.2 [ArchitecturalAxiom] Stable Architectural Identity}

\label{element:86ecebbf65fc46af9763c96840342823}
\index{[ArchitecturalAxiom] Stable Architectural Identity (I.10.2)}
{\small Principle / ArchitecturalAxiom · Solution Architect-reviewed Iteration 001; authorised ADD; resulting increment review pending · current}\par

Preserve stable architectural identities across transformations and round trips



Source entry: \textbf{NFR-004 — Stable Architectural Identity}. Primary subject: \texttt{\seqsplit{G01.03}}.



The source classification is Principle «ArchitecturalAxiom». The source identifier is a scoped register reference and does not determine architectural type. Applicability, release scope and unresolved thresholds remain as stated above.


{\footnotesize \href{../../../../../../../continuum/motivation/principle/architectural-axiom/86ecebbf65fc46af9763c96840342823.md}{Canonical record: 86ecebbf65fc46af9763c96840342823}; full qualifications and provenance.}\par
\section*{I.10.3 [ArchitecturalAxiom] Authoritative Projection Semantics}
\phantomsection
\label{navigation:I.10.3}
\addcontentsline{toc}{section}{I.10.3 [ArchitecturalAxiom] Authoritative Projection Semantics}

\label{element:d8ae43c850a646d485e6aba6c41a88e2}
\index{[ArchitecturalAxiom] Authoritative Projection Semantics (I.10.3)}
{\small Principle / ArchitecturalAxiom · Solution Architect-reviewed Iteration 001; authorised ADD; resulting increment review pending · current}\par

Preserve authoritative semantics in architectural projections



Source entry: \textbf{NFR-005 — Authoritative Projection Semantics}. Primary subject: \texttt{\seqsplit{G01.03}}.



The source classification is Principle «ArchitecturalAxiom». The source identifier is a scoped register reference and does not determine architectural type. Applicability, release scope and unresolved thresholds remain as stated above.


{\footnotesize \href{../../../../../../../continuum/motivation/principle/architectural-axiom/d8ae43c850a646d485e6aba6c41a88e2.md}{Canonical record: d8ae43c850a646d485e6aba6c41a88e2}; full qualifications and provenance.}\par
\section*{I.10.4 [ArchitecturalAxiom] Independent Architecture Verification}
\phantomsection
\label{navigation:I.10.4}
\addcontentsline{toc}{section}{I.10.4 [ArchitecturalAxiom] Independent Architecture Verification}

\label{element:82d29abf230d478aaf42ecb73ecbd01e}
\index{[ArchitecturalAxiom] Independent Architecture Verification (I.10.4)}
{\small Principle / ArchitecturalAxiom · Solution Architect-reviewed Iteration 001; authorised ADD; resulting increment review pending · current}\par

Independently verify implementations against architectural contracts



Source entry: \textbf{NFR-007 — Independent Architecture Verification}. Primary subject: \texttt{\seqsplit{G01.02}}.



The source classification is Principle «ArchitecturalAxiom». The source identifier is a scoped register reference and does not determine architectural type. Applicability, release scope and unresolved thresholds remain as stated above.


{\footnotesize \href{../../../../../../../continuum/motivation/principle/architectural-axiom/82d29abf230d478aaf42ecb73ecbd01e.md}{Canonical record: 82d29abf230d478aaf42ecb73ecbd01e}; full qualifications and provenance.}\par
\section*{I.10.5 [ArchitecturalAxiom] Defined Service Contracts}
\phantomsection
\label{navigation:I.10.5}
\addcontentsline{toc}{section}{I.10.5 [ArchitecturalAxiom] Defined Service Contracts}

\label{element:b7d07bb4ec064272898d2f0189cffefd}
\index{[ArchitecturalAxiom] Defined Service Contracts (I.10.5)}
{\small Principle / ArchitecturalAxiom · Solution Architect-reviewed Iteration 001; authorised ADD; resulting increment review pending · current}\par

Expose defined service contracts



Source entry: \textbf{NFR-008 — Defined Service Contracts}. Primary subject: \texttt{\seqsplit{G01.03}}.



The source classification is Principle «ArchitecturalAxiom». The source identifier is a scoped register reference and does not determine architectural type. Applicability, release scope and unresolved thresholds remain as stated above.


{\footnotesize \href{../../../../../../../continuum/motivation/principle/architectural-axiom/b7d07bb4ec064272898d2f0189cffefd.md}{Canonical record: b7d07bb4ec064272898d2f0189cffefd}; full qualifications and provenance.}\par
\section*{I.10.6 [ArchitecturalAxiom] Service Boundary Encapsulation}
\phantomsection
\label{navigation:I.10.6}
\addcontentsline{toc}{section}{I.10.6 [ArchitecturalAxiom] Service Boundary Encapsulation}

\label{element:376363587388401194ff967b0fe80a5a}
\index{[ArchitecturalAxiom] Service Boundary Encapsulation (I.10.6)}
{\small Principle / ArchitecturalAxiom · Solution Architect-reviewed Iteration 001; authorised ADD; resulting increment review pending · current}\par

Prevent implementation-specific technology leakage across service boundaries



Source entry: \textbf{NFR-009 — Service Boundary Encapsulation}. Primary subject: \texttt{\seqsplit{G01.03}}.



The source classification is Principle «ArchitecturalAxiom». The source identifier is a scoped register reference and does not determine architectural type. Applicability, release scope and unresolved thresholds remain as stated above.


{\footnotesize \href{../../../../../../../continuum/motivation/principle/architectural-axiom/376363587388401194ff967b0fe80a5a.md}{Canonical record: 376363587388401194ff967b0fe80a5a}; full qualifications and provenance.}\par
\section*{I.10.7 [ArchitecturalAxiom] Governed Information Identity}
\phantomsection
\label{navigation:I.10.7}
\addcontentsline{toc}{section}{I.10.7 [ArchitecturalAxiom] Governed Information Identity}

\label{element:c2e7ee7597e64275b3523befd72c0987}
\index{[ArchitecturalAxiom] Governed Information Identity (I.10.7)}
{\small Principle / ArchitecturalAxiom · Solution Architect-reviewed Iteration 001; authorised ADD; resulting increment review pending · current}\par

Governed information elements shall possess well-defined, traceable identities that remain consistent across their points of utilisation.



Source entry: \textbf{NFR-014 — Governed Information Identity}. Primary subject: \texttt{\seqsplit{G02.01}}.



The source classification is Principle «ArchitecturalAxiom». The source identifier is a scoped register reference and does not determine architectural type. Applicability, release scope and unresolved thresholds remain as stated above.


{\footnotesize \href{../../../../../../../continuum/motivation/principle/architectural-axiom/c2e7ee7597e64275b3523befd72c0987.md}{Canonical record: c2e7ee7597e64275b3523befd72c0987}; full qualifications and provenance.}\par
\section*{I.10.8 [ArchitecturalAxiom] Governed Information Versioning}
\phantomsection
\label{navigation:I.10.8}
\addcontentsline{toc}{section}{I.10.8 [ArchitecturalAxiom] Governed Information Versioning}

\label{element:2bd0132f072b42e9bf71bf483eb2deec}
\index{[ArchitecturalAxiom] Governed Information Versioning (I.10.8)}
{\small Principle / ArchitecturalAxiom · Solution Architect-reviewed Iteration 001; authorised ADD; resulting increment review pending · current}\par

Governed information shall support versioning, enabling individual versions to be identified, accessed and referenced unambiguously.



Source entry: \textbf{NFR-015 — Governed Information Versioning}. Primary subject: \texttt{\seqsplit{G02.03}}.



The source classification is Principle «ArchitecturalAxiom». The source identifier is a scoped register reference and does not determine architectural type. Applicability, release scope and unresolved thresholds remain as stated above.


{\footnotesize \href{../../../../../../../continuum/motivation/principle/architectural-axiom/2bd0132f072b42e9bf71bf483eb2deec.md}{Canonical record: 2bd0132f072b42e9bf71bf483eb2deec}; full qualifications and provenance.}\par
\section*{I.10.9 [ArchitecturalAxiom] Concurrency Policy Conformance}
\phantomsection
\label{navigation:I.10.9}
\addcontentsline{toc}{section}{I.10.9 [ArchitecturalAxiom] Concurrency Policy Conformance}

\label{element:7e3b05a2e3e547cbbd26b1e181af0d10}
\index{[ArchitecturalAxiom] Concurrency Policy Conformance (I.10.9)}
{\small Principle / ArchitecturalAxiom · Solution Architect-reviewed Iteration 001; authorised ADD; resulting increment review pending · current}\par

Concurrent activities shall produce outcomes consistent with their defined concurrency policies and applicable architectural and operational constraints.



Source entry: \textbf{NFR-025 — Concurrency Policy Conformance}. Primary subject: \texttt{\seqsplit{G06.03}}.



The source classification is Principle «ArchitecturalAxiom». The source identifier is a scoped register reference and does not determine architectural type. Applicability, release scope and unresolved thresholds remain as stated above.


{\footnotesize \href{../../../../../../../continuum/motivation/principle/architectural-axiom/7e3b05a2e3e547cbbd26b1e181af0d10.md}{Canonical record: 7e3b05a2e3e547cbbd26b1e181af0d10}; full qualifications and provenance.}\par
\section*{I.10.10 [ArchitecturalAxiom] Concurrent Task State Integrity}
\phantomsection
\label{navigation:I.10.10}
\addcontentsline{toc}{section}{I.10.10 [ArchitecturalAxiom] Concurrent Task State Integrity}

\label{element:26f3759e1d014a7cae5f265ea6530000}
\index{[ArchitecturalAxiom] Concurrent Task State Integrity (I.10.10)}
{\small Principle / ArchitecturalAxiom · Solution Architect-reviewed Iteration 001; authorised ADD; resulting increment review pending · current}\par

Preserve integrity of task state transitions under concurrency



Source entry: \textbf{NFR-028 — Concurrent Task State Integrity}. Primary subject: \texttt{\seqsplit{G07.01}}.



The source classification is Principle «ArchitecturalAxiom». The source identifier is a scoped register reference and does not determine architectural type. Applicability, release scope and unresolved thresholds remain as stated above.


{\footnotesize \href{../../../../../../../continuum/motivation/principle/architectural-axiom/26f3759e1d014a7cae5f265ea6530000.md}{Canonical record: 26f3759e1d014a7cae5f265ea6530000}; full qualifications and provenance.}\par
\section*{I.10.11 [ArchitecturalAxiom] Failure Observability}
\phantomsection
\label{navigation:I.10.11}
\addcontentsline{toc}{section}{I.10.11 [ArchitecturalAxiom] Failure Observability}

\label{element:2ed20179a98d44a68880958155e47923}
\index{[ArchitecturalAxiom] Failure Observability (I.10.11)}
{\small Principle / ArchitecturalAxiom · Solution Architect-reviewed Iteration 001; authorised ADD; resulting increment review pending · current}\par

Make processing failures observable and attributable



Source entry: \textbf{NFR-034 — Failure Observability}. Primary subject: \texttt{\seqsplit{G09.02}}.



The source classification is Principle «ArchitecturalAxiom». The source identifier is a scoped register reference and does not determine architectural type. Applicability, release scope and unresolved thresholds remain as stated above.


{\footnotesize \href{../../../../../../../continuum/motivation/principle/architectural-axiom/2ed20179a98d44a68880958155e47923.md}{Canonical record: 2ed20179a98d44a68880958155e47923}; full qualifications and provenance.}\par
\section*{I.10.12 [ArchitecturalAxiom] Security Information Integrity}
\phantomsection
\label{navigation:I.10.12}
\addcontentsline{toc}{section}{I.10.12 [ArchitecturalAxiom] Security Information Integrity}

\label{element:52083b57c01d4262bf965afa96489136}
\index{[ArchitecturalAxiom] Security Information Integrity (I.10.12)}
{\small Principle / ArchitecturalAxiom · Solution Architect-reviewed Iteration 001; authorised ADD; resulting increment review pending · current}\par

Preserve integrity of security-relevant information across interactions



Source entry: \textbf{NFR-042 — Security Information Integrity}. Primary subject: \texttt{\seqsplit{G03.04}}.



The source classification is Principle «ArchitecturalAxiom». The source identifier is a scoped register reference and does not determine architectural type. Applicability, release scope and unresolved thresholds remain as stated above.


{\footnotesize \href{../../../../../../../continuum/motivation/principle/architectural-axiom/52083b57c01d4262bf965afa96489136.md}{Canonical record: 52083b57c01d4262bf965afa96489136}; full qualifications and provenance.}\par
\section*{I.10.13 [ArchitecturalAxiom] Exchange Semantic Preservation}
\phantomsection
\label{navigation:I.10.13}
\addcontentsline{toc}{section}{I.10.13 [ArchitecturalAxiom] Exchange Semantic Preservation}

\label{element:602eb50793f448d8bdbc92d41f851877}
\index{[ArchitecturalAxiom] Exchange Semantic Preservation (I.10.13)}
{\small Principle / ArchitecturalAxiom · Solution Architect-reviewed Iteration 001; authorised ADD; resulting increment review pending · current}\par

Preserve semantics of exchanged information



Source entry: \textbf{NFR-050 — Exchange Semantic Preservation}. Primary subject: \texttt{\seqsplit{G05.02}}.



The source classification is Principle «ArchitecturalAxiom». The source identifier is a scoped register reference and does not determine architectural type. Applicability, release scope and unresolved thresholds remain as stated above.


{\footnotesize \href{../../../../../../../continuum/motivation/principle/architectural-axiom/602eb50793f448d8bdbc92d41f851877.md}{Canonical record: 602eb50793f448d8bdbc92d41f851877}; full qualifications and provenance.}\par
\section*{I.10.14 [ArchitecturalAxiom] Interoperability Encapsulation}
\phantomsection
\label{navigation:I.10.14}
\addcontentsline{toc}{section}{I.10.14 [ArchitecturalAxiom] Interoperability Encapsulation}

\label{element:e9e0b4f82b6544e7ba87651c578d82c2}
\index{[ArchitecturalAxiom] Interoperability Encapsulation (I.10.14)}
{\small Principle / ArchitecturalAxiom · Solution Architect-reviewed Iteration 001; authorised ADD; resulting increment review pending · current}\par

Avoid unnecessary exposure of implementation technology in interoperability mechanisms



Source entry: \textbf{NFR-051 — Interoperability Encapsulation}. Primary subject: \texttt{\seqsplit{G05.03}}.



The source classification is Principle «ArchitecturalAxiom». The source identifier is a scoped register reference and does not determine architectural type. Applicability, release scope and unresolved thresholds remain as stated above.


{\footnotesize \href{../../../../../../../continuum/motivation/principle/architectural-axiom/e9e0b4f82b6544e7ba87651c578d82c2.md}{Canonical record: e9e0b4f82b6544e7ba87651c578d82c2}; full qualifications and provenance.}\par
\section*{I.10.15 Contract Evolution Compatibility}
\phantomsection
\label{navigation:I.10.15}
\addcontentsline{toc}{section}{I.10.15 Contract Evolution Compatibility}

\label{element:0da415a142c640d9ab1c73d770df8ea3}
\index{Contract Evolution Compatibility (I.10.15)}
{\small Principle · Solution Architect-reviewed Iteration 001; authorised ADD; resulting increment review pending · current}\par

Minimise disruption to consumers when contracts evolve



Source entry: \textbf{NFR-054 — Contract Evolution Compatibility}. Primary subject: \texttt{\seqsplit{G05.04}}.



The source classification is Principle. The source identifier is a scoped register reference and does not determine architectural type. Applicability, release scope and unresolved thresholds remain as stated above.


{\footnotesize \href{../../../../../../../continuum/motivation/principle/0da415a142c640d9ab1c73d770df8ea3.md}{Canonical record: 0da415a142c640d9ab1c73d770df8ea3}; full qualifications and provenance.}\par
\section*{I.10.16 [ArchitecturalAxiom] Resource Contention Integrity}
\phantomsection
\label{navigation:I.10.16}
\addcontentsline{toc}{section}{I.10.16 [ArchitecturalAxiom] Resource Contention Integrity}

\label{element:db8d320dc4c0454080f8c0f3c4e52801}
\index{[ArchitecturalAxiom] Resource Contention Integrity (I.10.16)}
{\small Principle / ArchitecturalAxiom · Solution Architect-reviewed Iteration 001; authorised ADD; resulting increment review pending · current}\par

Prevent resource contention from compromising integrity



Source entry: \textbf{NFR-064 — Resource Contention Integrity}. Primary subject: \texttt{\seqsplit{G08.04}}.



The source classification is Principle «ArchitecturalAxiom». The source identifier is a scoped register reference and does not determine architectural type. Applicability, release scope and unresolved thresholds remain as stated above.


{\footnotesize \href{../../../../../../../continuum/motivation/principle/architectural-axiom/db8d320dc4c0454080f8c0f3c4e52801.md}{Canonical record: db8d320dc4c0454080f8c0f3c4e52801}; full qualifications and provenance.}\par
\section*{I.10.17 [ArchitecturalAxiom] Accepted Activity Durable Recording}
\phantomsection
\label{navigation:I.10.17}
\addcontentsline{toc}{section}{I.10.17 [ArchitecturalAxiom] Accepted Activity Durable Recording}

\label{element:6ee491b8c08d41868c3a7652fb794756}
\index{[ArchitecturalAxiom] Accepted Activity Durable Recording (I.10.17)}
{\small Principle / ArchitecturalAxiom · Solution Architect-reviewed Iteration 001; authorised ADD; resulting increment review pending · current}\par

Accepted processing activities shall be durably recorded before acknowledgement and retained until their associated processing obligations have been discharged.



Source entry: \textbf{NFR-067 — Accepted Activity Durable Recording}. Primary subject: \texttt{\seqsplit{G07.05}}.



The source classification is Principle «ArchitecturalAxiom». The source identifier is a scoped register reference and does not determine architectural type. Applicability, release scope and unresolved thresholds remain as stated above.


{\footnotesize \href{../../../../../../../continuum/motivation/principle/architectural-axiom/6ee491b8c08d41868c3a7652fb794756.md}{Canonical record: 6ee491b8c08d41868c3a7652fb794756}; full qualifications and provenance.}\par
\section*{I.10.18 [ArchitecturalAxiom] Native Operational Logging}
\phantomsection
\label{navigation:I.10.18}
\addcontentsline{toc}{section}{I.10.18 [ArchitecturalAxiom] Native Operational Logging}

\label{element:9d094b1b42644ab5aa2071a9460778e2}
\index{[ArchitecturalAxiom] Native Operational Logging (I.10.18)}
{\small Principle / ArchitecturalAxiom · Solution Architect-reviewed Iteration 001; authorised ADD; resulting increment review pending · current}\par

Provide a distinct, native capability for operational log capture and display, without dependency on external logging services.



Source entry: \textbf{NFR-075 — Native Operational Logging}. Primary subject: \texttt{\seqsplit{G09.05}}.



The source classification is Principle «ArchitecturalAxiom». The source identifier is a scoped register reference and does not determine architectural type. Applicability, release scope and unresolved thresholds remain as stated above.


{\footnotesize \href{../../../../../../../continuum/motivation/principle/architectural-axiom/9d094b1b42644ab5aa2071a9460778e2.md}{Canonical record: 9d094b1b42644ab5aa2071a9460778e2}; full qualifications and provenance.}\par
\section*{I.10.19 Responsive Interface Layouts}
\phantomsection
\label{navigation:I.10.19}
\addcontentsline{toc}{section}{I.10.19 Responsive Interface Layouts}

\label{element:2e276eda09a94670bee0619468947895}
\index{Responsive Interface Layouts (I.10.19)}
{\small Principle · Solution Architect-reviewed Iteration 001; authorised ADD; resulting increment review pending · current}\par

Provide responsive UI layouts adapting to supported device display characteristics and capabilities.



Source entry: \textbf{NFR-080 — Responsive Interface Layouts}. Primary subject: \texttt{\seqsplit{G12.01}}.



The source classification is Principle. The source identifier is a scoped register reference and does not determine architectural type. Applicability, release scope and unresolved thresholds remain as stated above.


{\footnotesize \href{../../../../../../../continuum/motivation/principle/2e276eda09a94670bee0619468947895.md}{Canonical record: 2e276eda09a94670bee0619468947895}; full qualifications and provenance.}\par
\section*{I.10.20 Contextual Navigation}
\phantomsection
\label{navigation:I.10.20}
\addcontentsline{toc}{section}{I.10.20 Contextual Navigation}

\label{element:39d698ff446d48ab9358844666fcafe5}
\index{Contextual Navigation (I.10.20)}
{\small Principle · Solution Architect-reviewed Iteration 001; authorised ADD; resulting increment review pending · current}\par

Core UIs should support contextual navigation where practical, preserving relevant operational context and authorisation.



Source entry: \textbf{NFR-082 — Contextual Navigation}. Primary subject: \texttt{\seqsplit{G12.04}}.



The source classification is Principle. The source identifier is a scoped register reference and does not determine architectural type. Applicability, release scope and unresolved thresholds remain as stated above.


{\footnotesize \href{../../../../../../../continuum/motivation/principle/39d698ff446d48ab9358844666fcafe5.md}{Canonical record: 39d698ff446d48ab9358844666fcafe5}; full qualifications and provenance.}\par
\section*{I.10.21 [ArchitecturalAxiom] Reference-Based Audit Records}
\phantomsection
\label{navigation:I.10.21}
\addcontentsline{toc}{section}{I.10.21 [ArchitecturalAxiom] Reference-Based Audit Records}

\label{element:9202f9b1adf84d4882a507a2d7280710}
\index{[ArchitecturalAxiom] Reference-Based Audit Records (I.10.21)}
{\small Principle / ArchitecturalAxiom · Solution Architect-reviewed Iteration 001; authorised ADD; resulting increment review pending · current}\par

Avoid duplicating persisted information in audit records where references suffice; exceptions when information itself causes/is the subject of the audit or no persisted instance exists (e.g. transformations).



Source entry: \textbf{NFR-085 — Reference-Based Audit Records}. Primary subject: \texttt{\seqsplit{G04.03}}.



The source classification is Principle «ArchitecturalAxiom». The source identifier is a scoped register reference and does not determine architectural type. Applicability, release scope and unresolved thresholds remain as stated above.


{\footnotesize \href{../../../../../../../continuum/motivation/principle/architectural-axiom/9202f9b1adf84d4882a507a2d7280710.md}{Canonical record: 9202f9b1adf84d4882a507a2d7280710}; full qualifications and provenance.}\par
\section*{I.10.22 [ArchitecturalAxiom] Independent Audit Service}
\phantomsection
\label{navigation:I.10.22}
\addcontentsline{toc}{section}{I.10.22 [ArchitecturalAxiom] Independent Audit Service}

\label{element:41fa01556ade4f22a35288b072113767}
\index{[ArchitecturalAxiom] Independent Audit Service (I.10.22)}
{\small Principle / ArchitecturalAxiom · Solution Architect-reviewed Iteration 001; authorised ADD; resulting increment review pending · current}\par

Audit Service must not depend on external event capture services for core audit generation, capture or management.



Source entry: \textbf{NFR-086 — Independent Audit Service}. Primary subject: \texttt{\seqsplit{G04.01}}.



The source classification is Principle «ArchitecturalAxiom». The source identifier is a scoped register reference and does not determine architectural type. Applicability, release scope and unresolved thresholds remain as stated above.


{\footnotesize \href{../../../../../../../continuum/motivation/principle/architectural-axiom/41fa01556ade4f22a35288b072113767.md}{Canonical record: 41fa01556ade4f22a35288b072113767}; full qualifications and provenance.}\par
\section*{I.10.23 [ArchitecturalAxiom] Large Information Object Storage}
\phantomsection
\label{navigation:I.10.23}
\addcontentsline{toc}{section}{I.10.23 [ArchitecturalAxiom] Large Information Object Storage}

\label{element:06ac0b10e0614c5cb33499c1e6c4edc7}
\index{[ArchitecturalAxiom] Large Information Object Storage (I.10.23)}
{\small Principle / ArchitecturalAxiom · Solution Architect-reviewed Iteration 001; authorised ADD; resulting increment review pending · current}\par

Information objects exceeding a defined size threshold shall be persisted in fit-for-purpose storage within the Harmonia ecosystem and exchanged by reference.



Source entry: \textbf{NFR-087 — Large Information Object Storage}. Primary subject: \texttt{\seqsplit{G02.04}}.



The source classification is Principle «ArchitecturalAxiom». The source identifier is a scoped register reference and does not determine architectural type. Applicability, release scope and unresolved thresholds remain as stated above.


{\small Recorded unresolved matters: largeObjectSizeThreshold; exact scope remains in the canonical record.}\par
{\footnotesize \href{../../../../../../../continuum/motivation/principle/architectural-axiom/06ac0b10e0614c5cb33499c1e6c4edc7.md}{Canonical record: 06ac0b10e0614c5cb33499c1e6c4edc7}; full qualifications and provenance.}\par
\section*{I.10.24 Robust MLLP Connection Management}
\phantomsection
\label{navigation:I.10.24}
\addcontentsline{toc}{section}{I.10.24 Robust MLLP Connection Management}

\label{element:e7b7d94a94194ef28a12051a550d3f3f}
\index{Robust MLLP Connection Management (I.10.24)}
{\small Principle · Solution Architect-reviewed Iteration 001; authorised ADD; resulting increment review pending · current}\par

MLLP connection management shall tolerate poorly behaving clients, including slow closure and repeated/multiple connection establishment, without compromising integrity.



Source entry: \textbf{NFR-090 — Robust MLLP Connection Management}. Primary subject: \texttt{\seqsplit{G05.01}}.



The source classification is Principle. The source identifier is a scoped register reference and does not determine architectural type. Applicability, release scope and unresolved thresholds remain as stated above.


{\footnotesize \href{../../../../../../../continuum/motivation/principle/e7b7d94a94194ef28a12051a550d3f3f.md}{Canonical record: e7b7d94a94194ef28a12051a550d3f3f}; full qualifications and provenance.}\par
\section*{I.10.25 Governed HL7 v2 Tolerance}
\phantomsection
\label{navigation:I.10.25}
\addcontentsline{toc}{section}{I.10.25 Governed HL7 v2 Tolerance}

\label{element:0968fc163ae54514bc1f02471fe6d866}
\index{Governed HL7 v2 Tolerance (I.10.25)}
{\small Principle · Solution Architect-reviewed Iteration 001; authorised ADD; resulting increment review pending · current}\par

Assess and report inbound HL7 v2.x structural conformance; support governed, configurable sending-system-specific tolerance of nonconformance when reliable interpretation is possible. Tolerated messages remain nonconformant.



Source entry: \textbf{NFR-091 — Governed HL7 v2 Tolerance}. Primary subject: \texttt{\seqsplit{G05.03}}.



The source classification is Principle. The source identifier is a scoped register reference and does not determine architectural type. Applicability, release scope and unresolved thresholds remain as stated above.


{\footnotesize \href{../../../../../../../continuum/motivation/principle/0968fc163ae54514bc1f02471fe6d866.md}{Canonical record: 0968fc163ae54514bc1f02471fe6d866}; full qualifications and provenance.}\par
\section*{I.10.26 [ArchitecturalAxiom] Operational State Reporting}
\phantomsection
\label{navigation:I.10.26}
\addcontentsline{toc}{section}{I.10.26 [ArchitecturalAxiom] Operational State Reporting}

\label{element:d7e3f0388d0043259585d1e4c5c948d6}
\index{[ArchitecturalAxiom] Operational State Reporting (I.10.26)}
{\small Principle / ArchitecturalAxiom · Solution Architect-reviewed Iteration 001; authorised ADD; resulting increment review pending · current}\par

Track/report subsystem, workflow, queue, persistence-service and audit-service state. Logging and OA\&M monitoring services require only basic health and availability monitoring.



Source entry: \textbf{NFR-093 — Operational State Reporting}. Primary subject: \texttt{\seqsplit{G09.01}}.



The source classification is Principle «ArchitecturalAxiom». The source identifier is a scoped register reference and does not determine architectural type. Applicability, release scope and unresolved thresholds remain as stated above.


{\footnotesize \href{../../../../../../../continuum/motivation/principle/architectural-axiom/d7e3f0388d0043259585d1e4c5c948d6.md}{Canonical record: d7e3f0388d0043259585d1e4c5c948d6}; full qualifications and provenance.}\par
\section*{I.10.27 Timezone Correctness}
\phantomsection
\label{navigation:I.10.27}
\addcontentsline{toc}{section}{I.10.27 Timezone Correctness}

\label{element:3da1fba95a00438db09c342dfba8b693}
\index{Timezone Correctness (I.10.27)}
{\small Principle · Solution Architect-reviewed Iteration 001; authorised ADD; resulting increment review pending · current}\par

Support concurrent timezones; preserve required timezone detail for local temporal values, distinguish absolute instants and handle daylight saving correctly.



Source entry: \textbf{NFR-095 — Timezone Correctness}. Primary subject: \texttt{\seqsplit{G12.05}}.



The source classification is Principle. The source identifier is a scoped register reference and does not determine architectural type. Applicability, release scope and unresolved thresholds remain as stated above.


{\footnotesize \href{../../../../../../../continuum/motivation/principle/3da1fba95a00438db09c342dfba8b693.md}{Canonical record: 3da1fba95a00438db09c342dfba8b693}; full qualifications and provenance.}\par
\section*{I.10.28 Target User Expertise}
\phantomsection
\label{navigation:I.10.28}
\addcontentsline{toc}{section}{I.10.28 Target User Expertise}

\label{element:1c0b0e15434d47b1b08ad86e087a45d1}
\index{Target User Expertise (I.10.28)}
{\small Principle · Solution Architect-reviewed Iteration 001; authorised ADD; resulting increment review pending · current}\par

R1.x/R2.x UIs target technically proficient and/or clinically trained users; novice/nontechnical/nonclinical user support is not required.



Source entry: \textbf{NFR-096 — Target User Expertise}. Primary subject: \texttt{\seqsplit{G12.01}}.



The source classification is Principle. The source identifier is a scoped register reference and does not determine architectural type. Applicability, release scope and unresolved thresholds remain as stated above.


{\footnotesize \href{../../../../../../../continuum/motivation/principle/1c0b0e15434d47b1b08ad86e087a45d1.md}{Canonical record: 1c0b0e15434d47b1b08ad86e087a45d1}; full qualifications and provenance.}\par
\section*{I.10.29 [ArchitecturalAxiom] Coherent Workflow Scheduler}
\phantomsection
\label{navigation:I.10.29}
\addcontentsline{toc}{section}{I.10.29 [ArchitecturalAxiom] Coherent Workflow Scheduler}

\label{element:99adfc9b5af049748df91755d395c247}
\index{[ArchitecturalAxiom] Coherent Workflow Scheduler (I.10.29)}
{\small Principle / ArchitecturalAxiom · Solution Architect-reviewed Iteration 001; authorised ADD; resulting increment review pending · current}\par

Use one coherent logically centralised (not necessarily physically centralised) scheduler for workflow activity.



Source entry: \textbf{NFR-097 — Coherent Workflow Scheduler}. Primary subject: \texttt{\seqsplit{G06.02}}.



The source classification is Principle «ArchitecturalAxiom». The source identifier is a scoped register reference and does not determine architectural type. Applicability, release scope and unresolved thresholds remain as stated above.


{\footnotesize \href{../../../../../../../continuum/motivation/principle/architectural-axiom/99adfc9b5af049748df91755d395c247.md}{Canonical record: 99adfc9b5af049748df91755d395c247}; full qualifications and provenance.}\par
\section*{I.10.30 [ArchitecturalAxiom] Coherent Calendar Service}
\phantomsection
\label{navigation:I.10.30}
\addcontentsline{toc}{section}{I.10.30 [ArchitecturalAxiom] Coherent Calendar Service}

\label{element:a9b1d031fa844057aa030f09e4fa6b01}
\index{[ArchitecturalAxiom] Coherent Calendar Service (I.10.30)}
{\small Principle / ArchitecturalAxiom · Solution Architect-reviewed Iteration 001; authorised ADD; resulting increment review pending · current}\par

Where needed, use one coherent logically centralised calendar for appointments, calendar scheduling and alarms; calendar and workflow scheduler are almost entirely independent.



Source entry: \textbf{NFR-098 — Coherent Calendar Service}. Primary subject: \texttt{\seqsplit{G06.02}}.



The source classification is Principle «ArchitecturalAxiom». The source identifier is a scoped register reference and does not determine architectural type. Applicability, release scope and unresolved thresholds remain as stated above.


{\footnotesize \href{../../../../../../../continuum/motivation/principle/architectural-axiom/a9b1d031fa844057aa030f09e4fa6b01.md}{Canonical record: a9b1d031fa844057aa030f09e4fa6b01}; full qualifications and provenance.}\par
\section*{I.10.31 Structured Information Encapsulation}
\phantomsection
\label{navigation:I.10.31}
\addcontentsline{toc}{section}{I.10.31 Structured Information Encapsulation}

\label{element:6973441f2d62448188af9bfce5ad9d9f}
\index{Structured Information Encapsulation (I.10.31)}
{\small Principle · Solution Architect-reviewed Iteration 001; authorised ADD; resulting increment review pending · current}\par

Prefer XML for structured information encapsulation where practical; allow JSON for UI interactions and other formats required by interoperability/contracts.



Source entry: \textbf{NFR-099 — Structured Information Encapsulation}. Primary subject: \texttt{\seqsplit{G05.03}}.



The source classification is Principle. The source identifier is a scoped register reference and does not determine architectural type. Applicability, release scope and unresolved thresholds remain as stated above.


{\footnotesize \href{../../../../../../../continuum/motivation/principle/6973441f2d62448188af9bfce5ad9d9f.md}{Canonical record: 6973441f2d62448188af9bfce5ad9d9f}; full qualifications and provenance.}\par
\section*{I.10.32 [ArchitecturalAxiom] Unbounded Clinical Information Retention}
\phantomsection
\label{navigation:I.10.32}
\addcontentsline{toc}{section}{I.10.32 [ArchitecturalAxiom] Unbounded Clinical Information Retention}

\label{element:4c270732fecf4e79bb83c938861ce05b}
\index{[ArchitecturalAxiom] Unbounded Clinical Information Retention (I.10.32)}
{\small Principle / ArchitecturalAxiom · Solution Architect-reviewed Iteration 001; authorised ADD; resulting increment review pending · current}\par

Impose no inherent temporal ceiling on patient/practitioner information retention; supporting infrastructure capacity and applicable governance still apply.



Source entry: \textbf{NFR-103 — Unbounded Clinical Information Retention}. Primary subject: \texttt{\seqsplit{G02.03}}.



The source classification is Principle «ArchitecturalAxiom». The source identifier is a scoped register reference and does not determine architectural type. Applicability, release scope and unresolved thresholds remain as stated above.


{\footnotesize \href{../../../../../../../continuum/motivation/principle/architectural-axiom/4c270732fecf4e79bb83c938861ce05b.md}{Canonical record: 4c270732fecf4e79bb83c938861ce05b}; full qualifications and provenance.}\par
\section*{I.10.33 R1.x Internationalisation}
\phantomsection
\label{navigation:I.10.33}
\addcontentsline{toc}{section}{I.10.33 R1.x Internationalisation}

\label{element:920606fdc7c14a27b7e39a08105c77b6}
\index{R1.x Internationalisation (I.10.33)}
{\small Principle · Solution Architect-reviewed Iteration 001; authorised ADD; resulting increment review pending · current}\par

R1.x must be architecturally multilingual (i18n), but only Australian English (en-AU) resources/content need be populated and delivered.



Source entry: \textbf{NFR-104 — R1.x Internationalisation}. Primary subject: \texttt{\seqsplit{G12.05}}.



The source classification is Principle. The source identifier is a scoped register reference and does not determine architectural type. Applicability, release scope and unresolved thresholds remain as stated above.


{\footnotesize \href{../../../../../../../continuum/motivation/principle/920606fdc7c14a27b7e39a08105c77b6.md}{Canonical record: 920606fdc7c14a27b7e39a08105c77b6}; full qualifications and provenance.}\par
\section*{I.10.34 [ArchitecturalAxiom] Component Recovery Consistency}
\phantomsection
\label{navigation:I.10.34}
\addcontentsline{toc}{section}{I.10.34 [ArchitecturalAxiom] Component Recovery Consistency}

\label{element:545465b03f794a509c932bd3126fbe57}
\index{[ArchitecturalAxiom] Component Recovery Consistency (I.10.34)}
{\small Principle / ArchitecturalAxiom · Solution Architect-reviewed Iteration 001; authorised ADD; resulting increment review pending · current}\par

Individual Harmonia components shall, where practical, support recovery to a known, operationally consistent state following failure.



Source entry: \textbf{NFR-107 — Component Recovery Consistency}. Primary subject: \texttt{\seqsplit{G07.03}}.



The source classification is Principle «ArchitecturalAxiom». The source identifier is a scoped register reference and does not determine architectural type. Applicability, release scope and unresolved thresholds remain as stated above.


{\footnotesize \href{../../../../../../../continuum/motivation/principle/architectural-axiom/545465b03f794a509c932bd3126fbe57.md}{Canonical record: 545465b03f794a509c932bd3126fbe57}; full qualifications and provenance.}\par
\section*{I.10.35 [ArchitecturalAxiom] Planned Restart Preservation}
\phantomsection
\label{navigation:I.10.35}
\addcontentsline{toc}{section}{I.10.35 [ArchitecturalAxiom] Planned Restart Preservation}

\label{element:0ce24ee595a445e5b1a95e54aa9f1eac}
\index{[ArchitecturalAxiom] Planned Restart Preservation (I.10.35)}
{\small Principle / ArchitecturalAxiom · Solution Architect-reviewed Iteration 001; authorised ADD; resulting increment review pending · current}\par

Planned system restarts and deployments shall preserve accepted governed work and its recoverable processing state.



Source entry: \textbf{NFR-109 — Planned Restart Preservation}. Primary subject: \texttt{\seqsplit{G07.03}}.



The source classification is Principle «ArchitecturalAxiom». The source identifier is a scoped register reference and does not determine architectural type. Applicability, release scope and unresolved thresholds remain as stated above.


{\footnotesize \href{../../../../../../../continuum/motivation/principle/architectural-axiom/0ce24ee595a445e5b1a95e54aa9f1eac.md}{Canonical record: 0ce24ee595a445e5b1a95e54aa9f1eac}; full qualifications and provenance.}\par
\section*{I.10.36 Degraded Mode Continuity}
\phantomsection
\label{navigation:I.10.36}
\addcontentsline{toc}{section}{I.10.36 Degraded Mode Continuity}

\label{element:51191ca802f54efb8554d20423199cd0}
\index{Degraded Mode Continuity (I.10.36)}
{\small Principle · Solution Architect-reviewed Iteration 001; authorised ADD; resulting increment review pending · current}\par

Harmonia should, where practical, continue providing unaffected capabilities when non-essential components or services become unavailable.



Source entry: \textbf{NFR-110 — Degraded Mode Continuity}. Primary subject: \texttt{\seqsplit{G07.05}}.



The source classification is Principle. The source identifier is a scoped register reference and does not determine architectural type. Applicability, release scope and unresolved thresholds remain as stated above.


{\footnotesize \href{../../../../../../../continuum/motivation/principle/51191ca802f54efb8554d20423199cd0.md}{Canonical record: 51191ca802f54efb8554d20423199cd0}; full qualifications and provenance.}\par
\section*{I.10.37 Processing Resource Isolation}
\phantomsection
\label{navigation:I.10.37}
\addcontentsline{toc}{section}{I.10.37 Processing Resource Isolation}

\label{element:41ce95e00e064d809ad4bda818705510}
\index{Processing Resource Isolation (I.10.37)}
{\small Principle · Solution Architect-reviewed Iteration 001; authorised ADD; resulting increment review pending · current}\par

Harmonia should, where practical, prevent resource exhaustion or excessive resource consumption within one processing context from adversely affecting the operation of unrelated processing contexts.



Source entry: \textbf{NFR-112 — Processing Resource Isolation}. Primary subject: \texttt{\seqsplit{G08.04}}.



The source classification is Principle. The source identifier is a scoped register reference and does not determine architectural type. Applicability, release scope and unresolved thresholds remain as stated above.


{\footnotesize \href{../../../../../../../continuum/motivation/principle/41ce95e00e064d809ad4bda818705510.md}{Canonical record: 41ce95e00e064d809ad4bda818705510}; full qualifications and provenance.}\par
\section*{I.10.38 Cross-Component Resource Contention}
\phantomsection
\label{navigation:I.10.38}
\addcontentsline{toc}{section}{I.10.38 Cross-Component Resource Contention}

\label{element:5bd9cc19d41e47d49978afaeea8343ed}
\index{Cross-Component Resource Contention (I.10.38)}
{\small Principle · Solution Architect-reviewed Iteration 001; authorised ADD; resulting increment review pending · current}\par

Harmonia shall endeavour to minimise cross-component resource contention and its impact on workflow execution, recognising that distributed architectures and non-real-time operating environments limit the ability to reliably enforce workload prioritisation.



Source entry: \textbf{NFR-115 — Cross-Component Resource Contention}. Primary subject: \texttt{\seqsplit{G08.04}}.



The source classification is Principle. The source identifier is a scoped register reference and does not determine architectural type. Applicability, release scope and unresolved thresholds remain as stated above.


{\footnotesize \href{../../../../../../../continuum/motivation/principle/5bd9cc19d41e47d49978afaeea8343ed.md}{Canonical record: 5bd9cc19d41e47d49978afaeea8343ed}; full qualifications and provenance.}\par
\section*{I.10.39 [ArchitecturalAxiom] Information Archival and Retrieval}
\phantomsection
\label{navigation:I.10.39}
\addcontentsline{toc}{section}{I.10.39 [ArchitecturalAxiom] Information Archival and Retrieval}

\label{element:2429db4dbecc4a03a66f5189d76a6c38}
\index{[ArchitecturalAxiom] Information Archival and Retrieval (I.10.39)}
{\small Principle / ArchitecturalAxiom · Solution Architect-reviewed Iteration 001; authorised ADD; resulting increment review pending · current}\par

Harmonia shall support a governed framework for information archival and retrieval.



Source entry: \textbf{NFR-117 — Information Archival and Retrieval}. Primary subject: \texttt{\seqsplit{G02.03}}.



The source classification is Principle «ArchitecturalAxiom». The source identifier is a scoped register reference and does not determine architectural type. Applicability, release scope and unresolved thresholds remain as stated above.


{\footnotesize \href{../../../../../../../continuum/motivation/principle/architectural-axiom/2429db4dbecc4a03a66f5189d76a6c38.md}{Canonical record: 2429db4dbecc4a03a66f5189d76a6c38}; full qualifications and provenance.}\par
\section*{I.10.40 Configurable Retention and Disposal}
\phantomsection
\label{navigation:I.10.40}
\addcontentsline{toc}{section}{I.10.40 Configurable Retention and Disposal}

\label{element:b1a46fce85b4448fba55987af146d61e}
\index{Configurable Retention and Disposal (I.10.40)}
{\small Principle · Solution Architect-reviewed Iteration 001; authorised ADD; resulting increment review pending · current}\par

Harmonia should support configurable information retention and disposal policies, avoiding arbitrary application-imposed limitations where practical and within the capabilities of the underlying infrastructure and middleware.



Source entry: \textbf{NFR-118 — Configurable Retention and Disposal}. Primary subject: \texttt{\seqsplit{G02.03}}.



The source classification is Principle. The source identifier is a scoped register reference and does not determine architectural type. Applicability, release scope and unresolved thresholds remain as stated above.


{\footnotesize \href{../../../../../../../continuum/motivation/principle/b1a46fce85b4448fba55987af146d61e.md}{Canonical record: b1a46fce85b4448fba55987af146d61e}; full qualifications and provenance.}\par
\section*{I.10.41 [ArchitecturalAxiom] Cryptographic Key Rotation}
\phantomsection
\label{navigation:I.10.41}
\addcontentsline{toc}{section}{I.10.41 [ArchitecturalAxiom] Cryptographic Key Rotation}

\label{element:6df99d4979a942fe833e1c43317b385b}
\index{[ArchitecturalAxiom] Cryptographic Key Rotation (I.10.41)}
{\small Principle / ArchitecturalAxiom · Solution Architect-reviewed Iteration 001; authorised ADD; resulting increment review pending · current}\par

Harmonia's cryptographic services shall support the rotation and replacement of cryptographic keys and certificates to maintain the security, integrity and trustworthiness of cryptographically protected information and communications.



Source entry: \textbf{NFR-123 — Cryptographic Key Rotation}. Primary subject: \texttt{\seqsplit{G03.05}}.



The source classification is Principle «ArchitecturalAxiom». The source identifier is a scoped register reference and does not determine architectural type. Applicability, release scope and unresolved thresholds remain as stated above.


{\footnotesize \href{../../../../../../../continuum/motivation/principle/architectural-axiom/6df99d4979a942fe833e1c43317b385b.md}{Canonical record: 6df99d4979a942fe833e1c43317b385b}; full qualifications and provenance.}\par
\section*{I.10.42 [ArchitecturalAxiom] Credential and Grant Revocation}
\phantomsection
\label{navigation:I.10.42}
\addcontentsline{toc}{section}{I.10.42 [ArchitecturalAxiom] Credential and Grant Revocation}

\label{element:370d9a6e0e00402aba9720580a2e365e}
\index{[ArchitecturalAxiom] Credential and Grant Revocation (I.10.42)}
{\small Principle / ArchitecturalAxiom · Solution Architect-reviewed Iteration 001; authorised ADD; resulting increment review pending · current}\par

Harmonia shall support governed revocation of authentication credentials and authorisation grants, and operational isolation of components, modules and subsystems through centrally governed enforcement of their access to Harmonia's information services.



Source entry: \textbf{NFR-124 — Credential and Grant Revocation}. Primary subject: \texttt{\seqsplit{G03.03}}.



The source classification is Principle «ArchitecturalAxiom». The source identifier is a scoped register reference and does not determine architectural type. Applicability, release scope and unresolved thresholds remain as stated above.


{\footnotesize \href{../../../../../../../continuum/motivation/principle/architectural-axiom/370d9a6e0e00402aba9720580a2e365e.md}{Canonical record: 370d9a6e0e00402aba9720580a2e365e}; full qualifications and provenance.}\par
\section*{I.10.43 [ArchitecturalAxiom] Security Anomaly Reporting}
\phantomsection
\label{navigation:I.10.43}
\addcontentsline{toc}{section}{I.10.43 [ArchitecturalAxiom] Security Anomaly Reporting}

\label{element:859f884874094c96a6ace15e17131a4c}
\index{[ArchitecturalAxiom] Security Anomaly Reporting (I.10.43)}
{\small Principle / ArchitecturalAxiom · Solution Architect-reviewed Iteration 001; authorised ADD; resulting increment review pending · current}\par

Harmonia shall report detectable security-related failures, violations and anomalous conditions through its governed operational monitoring and reporting capabilities.



Source entry: \textbf{NFR-127 — Security Anomaly Reporting}. Primary subject: \texttt{\seqsplit{G03.06}}.



The source classification is Principle «ArchitecturalAxiom». The source identifier is a scoped register reference and does not determine architectural type. Applicability, release scope and unresolved thresholds remain as stated above.


{\footnotesize \href{../../../../../../../continuum/motivation/principle/architectural-axiom/859f884874094c96a6ace15e17131a4c.md}{Canonical record: 859f884874094c96a6ace15e17131a4c}; full qualifications and provenance.}\par
\section*{I.10.44 Environment-Specific Configuration}
\phantomsection
\label{navigation:I.10.44}
\addcontentsline{toc}{section}{I.10.44 Environment-Specific Configuration}

\label{element:48e5415eff454e3e9e987506999e62f2}
\index{Environment-Specific Configuration (I.10.44)}
{\small Principle · Solution Architect-reviewed Iteration 001; authorised ADD; resulting increment review pending · current}\par

Harmonia should, where practical, support environment-specific configuration without requiring modification or rebuilding of application artefacts.



Source entry: \textbf{NFR-130 — Environment-Specific Configuration}. Primary subject: \texttt{\seqsplit{G10.03}}.



The source classification is Principle. The source identifier is a scoped register reference and does not determine architectural type. Applicability, release scope and unresolved thresholds remain as stated above.


{\footnotesize \href{../../../../../../../continuum/motivation/principle/48e5415eff454e3e9e987506999e62f2.md}{Canonical record: 48e5415eff454e3e9e987506999e62f2}; full qualifications and provenance.}\par
\section*{I.10.45 [ArchitecturalAxiom] Adjacent Release Coexistence}
\phantomsection
\label{navigation:I.10.45}
\addcontentsline{toc}{section}{I.10.45 [ArchitecturalAxiom] Adjacent Release Coexistence}

\label{element:3ce5568f2ded498386fd6fcfb142f89a}
\index{[ArchitecturalAxiom] Adjacent Release Coexistence (I.10.45)}
{\small Principle / ArchitecturalAxiom · Solution Architect-reviewed Iteration 001; authorised ADD; resulting increment review pending · current}\par

Harmonia shall support the concurrent operation of two adjacent solution releases against the same governed information environment, enabling controlled transition between releases without requiring in-place upgrades. Red/Blue deployment: two adjacent release generations can run concurrently against the same governed information environment, with controlled activity transition; in-place upgrade not required.



Source entry: \textbf{NFR-131 — Adjacent Release Coexistence}. Primary subject: \texttt{\seqsplit{G10.05}}.



The source classification is Principle «ArchitecturalAxiom». The source identifier is a scoped register reference and does not determine architectural type. Applicability, release scope and unresolved thresholds remain as stated above.


{\footnotesize \href{../../../../../../../continuum/motivation/principle/architectural-axiom/3ce5568f2ded498386fd6fcfb142f89a.md}{Canonical record: 3ce5568f2ded498386fd6fcfb142f89a}; full qualifications and provenance.}\par
\section*{I.10.46 [ArchitecturalAxiom] Three-Generation Release Compatibility}
\phantomsection
\label{navigation:I.10.46}
\addcontentsline{toc}{section}{I.10.46 [ArchitecturalAxiom] Three-Generation Release Compatibility}

\label{element:33d348c2936b401aa567d8fa531147ea}
\index{[ArchitecturalAxiom] Three-Generation Release Compatibility (I.10.46)}
{\small Principle / ArchitecturalAxiom · Solution Architect-reviewed Iteration 001; authorised ADD; resulting increment review pending · current}\par

Harmonia release N shall support concurrent operation with either release N−1 or release N+1 against the same governed information environment.Three-generation compatibility window N−1/N/N+1; N runs concurrently with either N+1 or N−1, not necessarily all three. General-purpose rollback and reversible database migration not required.



Source entry: \textbf{NFR-132 — Three-Generation Release Compatibility}. Primary subject: \texttt{\seqsplit{G10.05}}.



The source classification is Principle «ArchitecturalAxiom». The source identifier is a scoped register reference and does not determine architectural type. Applicability, release scope and unresolved thresholds remain as stated above.


{\footnotesize \href{../../../../../../../continuum/motivation/principle/architectural-axiom/33d348c2936b401aa567d8fa531147ea.md}{Canonical record: 33d348c2936b401aa567d8fa531147ea}; full qualifications and provenance.}\par
\section*{I.10.47 [ArchitecturalAxiom] Governed External Interaction Retry}
\phantomsection
\label{navigation:I.10.47}
\addcontentsline{toc}{section}{I.10.47 [ArchitecturalAxiom] Governed External Interaction Retry}

\label{element:49c13eb367b44565a011a77b90108cb8}
\index{[ArchitecturalAxiom] Governed External Interaction Retry (I.10.47)}
{\small Principle / ArchitecturalAxiom · Solution Architect-reviewed Iteration 001; authorised ADD; resulting increment review pending · current}\par

Harmonia shall support governed retry of failed external interactions where recovery through repetition is appropriate, consistent with the interaction's defined failure-handling, retry and idempotency policies.



Source entry: \textbf{NFR-135 — Governed External Interaction Retry}. Primary subject: \texttt{\seqsplit{G05.05}}.



The source classification is Principle «ArchitecturalAxiom». The source identifier is a scoped register reference and does not determine architectural type. Applicability, release scope and unresolved thresholds remain as stated above.


{\footnotesize \href{../../../../../../../continuum/motivation/principle/architectural-axiom/49c13eb367b44565a011a77b90108cb8.md}{Canonical record: 49c13eb367b44565a011a77b90108cb8}; full qualifications and provenance.}\par
\section*{I.10.48 [ArchitecturalAxiom] External Connection Configuration}
\phantomsection
\label{navigation:I.10.48}
\addcontentsline{toc}{section}{I.10.48 [ArchitecturalAxiom] External Connection Configuration}

\label{element:f24e63c49135411bbfbc06f9385aeea5}
\index{[ArchitecturalAxiom] External Connection Configuration (I.10.48)}
{\small Principle / ArchitecturalAxiom · Solution Architect-reviewed Iteration 001; authorised ADD; resulting increment review pending · current}\par

Harmonia shall support the configuration of external-system connections independently of the implementation of the components responsible for those interactions.



Source entry: \textbf{NFR-139 — External Connection Configuration}. Primary subject: \texttt{\seqsplit{G05.01}}.



The source classification is Principle «ArchitecturalAxiom». The source identifier is a scoped register reference and does not determine architectural type. Applicability, release scope and unresolved thresholds remain as stated above.


{\footnotesize \href{../../../../../../../continuum/motivation/principle/architectural-axiom/f24e63c49135411bbfbc06f9385aeea5.md}{Canonical record: f24e63c49135411bbfbc06f9385aeea5}; full qualifications and provenance.}\par
\section*{I.10.49 [ArchitecturalAxiom] Governed OA\&M Information}
\phantomsection
\label{navigation:I.10.49}
\addcontentsline{toc}{section}{I.10.49 [ArchitecturalAxiom] Governed OA\&M Information}

\label{element:43593f7a9566410ebc1aa30dbe935848}
\index{[ArchitecturalAxiom] Governed OA\&M Information (I.10.49)}
{\small Principle / ArchitecturalAxiom · Solution Architect-reviewed Iteration 001; authorised ADD; resulting increment review pending · current}\par

Harmonia's operational administration, monitoring and management (OA\&M) information shall constitute a governed information set, subject to the applicable information governance principles, mandates and requirements, irrespective of the technologies or mechanisms used for its collection, persistence and management.



Source entry: \textbf{NFR-142 — Governed OA\&M Information}. Primary subject: \texttt{\seqsplit{G09.05}}.



The source classification is Principle «ArchitecturalAxiom». The source identifier is a scoped register reference and does not determine architectural type. Applicability, release scope and unresolved thresholds remain as stated above.


{\footnotesize \href{../../../../../../../continuum/motivation/principle/architectural-axiom/43593f7a9566410ebc1aa30dbe935848.md}{Canonical record: 43593f7a9566410ebc1aa30dbe935848}; full qualifications and provenance.}\par
\section*{I.10.50 Root Cause Investigation}
\phantomsection
\label{navigation:I.10.50}
\addcontentsline{toc}{section}{I.10.50 Root Cause Investigation}

\label{element:9cdbf642158c4c37a6f3f10a6dd9ebfe}
\index{Root Cause Investigation (I.10.50)}
{\small Principle · Solution Architect-reviewed Iteration 001; authorised ADD; resulting increment review pending · current}\par

Harmonia should support the investigation and analysis of operational failures, degraded conditions and abnormal behaviour to assist in identifying their underlying causes, potentially incorporating automated and AI-assisted analysis within governed Assurance workflows.



Source entry: \textbf{NFR-145 — Root Cause Investigation}. Primary subject: \texttt{\seqsplit{G09.02}}.



The source classification is Principle. The source identifier is a scoped register reference and does not determine architectural type. Applicability, release scope and unresolved thresholds remain as stated above.


{\footnotesize \href{../../../../../../../continuum/motivation/principle/9cdbf642158c4c37a6f3f10a6dd9ebfe.md}{Canonical record: 9cdbf642158c4c37a6f3f10a6dd9ebfe}; full qualifications and provenance.}\par
\section*{I.10.51 Repeatable Performance Testing}
\phantomsection
\label{navigation:I.10.51}
\addcontentsline{toc}{section}{I.10.51 Repeatable Performance Testing}

\label{element:1c2eb4d9aef4425288528907fed7ebc2}
\index{Repeatable Performance Testing (I.10.51)}
{\small Principle · Solution Architect-reviewed Iteration 001; authorised ADD; resulting increment review pending · current}\par

Harmonia should, where practical, support repeatable performance and capacity testing using representative workloads and operational conditions, enabling meaningful comparison of results across solution releases.



Source entry: \textbf{NFR-149 — Repeatable Performance Testing}. Primary subject: \texttt{\seqsplit{G11.04}}.



The source classification is Principle. The source identifier is a scoped register reference and does not determine architectural type. Applicability, release scope and unresolved thresholds remain as stated above.


{\footnotesize \href{../../../../../../../continuum/motivation/principle/1c2eb4d9aef4425288528907fed7ebc2.md}{Canonical record: 1c2eb4d9aef4425288528907fed7ebc2}; full qualifications and provenance.}\par
\section*{I.10.52 [ArchitecturalAxiom] Distributed Temporal Consistency}
\phantomsection
\label{navigation:I.10.52}
\addcontentsline{toc}{section}{I.10.52 [ArchitecturalAxiom] Distributed Temporal Consistency}

\label{element:ed3e8d7923d14f8cbc6f08d3673248f2}
\index{[ArchitecturalAxiom] Distributed Temporal Consistency (I.10.52)}
{\small Principle / ArchitecturalAxiom · Solution Architect-reviewed Iteration 001; authorised ADD; resulting increment review pending · current}\par

Harmonia shall support the consistent representation, interpretation and management of temporal information across distributed operational domains, including the identification and management of event horizons, temporal ordering and temporal uncertainty where relevant to governed processing activities.



Source entry: \textbf{NFR-153 — Distributed Temporal Consistency}. Primary subject: \texttt{\seqsplit{G13.02}}.



The source classification is Principle «ArchitecturalAxiom». The source identifier is a scoped register reference and does not determine architectural type. Applicability, release scope and unresolved thresholds remain as stated above.


{\footnotesize \href{../../../../../../../continuum/motivation/principle/architectural-axiom/ed3e8d7923d14f8cbc6f08d3673248f2.md}{Canonical record: ed3e8d7923d14f8cbc6f08d3673248f2}; full qualifications and provenance.}\par
\section*{I.10.53 [ArchitecturalAxiom] External Information Sovereignty}
\phantomsection
\label{navigation:I.10.53}
\addcontentsline{toc}{section}{I.10.53 [ArchitecturalAxiom] External Information Sovereignty}

\label{element:aa179ef5bdf44b4eab3f2c855a8798d2}
\index{[ArchitecturalAxiom] External Information Sovereignty (I.10.53)}
{\small Principle / ArchitecturalAxiom · Solution Architect-reviewed Iteration 001; authorised ADD; resulting increment review pending · current}\par

Harmonia shall preserve the authority of external information sources by supporting governed references to externally managed information and enabling its retrieval through Pyria when required, without requiring that information to be incorporated into Harmonia's persistent information holdings.



Source entry: \textbf{NFR-154 — External Information Sovereignty}. Primary subject: \texttt{\seqsplit{G02.05}}.



The source classification is Principle «ArchitecturalAxiom». The source identifier is a scoped register reference and does not determine architectural type. Applicability, release scope and unresolved thresholds remain as stated above.


{\footnotesize \href{../../../../../../../continuum/motivation/principle/architectural-axiom/aa179ef5bdf44b4eab3f2c855a8798d2.md}{Canonical record: aa179ef5bdf44b4eab3f2c855a8798d2}; full qualifications and provenance.}\par
\section*{I.10.54 [ArchitecturalAxiom] Information Residency Governance}
\phantomsection
\label{navigation:I.10.54}
\addcontentsline{toc}{section}{I.10.54 [ArchitecturalAxiom] Information Residency Governance}

\label{element:349f526a401e4e299b854cbfedad0fca}
\index{[ArchitecturalAxiom] Information Residency Governance (I.10.54)}
{\small Principle / ArchitecturalAxiom · Solution Architect-reviewed Iteration 001; authorised ADD; resulting increment review pending · current}\par

Harmonia shall support the governance of information residency through its persistence framework, enabling the physical storage location of governed information to be determined and controlled by the applicable Mnemosyne PersistenceService, consistent with relevant information governance and deployment requirements.



Source entry: \textbf{NFR-155 — Information Residency Governance}. Primary subject: \texttt{\seqsplit{G02.05}}.



The source classification is Principle «ArchitecturalAxiom». The source identifier is a scoped register reference and does not determine architectural type. Applicability, release scope and unresolved thresholds remain as stated above.


{\footnotesize \href{../../../../../../../continuum/motivation/principle/architectural-axiom/349f526a401e4e299b854cbfedad0fca.md}{Canonical record: 349f526a401e4e299b854cbfedad0fca}; full qualifications and provenance.}\par
\section*{I.10.55 [ArchitecturalAxiom] Distributed Authorisation Enforcement}
\phantomsection
\label{navigation:I.10.55}
\addcontentsline{toc}{section}{I.10.55 [ArchitecturalAxiom] Distributed Authorisation Enforcement}

\label{element:172720628fa540928e34b52562c7067f}
\index{[ArchitecturalAxiom] Distributed Authorisation Enforcement (I.10.55)}
{\small Principle / ArchitecturalAxiom · Solution Architect explicitly approved inclusion; incorporated-result/visual review pending · current}\par

\par\medskip\textbf{Approved definition}\par


Every Harmonia component and subsystem SHALL enforce applicable authorisation requirements through a Themis-governed Policy Enforcement Point (PEP) before permitting a governed operation.



The Solution Architect's complete Subject/Views Task explicitly approves this
native Principle / ArchitecturalAxiom for the governed Continuum.
The incoming candidate/pending-ADD caption records its earlier source standing;
it does not override this subsequent inclusion approval. Original source wording,
native type, specialisation and Exchange identity evidence remain recoverable in
the retained immutable 004 export. This incorporation supplies neither HCRL closure
nor actual Archi visual acceptance.



This element retains its own bounded meaning. Broader existing security,
authorisation, messaging or stakeholder concerns do not establish identity
equivalence. No unrelated incoming properties, relationships, classifications,
Capability allocations or OpenSpec candidates are incorporated.


{\footnotesize \href{../../../../../../../continuum/motivation/principle/architectural-axiom/172720628fa540928e34b52562c7067f.md}{Canonical record: 172720628fa540928e34b52562c7067f}; full qualifications and provenance.}\par
\section*{I.10.56 [ArchitecturalAxiom] Passport Propagation / Availability}
\phantomsection
\label{navigation:I.10.56}
\addcontentsline{toc}{section}{I.10.56 [ArchitecturalAxiom] Passport Propagation / Availability}

\label{element:284b53f4a3c54acdba9e707124e997b8}
\index{[ArchitecturalAxiom] Passport Propagation / Availability (I.10.56)}
{\small Principle / ArchitecturalAxiom · Solution Architect explicitly approved inclusion; incorporated-result/visual review pending · current}\par

\par\medskip\textbf{Approved definition}\par


Harmonia SHALL ensure that the applicable Themis Passport is available across internal component and subsystem interactions, preserving the authorisation capability and its governing context. The Passport is a preconstructed Themis snapshot of effective authorisation capability, not identity credentials; availability may be by governed reference.



The Solution Architect's complete Subject/Views Task explicitly approves this
native Principle / ArchitecturalAxiom for the governed Continuum.
The incoming candidate/pending-ADD caption records its earlier source standing;
it does not override this subsequent inclusion approval. Original source wording,
native type, specialisation and Exchange identity evidence remain recoverable in
the retained immutable 004 export. This incorporation supplies neither HCRL closure
nor actual Archi visual acceptance.



This element retains its own bounded meaning. Broader existing security,
authorisation, messaging or stakeholder concerns do not establish identity
equivalence. No unrelated incoming properties, relationships, classifications,
Capability allocations or OpenSpec candidates are incorporated.


{\footnotesize \href{../../../../../../../continuum/motivation/principle/architectural-axiom/284b53f4a3c54acdba9e707124e997b8.md}{Canonical record: 284b53f4a3c54acdba9e707124e997b8}; full qualifications and provenance.}\par
